% Auto-generated by paper/render_tactical.py
% 47 tactical cards grouped by parent methodology.

\subsection{Audit and Relax Implicit Assumption}
\textit{Parent methodology id: \texttt{assumption\_audit\_and\_relax}; 11 tactical clusters.}

\begin{tcolorbox}[colback=bgskill!50, colframe=cgreen!70, title={\textbf{C00}: Structural Constraint Exploitation for Identifiability \textcolor{gray}{\small (O19/H0/R10, conf: \textbf{high})}}, breakable, before skip=4pt, after skip=4pt]
\textbf{Tactical pattern.} The shared move is: take a structural object that prior identifiability theory either eliminated for tractability, treated as a complicating nuisance, or never connected to identification at all --- and re-derive identifiability proofs by exploiting that object as a *source* of constraints rather than a source of ambiguity. Concrete instantiations include modeling instantaneous dependencies as a DAG and using acyclicity to pin down the parameter (ICLR\_2025\_1525), using the geometric support changes induced by interventions on latent factors (ICML\_2023\_0476, NeurIPS\_2023\_0640), exploiting the covariance geometry of single-node interventions for linear-causal-from-nonlinear-mixing recovery (NeurIPS\_2023\_0675), lifting causal effect identification to exchangeable data via the de Finetti mixing measure (NeurIPS\_2024\_1202), substituting an architectural additivity constraint for distributional auxiliary supervision (NeurIPS\_2023\_0587), and reducing deep polynomial-network identifiability to a partially symmetric tensor decomposition (NeurIPS\_2025\_1887). The tactical signature is always: name the discarded/hidden structure, encode it formally, then prove that the formal encoding implies recovery up to a stated equivalence class.

\textbf{Differentiation.} \textit{Cluster 9 (Structural Inductive Bias Formalization) is the closest sibling --- both audit a "negligible" structural property and re-formalize it. The split is what the formalization buys: cluster 9 produces *explanatory/predictive* theorems (random features beat linear under spiked covariance, NTK explains task arithmetic, GNN extrapolation is governed by an architectural invariant), whereas this cluster produces *identifiability/uniqueness* theorems that recover hidden objects (latent causes, ICA sources, causal effects, network weights) up to a stated equivalence class. A second contrast is cluster 30 (Objective Auditing via Explicit Constraint Injection): cluster 30 takes the discovered structure and bakes it into a training loss, while this cluster takes it and bakes it into a proof --- papers here typically derive theorems first and only then propose an estimator. A third contrast is cluster 5 (Structural Dependency Formalization for Tractable Fairness), which formalizes hidden dependencies to produce tractable *fairness objectives*; this cluster's outputs are recovery guarantees, not optimization-friendly criteria.}

\textbf{When to pick.} Pick this tactic when the target quantity is a hidden object (latent causes, ICA sources, network weights, treatment effects) whose recovery is currently blocked by a non-identifiability result, AND when prior work has explicitly *thrown away* or *oversimplified* a piece of structure (acyclicity, covariance, support geometry, exchangeability, additivity, network topology) that you can name. The signal is the phrase "we assume X to avoid this complication" in a prior paper --- that X is your candidate constraint to put back.

\textbf{Tactical failure.} Reject papers in this cluster fail in two predictable ways. First, they wrap an ML estimator around an existing identifiability result without discovering any new structural constraint, so the contribution reads as engineering rather than identification: ICLR\_2025\_1481 learns low-dimensional representations of high-dim instruments but inherits standard IV exclusion as an assumption rather than relaxing it via newly-recognized structure; ICLR\_2024\_0952 imports the Riesz representer for stability but does not exploit a new feature of the survival functional; ICLR\_2023\_0291 extends linear ICA to multi-view but its identifiability conditions reduce to standard diversity counts. Second, they substitute one strong assumption with another that *appears* weaker but is structurally just as restrictive, without either being testable or providing a payoff: ICLR\_2025\_1385's rank-preservation across counterfactual worlds is a near-universal ordering claim disguised as an ordinal relaxation; NeurIPS\_2024\_1234 invokes "switchable mechanisms" for soft interventions without nailing down the conditions under which the switches are themselves identifiable; ICLR\_2023\_0267 sets up a chicken-and-egg variational scheme where the joint identifiability of variables and graph is asserted rather than proven. The third (rarer) variant is ambitious joint-identification claims (NeurIPS\_2025\_1901 for hidden-driver + observed causal structure, ICLR\_2023\_0170 for unobserved confounders via mismatched pairs) where the structural constraint invoked is too vague to support a clean theorem. The discriminator is whether the paper produces a new identifiability theorem with explicit, checkable conditions, or merely an estimator that operates inside the old conditions.

\textbf{Examples.}
\begin{itemize}[leftmargin=12pt,topsep=2pt,itemsep=1pt]
  \item \textbf{[Oral]} \texttt{ICLR\_2025\_1525} \textemdash\ Canonical instance of putting a previously-discarded structural object back in: instantaneous dependencies were typically removed for tractability, and the paper proves they instead supply additional DAG-acyclicity constraints that drive identification.
  \\ \textcolor{gray}{\footnotesize \textit{Lesson:} Putting back a structural object the field eliminated for tractability is the highest-yield move --- instantaneous dependencies were a 'nuisance' until they became the identifiability resource.}
  \item \textbf{[Oral]} \texttt{NeurIPS\_2023\_0675} \textemdash\ Exemplifies the geometric-signature subvariant --- quadratic forms of the observable covariance under single-node interventions are shown to uniquely pin down the linear causal components even through general nonlinear mixing.
  \\ \textcolor{gray}{\footnotesize \textit{Lesson:} Geometric signatures (here: quadratic forms of intervened covariance) can pin down latent structure where pure conditional-independence machinery cannot --- pick the right algebraic object for the identification.}
  \item \textbf{[Oral]} \texttt{NeurIPS\_2024\_1202} \textemdash\ Lifts the causal-effect framework from i.i.d. to exchangeable data by treating de Finetti's mixing measure as the structural-equation-system surrogate, the cleanest 'import a representation theorem to enlarge the identifiable regime' move in the cluster.
  \\ \textcolor{gray}{\footnotesize \textit{Lesson:} Importing a representation theorem (here: de Finetti) at the right level of abstraction enlarges the identifiable regime --- exchangeability becomes a structural equation surrogate.}
  \item \textbf{[Oral]} \texttt{NeurIPS\_2023\_0587} \textemdash\ Demonstrates substituting a single architectural constraint (additive decoders) for distributional/temporal auxiliary signals and proving block-wise identifiability from rigidity alone --- exactly the cluster's 'structure as identification source' pattern.
  \\ \textcolor{gray}{\footnotesize \textit{Lesson:} A single architectural constraint (additive decoders) can replace heavier distributional / temporal supervision and still prove block-wise identifiability --- structure as identification source.}
  \item \textbf{[Reject]} \texttt{ICLR\_2025\_1481} \textemdash\ Wraps representation learning around standard IV partial-identification without discovering a new structural constraint that the high-dim setting unlocks --- the exclusion restriction is assumed rather than re-derived.
  \\ \textcolor{gray}{\footnotesize \textit{Lesson:} Wrapping ML estimators around an existing identifiability result without discovering a new structural constraint is incremental --- the high-dim setting must unlock something the low-dim version doesn't.}
  \item \textbf{[Reject]} \texttt{ICLR\_2025\_1385} \textemdash\ Substitutes one strong assumption (cardinal counterfactual model) with another (rank invariance across all worlds) that is global and untestable, illustrating the failure mode of pseudo-weak structural claims.
  \\ \textcolor{gray}{\footnotesize \textit{Lesson:} Substituting one strong assumption with another global untestable assumption is pseudo-relaxation --- the new condition must be locally testable or weaker, not just relabeled.}
  \item \textbf{[Reject]} \texttt{ICLR\_2024\_0952} \textemdash\ Imports the Riesz representer to stabilize survival-effect estimation, a methodological improvement that does not reveal or exploit a new identifying structure in the survival functional itself.
  \\ \textcolor{gray}{\footnotesize \textit{Lesson:} Importing a stabilization tool (Riesz representer) is methodological improvement, not structural exploitation --- the cluster requires revealing or exploiting a NEW identifying structure.}
\end{itemize}
\end{tcolorbox}

\begin{tcolorbox}[colback=bgskill!50, colframe=cgreen!70, title={\textbf{C03}: Selective Retention Under Non-Stationarity \textcolor{gray}{\small (O15/H1/R16, conf: \textbf{high})}}, breakable, before skip=4pt, after skip=4pt]
\textbf{Tactical pattern.} The shared move is to audit the implicit assumption that adaptation should overwrite prior state --- that "more updates strictly improve the model" --- and replace it with an explicit *retention partition* over parameters, components, exemplars, or update opportunities. Concretely, papers (i) name a concrete unit of accumulated knowledge (LoRA adapters, policy networks, BatchNorm affine params, replay exemplars, neuron activations, dual-GP inducing points, function-vector subspaces), (ii) introduce a gate that decides per-unit whether to freeze, blend (EMA / weighted composition / multi-step TD prior), reset, or refresh that unit when the environment shifts, and (iii) trigger the gate from a non-stationarity signal (entropy of outputs, SNR of recent experience, dormancy statistics, function-vector cosine, mean/covariance drift). The intervention is at the *parameter/component dynamics* level --- not at the loss function or the theoretical guarantee.

\textbf{Differentiation.} \textit{Versus sibling cluster 30 (Objective Auditing via Explicit Constraint Injection), which reformulates the *training loss* by adding a constraint that encodes a suppressed property (variance/causal invariance), this cluster leaves the per-step objective largely intact and intervenes on *which parameters or examples participate in the update*. Versus sibling cluster 39 (Objective Realignment for Robustness via Geometric Targeting), which retargets the loss surface toward worst-case geometry, this cluster's audit lever is temporal --- it gates updates over a sequence rather than reshaping a single optimization. Versus sibling cluster 9 (Structural Inductive Bias Formalization), which uses a hidden structural property to *explain* a phenomenon, this cluster operationalizes its audit into a runtime selection mechanism (freeze/reset/blend) rather than a derivation. The signature verb here is "retain selectively," not "constrain" or "prove."}

\textbf{When to pick.} Pick this instantiation when the problem has a temporal axis with non-stationary distribution shift (sequential tasks, test-time shift, long-horizon RL) AND the friction is that the same gradient that helps now hurts later --- i.e., the bottleneck is a retention/update tradeoff over identifiable model components. If the audit instead points at a missing term in the loss or a missing property of the data-generating process, route to clusters 30 or 34.

\textbf{Tactical failure.} Reject papers in this cluster fail when the retention mechanism is a heuristic blend of validated parts rather than a reframing keyed to a newly-named structural lever. ICLR\_2025\_1496 (diversity-aware replay) and ICLR\_2023\_0308 (PCA exemplar selection) propose buffer-selection rules that read as engineering refinements over herding/random --- no new audit, just a different scoring function. ICLR\_2025\_1491 (EMA with adaptive beta) and ICLR\_2024\_0955 (structural-knowledge regularizer) bolt a smoothing/regularization term onto an existing pipeline without showing the term resolves a previously invisible failure. ICML\_2025\_1805 (TD($\lambda$) continual learning) and NeurIPS\_2024\_1327 (variational continual TTA) import a sophisticated formalism (multi-step returns, Bayesian posterior) but the empirical lift over single-step VCL or entropy-min TTA is incremental and the audit ("error accumulates from single-step updates") is generic. NeurIPS\_2023\_0597 and ICLR\_2025\_1343 stack independently-validated components and discover interference too late. The common pattern: the retention partition is *unmotivated by a measured non-stationarity property* of the system (unlike dormancy fraction in ICML\_2023\_0549 or function-vector cosine in ICLR\_2025\_1642), so the mechanism reads as one more anti-forgetting trick.

\textbf{Examples.}
\begin{itemize}[leftmargin=12pt,topsep=2pt,itemsep=1pt]
  \item \textbf{[Oral]} \texttt{ICML\_2024\_1081} \textemdash\ Inverts the assumption that policy updates always help by formally tying update suppression to a measurable SNR drop in non-stationary RL --- a paradigmatic 'when-to-not-update' retention gate.
  \\ \textcolor{gray}{\footnotesize \textit{Lesson:} Tie update suppression to a measurable structural lever (SNR drop) --- when-to-not-update becomes a derivable gate, not a heuristic.}
  \item \textbf{[Oral]} \texttt{ICML\_2024\_1123} \textemdash\ Achieves no-forgetting by architectural composition over frozen prior policies, making selective retention a structural guarantee rather than a regularization penalty.
  \\ \textcolor{gray}{\footnotesize \textit{Lesson:} Architectural composition over frozen prior policies makes selective retention a structural guarantee --- beats regularization penalties hands down.}
  \item \textbf{[Oral]} \texttt{ICML\_2023\_0549} \textemdash\ Names a measurable retention failure (dormant neurons) and introduces a per-unit reset rule keyed to that statistic --- the audit-then-recycle pattern in its purest form.
  \\ \textcolor{gray}{\footnotesize \textit{Lesson:} Naming a measurable retention failure (dormant neurons) and adding a per-unit reset rule keyed to it is the audit-then-recycle pattern in pure form.}
  \item \textbf{[Oral]} \texttt{ICLR\_2025\_1577} \textemdash\ Replaces parameter-growing incremental learning with a fixed adapter pool plus class-specific importance weights, reframing retention as basis-reuse rather than accumulation.
  \\ \textcolor{gray}{\footnotesize \textit{Lesson:} Reframe retention as basis-reuse (fixed adapter pool + class-specific weights) instead of accumulation --- the right reframe replaces the parameter-growing assumption.}
  \item \textbf{[Reject]} \texttt{ICLR\_2025\_1496} \textemdash\ The retention mechanism (diversity-aware replay subset selection) is a buffer-scoring heuristic with no audit of a structural non-stationarity property the existing buffer rules miss.
  \\ \textcolor{gray}{\footnotesize \textit{Lesson:} A buffer-scoring heuristic without an audit of a specific structural non-stationarity property is a knob, not a retention reframe.}
  \item \textbf{[Reject]} \texttt{ICLR\_2025\_1491} \textemdash\ Uses EMA with adaptive $\beta$ as the sole anti-forgetting lever --- a smoothing knob added to instruction tuning rather than a reframing of which components should be retained and why.
  \\ \textcolor{gray}{\footnotesize \textit{Lesson:} An EMA $\beta$ knob is smoothing, not retention reframing --- adaptive smoothing alone does not identify which components should be retained and why.}
  \item \textbf{[Reject]} \texttt{ICML\_2025\_1805} \textemdash\ Imports TD($\lambda$)-style multi-step regularization into VCL but the audit ('single-step error accumulates') yields a formal extension with incremental gains rather than a newly-identified retention lever.
  \\ \textcolor{gray}{\footnotesize \textit{Lesson:} Importing a multi-step regularizer for incremental gains is a formal extension, not a newly-identified retention lever --- the cluster wants a NEW lever, not a refinement.}
\end{itemize}
\end{tcolorbox}

\begin{tcolorbox}[colback=bgskill!50, colframe=cgreen!70, title={\textbf{C05}: Structural Dependency Formalization for Tractable Fairness \textcolor{gray}{\small (O9/H0/R11, conf: \textbf{high})}}, breakable, before skip=4pt, after skip=4pt]
\textbf{Tactical pattern.} Locate a fairness setting whose existing framework presupposes that some structural element is fixed or absent --- static individual features, exchangeable train/test data, a single binary protected attribute, no feedback from decisions to future inputs, no latent confounder, no causal mediator --- then write down an explicit structural model of that element (a structural causal model with latents, a dynamical system over data distributions, a density-ratio-reweighted population, a Wasserstein/RKHS lift over the joint attribute space, a response function mapping predictions back to features) and *derive* a fairness objective or guarantee from the extended model that could not be expressed in the original. The signature verbs are "formalize the dependency, then re-derive" --- counterfactual fairness on a graph SCM with hidden confounders (ICLR\_2024\_0799), bias amplification as a stability condition of a feedback dynamical system (ICML\_2023\_0434), individual fairness as Wasserstein invariance over transport balls (ICLR\_2021\_0065), reciprocity as bilateral marginal contribution-vs-benefit (ICML\_2023\_0473), conformal coverage under feedback covariate shift via jackknife+ (ICML\_2023\_0477).

\textbf{Differentiation.} \textit{The closest sibling is cluster 30 (Objective Auditing via Explicit Constraint Injection), which also reformulates objectives --- but cluster 30 acts *on the loss* by adding an auxiliary term, regularizer, or adversarial penalty that directly enforces a missing statistical property (variance, invariance, independence). This cluster acts *upstream of the loss*: it builds a richer model of the data-generating or decision-making process (an SCM with mediators, a dynamical recursion, a covariate-shifted population, a bilateral exchange) and the fairness criterion falls out of that extended model --- e.g., ICML\_2023\_0570 derives a harm-detection test by *comparing two model populations* (group-aware vs group-blind), not by adding a constraint to one. It also differs from sibling 0 (Structural Constraint Exploitation for Identifiability), which formalizes hidden structure to *recover latent variables* in causal/representation learning; here the formalization serves a *normative criterion* (who is harmed, what counts as fair under feedback). And unlike sibling 39 (Objective Realignment for Robustness via Geometric Targeting), the target is distributional/social fairness, not worst-case loss geometry.}

\textbf{When to pick.} Pick this tactic when the fairness gap you want to address is invisible inside the existing framework --- i.e., the framework treats as fixed or absent the very thing causing the disparity (features that respond to predictions, a covariate shift between deploy and train, a confounder upstream of both attribute and outcome, a feedback loop from outputs to future inputs). The signal is that you need to write down a generative or dynamical model *before* you can even state the fairness target, not just add a term to an existing one.

\textbf{Tactical failure.} The dominant failure is stopping at the structural reformulation without earning a result that only the reformulation can deliver: a definitional move dressed as a framework. ICLR\_2024\_0800 extends the SCM with a "response function" mapping predictions back to features but reviewers read the counterfactual-over-trajectory criterion as a relabeling rather than a tractable estimator with new identifiability. ICLR\_2023\_0388 ports density-ratio reweighting from unsupervised domain adaptation to fairness almost verbatim, so the contribution reads as recombination rather than a fairness-specific advance. ICLR\_2025\_1542 names "social determinants of opportunity" as latent causal mediators, but the operational procedure reduces to standard mediation adjustment with sociological framing layered on top. ICLR\_2024\_0799 stacks hidden confounders + identifiable VAE + Gaussian-mixture priors on graphs, but the identifiability assumptions are strong, unverified, and the fairness gain over confounder-blind baselines is not isolated. ICLR\_2023\_0313 lifts cross-covariance to an RKHS for multi-type attributes but never ties the new measure to a fairness guarantee that binary-attribute methods provably miss. ICLR\_2023\_0381 combines unbiasedness with optimistic exploration in the bank-loan loop, but the regret bound and the fairness improvement are not shown to be jointly necessary consequences of the formalization. The shared pattern: the structural model is written down, the fairness object is renamed in its language, but no estimator, identifiability theorem, or empirical lift exclusive to the extended model is delivered.

\textbf{Examples.}
\begin{itemize}[leftmargin=12pt,topsep=2pt,itemsep=1pt]
  \item \textbf{[Oral]} \texttt{ICML\_2023\_0434} \textemdash\ Canonical instance: models the train$\to$deploy$\to$retrain pipeline as a dynamical system over data distributions and derives bias amplification as a stability condition (uniform faithfulness), so the fairness object only exists once the feedback dependency is explicit.
  \\ \textcolor{gray}{\footnotesize \textit{Lesson:} Modeling the train$\to$deploy$\to$retrain pipeline as a dynamical system makes bias amplification a stability condition (uniform faithfulness) --- the fairness object only exists once feedback is explicit.}
  \item \textbf{[Oral]} \texttt{ICLR\_2021\_0065} \textemdash\ Formalizes individual fairness as Wasserstein invariance over transport balls around each point --- the tractable regularizer is a derivative of the new structural object, not a constraint bolted onto an existing loss.
  \\ \textcolor{gray}{\footnotesize \textit{Lesson:} Wasserstein invariance over transport balls makes the regularizer a derivative of the new structural object, not a constraint bolted onto an existing loss.}
  \item \textbf{[Oral]} \texttt{ICML\_2023\_0570} \textemdash\ Builds the harm criterion by structurally comparing a group-aware predictor to a group-blind counterfactual model, exposing a fairness failure mode that aggregate accuracy and constraint-injection methods both hide.
  \\ \textcolor{gray}{\footnotesize \textit{Lesson:} Comparing group-aware vs group-blind counterfactual models surfaces a fairness failure mode that aggregate-accuracy and constraint-injection methods both hide.}
  \item \textbf{[Oral]} \texttt{ICML\_2023\_0473} \textemdash\ Replaces the unilateral data-valuation framing with a bilateral contribution$\leftrightarrow$benefit structural model, then proves reciprocity guarantees that no single-direction objective could express.
  \\ \textcolor{gray}{\footnotesize \textit{Lesson:} Bilateral contribution$\leftrightarrow$benefit structural model proves reciprocity guarantees no single-direction objective could express --- pick the right structural relation, not just a richer SCM.}
  \item \textbf{[Reject]} \texttt{ICLR\_2024\_0800} \textemdash\ Augments the SCM with a feature-response function for strategic agents but the resulting dynamic counterfactual-fairness criterion reads as definitional --- no new identification result or estimator is produced from the extended structure.
  \\ \textcolor{gray}{\footnotesize \textit{Lesson:} Augmenting an SCM with a strategic-agent feature-response function is definitional unless it produces a new identification result or estimator the original SCM cannot.}
  \item \textbf{[Reject]} \texttt{ICLR\_2023\_0388} \textemdash\ Re-estimates the fairness tradeoff under covariate shift by importing density-ratio reweighting from unsupervised domain adaptation almost unchanged, so the structural reformulation amounts to recombination rather than a fairness-specific derivation.
  \\ \textcolor{gray}{\footnotesize \textit{Lesson:} Importing density-ratio reweighting from UDA almost unchanged is recombination, not fairness-specific structural reformulation.}
  \item \textbf{[Reject]} \texttt{ICLR\_2025\_1542} \textemdash\ Posits community-level latent mediators ('social determinants of opportunity') as a causal extension, but the operational procedure collapses to standard mediation analysis with a sociological label, not a guarantee unique to the new structure.
  \\ \textcolor{gray}{\footnotesize \textit{Lesson:} A latent-mediator extension that reduces operationally to standard mediation analysis with a sociological label is a relabel, not a structural guarantee.}
\end{itemize}
\end{tcolorbox}

\begin{tcolorbox}[colback=bgskill!50, colframe=cgreen!70, title={\textbf{C09}: Structural Inductive Bias Formalization \textcolor{gray}{\small (O16/H1/R2, conf: \textbf{low})}}, breakable, before skip=4pt, after skip=4pt]
\textbf{Tactical pattern.} Hold the architecture/algorithm fixed and derive an EXACT closed-form mathematical object --- kernel spectrum, analytic score function, partition-function equivalence, frequency-domain interpolant, harmonic basis, or minimal sufficient geometric assumption --- that pins down the latent structural bias making the system work. Concrete verbs/objects recur across the cluster: "compute NTK eigenvalue decay" (ICLR\_2023\_0157), "derive the analytic score implied by locality+equivariance" (ICML\_2025\_1670), "prove convergence to minimum-degree Fourier interpolant" (ICML\_2023\_0460), "construct explicit RBM weight matrix whose partition function equals the HN's" (ICLR\_2021\_0047), "formalize the expansion assumption and prove sufficiency" (ICLR\_2021\_0070), "apply the Kolmogorov-Arnold theorem to determine where learnable nonlinearities go" (ICLR\_2025\_1464). The deliverable is a quantitative prediction or exact equivalence (memorization threshold, eigenvalue rate, phase boundary, equivalent formalism) --- not a new training objective, not a new architecture, not an identifiability theorem about latent variables.

\textbf{Differentiation.} \textit{The closest sibling is Cluster 30 (Objective Auditing via Explicit Constraint Injection), which also locates an implicit structural property --- but cluster 30's deliverable is a reformulated training objective or auxiliary loss that ENFORCES the property (e.g., VICReg variance/covariance regularizers, Bayesian causal-invariance prior). Cluster 9 instead leaves the training pipeline alone and DERIVES a closed-form characterization that EXPLAINS why empirically-successful systems already exhibit the property (e.g., ICLR\_2024\_0840 derives the harmonic-basis memorization threshold of an unmodified diffusion denoiser; ICML\_2025\_1670 derives the analytic score function implied by existing CNN inductive biases). Cluster 9 also differs from Cluster 0 (Structural Constraint Exploitation for Identifiability), whose payoff is recovering latent causes/parameters --- cluster 9's payoff is a mechanistic predictive theory of generalization, extrapolation, or creativity in the observable. If a paper rewrites the loss it is cluster 30; if it claims latent-variable identifiability it is cluster 0; if it derives an exact analytic prediction matching trained-model behavior, it is cluster 9.}

\textbf{When to pick.} Pick this cluster when the empirical phenomenon is already stable and the open question is "what closed-form mathematical object does the unmodified system actually compute," AND a well-developed external formalism (NTK theory, harmonic analysis, circuit theory, Boolean Fourier analysis, decomposition theorems) can be ported in to give an exact equivalence or quantitative rate. Avoid this lens if you intend to ship a new objective or architecture --- those belong to cluster 30 or to objective-realignment clusters.

\textbf{Tactical failure.} With only 2 reject papers in the cluster, the failure pattern is necessarily speculative --- confidence is low. The two visible failures both involve formalizations that fall short of an EXACT, generalizable closed-form characterization. ICLR\_2025\_1548 derives a phase transition for random features vs. linear models, but only inside a heavily idealized spiked-covariance + proportional-high-dimensional regime; the "exact" characterization is so regime-bound that reviewers read it as a narrow technical extension of an existing equivalence rather than a new structural insight, and the qualitative claim (alignment helps expressive models) is not surprising once the regime is fixed. ICLR\_2024\_0923 reformulates post-nonlinear causal fitting as a maximal-correlation constraint --- but the deliverable is an optimization reformulation, not a closed-form prediction or exact equivalence, so it lands in an awkward middle ground between cluster 9 (no analytic object derived) and cluster 30 (no replacement objective trained at scale). The pattern: papers that promise formalization but deliver either (a) an idealized-regime equivalence with no genuinely new structural object, or (b) a constraint reformulation lacking an exact closed-form prediction, get rejected for shipping the framing without the mathematical payoff.

\textbf{Examples.}
\begin{itemize}[leftmargin=12pt,topsep=2pt,itemsep=1pt]
  \item \textbf{[Oral]} \texttt{ICLR\_2024\_0840} \textemdash\ Identifies geometry-adaptive harmonic representations as the specific inductive bias of diffusion denoisers and derives an exact memorization-vs-generalization threshold from manifold dimensionality --- a quantitative prediction, not a new loss.
  \\ \textcolor{gray}{\footnotesize \textit{Lesson:} Identify geometry-adaptive harmonic representations as the inductive bias and derive an exact memorization-vs-generalization threshold from manifold dimensionality --- a quantitative prediction, not a new loss.}
  \item \textbf{[Oral]} \texttt{ICML\_2025\_1670} \textemdash\ Derives the EXACT analytic score function implied by convolutional locality and equivariance and shows it predicts the actual outputs of trained diffusion models --- an unusually literal closed-form characterization of architectural bias.
  \\ \textcolor{gray}{\footnotesize \textit{Lesson:} Derive the EXACT analytic score function implied by convolutional locality + equivariance and show it predicts trained model outputs --- unusually literal closed-form characterization.}
  \item \textbf{[Oral]} \texttt{ICLR\_2023\_0157} \textemdash\ Computes explicit GP/NTK kernel spectra for ResNets, separating conditioning from spectral bias via closed-form eigenvalue decay rates rather than reformulating training.
  \\ \textcolor{gray}{\footnotesize \textit{Lesson:} Compute explicit GP/NTK kernel spectra for ResNets --- separating conditioning from spectral bias via closed-form eigenvalue decay rates beats reformulating training.}
  \item \textbf{[Oral]} \texttt{ICLR\_2021\_0070} \textemdash\ Distills three empirically-successful methods into a single minimal expansion assumption and proves it is exactly sufficient for correctness --- the canonical 'identify minimal structural property, derive exact guarantee' move.
  \\ \textcolor{gray}{\footnotesize \textit{Lesson:} Distill three empirically-successful methods into one minimal expansion assumption and prove it exactly sufficient --- identifying minimal structural property derives exact guarantees.}
  \item \textbf{[Reject]} \texttt{ICLR\_2025\_1548} \textemdash\ Derives a phase transition only inside an idealized spiked-covariance proportional-limit regime, so the formalization yields a narrow equivalence rather than a structural insight that survives outside the assumption.
  \\ \textcolor{gray}{\footnotesize \textit{Lesson:} A phase-transition derived only inside an idealized spiked-covariance proportional limit yields narrow equivalence rather than structural insight surviving outside the assumption.}
  \item \textbf{[Reject]} \texttt{ICLR\_2024\_0923} \textemdash\ Replaces a non-convex objective with a maximal-correlation constraint but stops at reformulation --- it never produces the closed-form prediction or exact equivalence that distinguishes successful papers in this cluster.
  \\ \textcolor{gray}{\footnotesize \textit{Lesson:} Replacing a non-convex objective with a maximal-correlation constraint stops at reformulation --- closed-form prediction or exact equivalence is what the cluster requires.}
\end{itemize}
\end{tcolorbox}

\begin{tcolorbox}[colback=bgskill!50, colframe=cgreen!70, title={\textbf{C10}: Adaptive Dual-Guarantee Robust Objective Design \textcolor{gray}{\small (O8/H1/R4, conf: \textbf{medium})}}, breakable, before skip=4pt, after skip=4pt]
\textbf{Tactical pattern.} The shared move is to engineer a SINGLE objective or algorithm whose internal mechanism --- adaptive regularizer curvature, dual/Lagrangian reformulation, soft distance-penalized uncertainty set, or prediction-plus-fallback meta-procedure --- provably yields optimality in BOTH a favorable regime (stochastic, predictable, low-noise) and an unfavorable regime (adversarial, worst-case, high-variance) without needing to know in advance which regime is active. Concretely: ICML\_2023\_0419 picks Tsallis/Shannon entropy whose curvature lets FTRL exploit stochastic gaps yet retain adversarial regret; NeurIPS\_2025\_1933 derives the exact uncertainty set for which an entropy heuristic is the Lagrangian dual of a DRO problem; ICML\_2023\_0523 swaps a hard uncertainty-set minimax for soft distance-penalized satisficing; ICLR\_2025\_1432 transports the flat-minima/generalization correspondence into a flatness/transition-robustness theorem; ICML\_2023\_0453 reuses exponential smoothing as a data-free PAC-Bayes prior so bounds stay valid for unbounded importance weights. The verbs are "design regularizer whose curvature switches behavior," "reverse-engineer the implicit DRO problem," "replace hard set with soft penalty," and "interpolate via duality" --- not "add an auxiliary loss" or "pick a worst case to optimize."

\textbf{Differentiation.} \textit{The closest sibling is cluster 39 (Objective Realignment for Robustness via Geometric Targeting), which also reformulates training objectives toward a robustness criterion --- but cluster 39 picks ONE geometric target and aligns the loss to it (Bellman infinity-error replacing MSE in ICML\_2024\_1135; PAGs as a standalone objective in ICML\_2023\_0445), accepting the conservatism that comes with it. Cluster 10 explicitly refuses the single-target framing: its papers demand that the SAME algorithm be tight under stochastic AND adversarial regimes simultaneously (ICML\_2023\_0419), or that "robust" not collapse to "pessimistic" (ICML\_2023\_0523's satisficing). It also differs from cluster 30 (Objective Auditing via Explicit Constraint Injection), which bolts on auxiliary statistical terms (variance, covariance, causal invariance) to enforce one suppressed property; cluster 10's reformulations do not add a new constraint --- they redesign the regularizer/uncertainty-set geometry so dual-regime optimality falls out of the math itself.}

\textbf{When to pick.} Pick this tactic when the problem cannot be assumed to live in one regime --- you need provable performance whether the environment turns out to be stochastic-and-benign or adversarial-and-worst-case, and you cannot tune two separate algorithms because the regime is unknown at deployment. The deciding signal is the presence of a competing-regime axis (online/offline, adversarial/stochastic, prediction-trustworthy/prediction-broken) plus a hard requirement that one algorithm dominate across both ends.

\textbf{Tactical failure.} Reject papers in this cluster invoke the dual-guarantee FRAME but ship off-the-shelf compositions instead of inventing a new adaptive mechanism. ICLR\_2023\_0304 ports online-learning regret machinery into the robust MDP setting and calls it best-of-both --- but the algorithm is just FTRL versus the worst-case transition, with no regularizer whose behavior actually adapts to observed structure, so reviewers saw a routine combination rather than a novel dual-regime objective. ICLR\_2024\_0919 wraps an online knapsack solver with prediction-augmented + robust-fallback partitioning, which is the canonical learning-augmented template; the consistency-robustness frontier for fractional knapsack was already largely mapped, leaving the contribution incremental. NeurIPS\_2025\_1941 extends PAC-Bayes to Markov chains via mixing time --- a standard concentration-bound move that produces a single-regime certificate, not a mechanism that interpolates regimes. NeurIPS\_2025\_1853 tightens deep actor-critic risk bounds using validation rollouts, again a single-regime tightening rather than dual-regime optimality. The shared failure: borrowing the rhetoric of "best of both" or "robust + adaptive" without designing a regularizer, dual, or geometric correspondence whose mathematics organically delivers the dual guarantee.

\textbf{Examples.}
\begin{itemize}[leftmargin=12pt,topsep=2pt,itemsep=1pt]
  \item \textbf{[Oral]} \texttt{ICML\_2023\_0419} \textemdash\ Paradigmatic instance: a single entropy regularizer whose curvature plus adaptive step size proves polylog regret under stochastic feedback AND root-T regret under adversarial feedback --- exactly the dual-regime adaptive-regularizer move that defines this cluster.
  \\ \textcolor{gray}{\footnotesize \textit{Lesson:} A single entropy regularizer with curvature + adaptive step size proving polylog regret in stochastic AND root-T regret in adversarial --- one mechanism, two regimes, no separate worst-case program.}
  \item \textbf{[Oral]} \texttt{NeurIPS\_2025\_1933} \textemdash\ Reverse-engineers the implicit DRO Lagrangian-dual that a heuristic state-entropy regularizer is secretly solving, recovering the precise uncertainty sets for which it is robust --- the duality-as-bridge tactic central to the cluster.
  \\ \textcolor{gray}{\footnotesize \textit{Lesson:} Reverse-engineer the implicit DRO Lagrangian-dual that a heuristic regularizer is secretly solving --- recover the precise uncertainty sets for which it is robust.}
  \item \textbf{[Oral]} \texttt{ICML\_2023\_0522} \textemdash\ Identifies a structural decoupling in the dual of the budget-pacing problem that allows a single sample to yield a worst-case-robust policy, instantiating the dual-characterization route to dual guarantees.
  \\ \textcolor{gray}{\footnotesize \textit{Lesson:} Identify a structural decoupling in the dual of the budget-pacing problem --- a single sample yields a worst-case-robust policy by exploiting dual structure.}
  \item \textbf{[Oral]} \texttt{ICLR\_2025\_1432} \textemdash\ Establishes a flatness-in-policy-parameters $\leftrightarrow$ transition-perturbation-robustness correspondence, importing supervised generalization geometry as the mechanism that delivers robustness without a separate worst-case program.
  \\ \textcolor{gray}{\footnotesize \textit{Lesson:} Establish flatness-in-policy-parameters $\leftrightarrow$ transition-perturbation-robustness --- import supervised generalization geometry as the mechanism delivering robustness without separate worst-case programs.}
  \item \textbf{[Reject]} \texttt{ICLR\_2023\_0304} \textemdash\ Frames online robust MDPs as an adversarial game with bandit exploration but the resulting algorithm has no adaptive regularizer or dual reformulation that organically yields dual-regime optimality --- it is online-learning machinery applied to robust MDPs.
  \\ \textcolor{gray}{\footnotesize \textit{Lesson:} Online robust MDPs framed as adversarial games with bandit exploration --- applying online-learning machinery without a new adaptive regularizer fails the dual-regime mechanism criterion.}
  \item \textbf{[Reject]} \texttt{ICLR\_2024\_0919} \textemdash\ Uses the canonical prediction-augmented-plus-robust-fallback partitioning template for online knapsack rather than inventing a new mechanism, leaving the consistency-robustness contribution incremental on an already-mapped frontier.
  \\ \textcolor{gray}{\footnotesize \textit{Lesson:} Prediction-augmented + robust-fallback partitioning is the canonical template --- incremental on an already-mapped frontier without a new mechanism.}
  \item \textbf{[Reject]} \texttt{NeurIPS\_2025\_1941} \textemdash\ Tightens PAC-Bayes via Markov-chain mixing time --- a single-regime concentration improvement rather than a mechanism that interpolates between competing regimes, missing the dual-guarantee design move.
  \\ \textcolor{gray}{\footnotesize \textit{Lesson:} PAC-Bayes tightening via Markov-chain mixing is single-regime concentration improvement --- missing the dual-guarantee design move that interpolates between regimes.}
\end{itemize}
\end{tcolorbox}

\begin{tcolorbox}[colback=bgskill!50, colframe=cgreen!70, title={\textbf{C12}: Structural Sensitivity Analysis for Mechanism Targeting \textcolor{gray}{\small (O21/H0/R6, conf: \textbf{high})}}, breakable, before skip=4pt, after skip=4pt]
\textbf{Tactical pattern.} Re-derive the sensitivity calculus of a differentially private pipeline to identify the *operation* --- not the raw data --- where noise should be injected or budget allocated. Concretely: locate the actual privacy-sensitive step (merge of sketches, policy update, multi-epoch matrix factorization, search response, gating selection, hyperparameter loop, foundation-model adapter), formalize its sensitivity (often jointly across passes/queries/users rather than per-record), then design a mechanism whose noise scales with that re-targeted sensitivity rather than with the conventional per-observation surrogate. Variants include moving noise from individual observations to aggregate decisions (ICML\_2024\_1103), from full-model gradients to a low-dimensional adapter subspace (ICML\_2024\_1144), from per-pass to joint-pass factorization (ICML\_2023\_0490), from post-processing to merge operations (ICML\_2023\_0536), and from sequential to parallel auditing within one run (NeurIPS\_2023\_0708, ICML\_2025\_1677). The shared verb is "re-target the mechanism" backed by a tight new sensitivity bound.

\textbf{Differentiation.} \textit{Among siblings under "Audit and Relax Implicit Assumption," the closest neighbor is cluster 30 (Objective Auditing via Explicit Constraint Injection), which audits the *training loss* and adds a regularizer to enforce a suppressed property. This cluster is structurally different: it audits the *mechanism's sensitivity profile* --- where in the computation an individual record actually leaks --- and relocates the noise/budget injection point, leaving the underlying objective untouched. It also differs from cluster 39 (Objective Realignment for Robustness via Geometric Targeting), which retargets a *robustness criterion*; here the target is privacy accounting, and the move is operational (which step is the mechanism applied to) rather than criterion-level. Unlike cluster 0 (identifiability via structural constraints), the structure being exploited is the sensitivity graph of a stochastic mechanism, not an information-theoretic identifiability source.}

\textbf{When to pick.} Pick this instantiation when your problem already has a working but loose privacy mechanism applied at a "default" location (per-sample gradient, per-query response, per-pass noise) and you can identify a downstream aggregation, composition, or response operation whose per-individual sensitivity is provably lower or qualitatively different. Signal: the conventional analysis treats some pipeline stage as "post-processing" or "public" when it actually consumes private information, or the response type is structurally richer (vector/solution/stream) than the scalar setting the existing bound was proven for.

\textbf{Tactical failure.} Papers attempting this tactic fail when they either (a) point out a sensitivity gap without delivering a new tight mechanism, or (b) repurpose DP machinery for a different goal without re-deriving the sensitivity claim that justifies the move. NeurIPS\_2024\_1164 catalogues subtle accounting mistakes in subsampled composition but functions as a clarification/correction rather than producing a new mechanism with strictly improved utility --- reviewers see it as an audit note, not a methodological contribution. NeurIPS\_2023\_0576 proposes a linear-scaling heuristic for DP hyperparameter cost but relies on extrapolation from non-DP scaling laws rather than a formal sensitivity analysis of the tuning loop, so the privacy-cost saving is empirical rather than provable. NeurIPS\_2024\_1249 reinterprets the exponential mechanism on MoE gating as LDP-flavored regularization, but the contribution is a generalization-bound view rather than a tighter privacy-utility frontier, leaving the "structural sensitivity" insight without an operational payoff. ICLR\_2023\_0277 (MaSS) and ICLR\_2023\_0235 (Fed-Cor) likewise relocate where suppression/aggregation happens but lean on minimax training or shared-transformation engineering instead of a closed-form sensitivity bound, so the privacy guarantee remains heuristic.

\textbf{Examples.}
\begin{itemize}[leftmargin=12pt,topsep=2pt,itemsep=1pt]
  \item \textbf{[Oral]} \texttt{ICML\_2024\_1103} \textemdash\ Canonical example of relocating the noise from individual observations to the aggregate policy update --- the per-individual sensitivity of the aggregate is provably lower, and the paper re-derives the privacy-utility trade-off at this new injection point.
  \\ \textcolor{gray}{\footnotesize \textit{Lesson:} Relocate noise from individual observations to aggregate policy update --- per-individual sensitivity of the aggregate is provably lower; re-derive privacy-utility at the new injection point.}
  \item \textbf{[Oral]} \texttt{ICML\_2023\_0490} \textemdash\ Reformalizes sensitivity jointly across multiple epochs as an adaptive stream rather than treating each pass independently, then designs a matrix-factorization mechanism that exploits this joint sensitivity --- tactically the cluster's signature move.
  \\ \textcolor{gray}{\footnotesize \textit{Lesson:} Reformalize sensitivity jointly across multiple epochs as adaptive stream --- design matrix-factorization mechanisms exploiting joint sensitivity, not per-pass independence.}
  \item \textbf{[Oral]} \texttt{ICML\_2023\_0536} \textemdash\ Identifies that the merge operation on DP sketches --- assumed to be private post-processing --- actually leaks information, and redesigns merge itself as a privacy mechanism with its own calibrated noise.
  \\ \textcolor{gray}{\footnotesize \textit{Lesson:} Show the merge operation on DP sketches (assumed post-processing) actually leaks --- redesign merge itself as a privacy mechanism with calibrated noise.}
  \item \textbf{[Oral]} \texttt{NeurIPS\_2023\_0689} \textemdash\ Shows that a combinatorial problem's intrinsic low-sensitivity structure makes generic composition unnecessary, and constructs a direct mechanism whose noise targets the actual sensitivity rather than a worst-case surrogate.
  \\ \textcolor{gray}{\footnotesize \textit{Lesson:} Combinatorial intrinsic low-sensitivity makes generic composition unnecessary --- construct direct mechanisms targeting actual sensitivity, not worst-case surrogates.}
  \item \textbf{[Reject]} \texttt{NeurIPS\_2024\_1164} \textemdash\ Surfaces real sensitivity-accounting errors in composed subsampled mechanisms but stops at correction rather than producing a new mechanism with tighter, re-targeted sensitivity --- diagnostic without operational payoff.
  \\ \textcolor{gray}{\footnotesize \textit{Lesson:} Surfacing sensitivity-accounting errors without producing a new tighter mechanism is diagnostic without operational payoff.}
  \item \textbf{[Reject]} \texttt{NeurIPS\_2023\_0576} \textemdash\ Proposes a scaling rule that lets hyperparameter search be done at small scale to save privacy budget, but the saving is an empirical extrapolation rather than a formal sensitivity bound on the tuning loop.
  \\ \textcolor{gray}{\footnotesize \textit{Lesson:} An empirical scaling rule for HP-search saving is extrapolation, not a formal sensitivity bound on the tuning loop.}
  \item \textbf{[Reject]} \texttt{NeurIPS\_2024\_1249} \textemdash\ Reinterprets exponential-mechanism gating as LDP regularization and derives generalization bounds, but the move yields a learning-theory result rather than a tighter privacy-utility mechanism, missing the cluster's operational payoff.
  \\ \textcolor{gray}{\footnotesize \textit{Lesson:} Reinterpreting exponential-mechanism gating as LDP regularization yields a learning-theory result rather than a tighter privacy-utility mechanism --- wrong methodology family.}
\end{itemize}
\end{tcolorbox}

\begin{tcolorbox}[colback=bgskill!50, colframe=cgreen!70, title={\textbf{C19}: Relational Structure Critique and Realignment \textcolor{gray}{\small (O2/H2/R14, conf: \textbf{low})}}, breakable, before skip=4pt, after skip=4pt]
\textbf{Tactical pattern.} The specific move is to locate where a model's STRUCTURAL/topological encoding of relationships (the graph proxy chosen, the message-passing rule, the attention scope, the connectivity of pipeline components) is misaligned with the data's actual relational organization, and then surgically replace or constrain that encoding while leaving the learning objective largely intact. Concrete instantiations: substitute a verbose structural proxy with a leaner relational graph (data-flow graph replacing AST in GraphCodeBERT, ICLR\_2021\_0029), bound an over-broad attention range to local neighborhoods (Graph Transformers' over-globalizing fix in ICML\_2024\_1056), inject a learnable per-edge linear map to prevent dimensional collapse in aggregation (sheaf neural nets in ICLR\_2025\_1400), or wire previously independent components into a single relationally-aware pipeline (text encoder prompted with graph neighborhoods in NeurIPS\_2023\_0713). The verb is "rewire the relational encoding"; the object is edges, neighborhoods, or topological scope --- not the loss.

\textbf{Differentiation.} \textit{Unlike Sibling 30 (Objective Auditing via Explicit Constraint Injection), which leaves architecture intact and adds a regularizer/auxiliary loss to inject a missing statistical property, Cluster 19 keeps the objective largely standard and instead modifies the STRUCTURAL INPUT or PROPAGATION RULE the model receives --- what graph it sees, how far attention reaches, whether edges carry transformations. Unlike Sibling 34 (Generative Process Realignment via Structural Assumption Audit), which audits forward/reverse process design in generative pipelines, Cluster 19 targets discriminative/representational architectures where relational topology between entities is the unit of audit. And unlike Sibling 9 (Structural Inductive Bias Formalization), which formalizes a latent structural property to derive a predictive theory, Cluster 19 takes an engineering action --- swapping a graph proxy or constraining a propagation operator --- without requiring formal characterization as the deliverable.}

\textbf{When to pick.} Pick this tactical instantiation when the failure of the existing method traces specifically to the topology/scope of its relational encoding (wrong graph proxy, too-global aggregation, missing inter-component relational link), rather than to its loss, its independence assumption, or its generative direction. The diagnostic signal is that you can point at an edge set, an attention mask, or a message-passing radius and say "this is the wrong structural object for the data."

\textbf{Tactical failure.} Reject papers in this cluster cluster around four specific failure patterns. (1) Cross-domain transplant without demonstrating structural necessity: ICML\_2025\_1827 imports the recommender Wide \& Deep architecture into node classification by analogy alone, and ICLR\_2025\_1400 transplants sheaf-theoretic restriction maps into recommendation without isolating that the dimensional-collapse story is the actual bottleneck --- reviewers cannot tell whether the structural realignment matters or any expressivity bump would suffice. (2) Composite structural fix where the proposed relational change is entangled with several auxiliary tweaks, leaving no ablation that isolates the relational mechanism: ICLR\_2024\_0798 (correlated dense associative memories) explicitly notes the modification "appears as a routine extension" without isolating ablations, and ICLR\_2024\_0776 stacks ranking + attribute features + hierarchical partitioning so the relational claim is unidentifiable. (3) Benchmark/evaluation reframing without a compelling structural problem driving it --- ICLR\_2024\_0932 (recent link classification) and ICLR\_2024\_0917 (overlap-based generative evaluation) propose new relational task formulations but reviewers see them as gap-filling rather than exposing a structural defect. (4) Borrowed scaffolds (knowledge graphs over features in ICLR\_2023\_0357, learned constraint functions in ICLR\_2023\_0239, differentiable coarsening in ICLR\_2023\_0270) where the auxiliary relational structure is plausible but its source-of-truth is unverified, so the structural proxy itself becomes the new unjustified assumption.

\textbf{Examples.}
\begin{itemize}[leftmargin=12pt,topsep=2pt,itemsep=1pt]
  \item \textbf{[Oral]} \texttt{ICML\_2024\_1056} \textemdash\ Exemplifies constraining an over-broad relational scope: formalizes that softmax-based global attention systematically dilutes locally-concentrated signal in graphs and bounds attention range --- a structural surgery on the propagation operator, not the loss.
  \\ \textcolor{gray}{\footnotesize \textit{Lesson:} Formalize that softmax-based global attention dilutes locally-concentrated graph signal and bound attention range --- structural surgery on the propagation operator, not the loss.}
  \item \textbf{[Oral]} \texttt{ICLR\_2023\_0375} \textemdash\ Exemplifies redesigning the message-passing rule along the relational edge: decomposes propagated content into class-invariant and class-variant components so structurally incompatible neighbors share only the safe channel --- modifying what flows along edges rather than which edges exist.
  \\ \textcolor{gray}{\footnotesize \textit{Lesson:} Decompose propagated content into class-invariant and class-variant components --- modify what flows along edges rather than which edges exist.}
  \item \textbf{[Reject]} \texttt{ICML\_2025\_1827} \textemdash\ Cross-domain transplant failure: borrows Wide \& Deep from recommendation by analogy without demonstrating that the memorization/generalization split corresponds to a real structural defect in node classification.
  \\ \textcolor{gray}{\footnotesize \textit{Lesson:} Borrowing Wide \& Deep from recommendation by analogy without demonstrating the memorization/generalization split corresponds to a real structural defect in the target setting fails.}
  \item \textbf{[Reject]} \texttt{ICLR\_2024\_0798} \textemdash\ Composite-fix failure: pairs excitatory and inhibitory associative rules to encode item correlations but lacks ablations isolating that the independence assumption --- rather than capacity or stability --- is the binding structural mismatch.
  \\ \textcolor{gray}{\footnotesize \textit{Lesson:} Composite excitatory/inhibitory associative rules need ablations isolating that the independence assumption --- not capacity or stability --- is the binding mismatch.}
  \item \textbf{[Reject]} \texttt{ICLR\_2025\_1400} \textemdash\ Imported-machinery failure: applies sheaf-theoretic restriction maps to recommender oversmoothing without isolating whether sheaf structure specifically (versus any per-edge transformation) is what the data demands.
  \\ \textcolor{gray}{\footnotesize \textit{Lesson:} Sheaf-theoretic restriction maps need to isolate whether sheaf structure specifically (versus any per-edge transformation) is what the data demands.}
\end{itemize}
\end{tcolorbox}

\begin{tcolorbox}[colback=bgskill!50, colframe=cgreen!70, title={\textbf{C30}: Objective Auditing via Explicit Constraint Injection \textcolor{gray}{\small (O9/H6/R16, conf: \textbf{high})}}, breakable, before skip=4pt, after skip=4pt]
\textbf{Tactical pattern.} Papers in this cluster diagnose a specific structural property (decorrelation, expert orthogonality, sharpness, ensemble disagreement, calibrated abstention, missing-data propensity) that the standard training loss either suppresses, ignores, or only enforces implicitly through architectural side-effects, then **inject an explicit, named auxiliary term** --- a regularizer, an adversarial min-max term, an inverse-propensity reweighting, or a second loss head --- that directly optimizes that property alongside the original objective. The verbs are concrete: "add orthogonality loss between experts" (1836), "add disagreement loss on unlabeled OOD" (0166), "decompose non-degeneracy into variance/covariance/invariance terms" (0148), "reformulate objective as min-max over a local neighborhood" (0067), "treat label availability as missing-data and reweight by IPW" (0414), "keep the pretraining LM head live as an auxiliary regularizer during RM training" (1289). Crucially, the injected term is a *first-class* optimization target, not a hyperparameter knob --- papers usually prove or demonstrate that without the term the property collapses, and ablate to show the term is what closes the gap.

\textbf{Differentiation.} \textit{Versus sibling cluster 39 (Objective Realignment for Robustness via Geometric Targeting), which **replaces** the loss to match a geometric/worst-case robustness criterion (e.g., MSE$\to$infinity-error for Q-learning), this cluster **adds onto** the existing loss with a new term enforcing a previously-implicit property --- the original objective stays in place. Versus sibling 34 (Generative Process Realignment), which restructures the *architecture* of forward/reverse generative processes (unify or decouple components), this cluster leaves the model architecture and pipeline intact and changes only the loss function. Versus sibling 0 (Structural Constraint Exploitation for Identifiability), which exploits already-present structural constraints to derive *identifiability* theorems, this cluster typically delivers an empirically optimizable regularizer rather than a non-identifiability proof. Versus sibling 9 (Structural Inductive Bias Formalization), which is purely theoretical characterization of latent properties, this cluster is action-oriented: the audit's deliverable is a new term in a real loss.}

\textbf{When to pick.} Pick this tactical instantiation when the property you want is already (weakly) emergent from existing training but only as a side-effect of architectural choices, initializations, or data idiosyncrasies --- and you can write that property down as a measurable statistic on representations, gradients, or output distributions. Choose a sibling tactic instead if the audit's payoff is an identifiability theorem (cluster 0) or if the geometry of the loss itself, not a missing term, is the binding constraint (cluster 39).

\textbf{Tactical failure.} The dominant failure mode is **injecting a plausible-looking auxiliary term without proving it is the binding constraint** that closes the observed gap. ICLR\_2025\_1644 (VC regularization) proposes a VICReg-style decorrelation regularizer in a supervised setting but reviewers cannot tell whether the suppressed statistical property --- rather than implicit capacity expansion --- is what drives the gain. ICLR\_2024\_0767 (AdaFlood) replaces a global flood constant with a per-instance one, an obvious natural extension whose ablations do not isolate the heterogeneity-uniformity assumption from confounds like data quality or model capacity. ICML\_2025\_1722 (House of Cards) builds an asymmetric-ablation diagnostic and a curriculum dropout fix, but the downstream gains may come from generic regularization rather than the targeted dominant-parameter mechanism. ICLR\_2025\_1566 (Causal Invariant BNN) reformulates the objective to encode an ignored structural assumption but the failure mode appears only on bespoke metrics designed to surface it, leaving open whether it bites in normal evaluation. NeurIPS\_2024\_1178 (Collision CE) replaces standard CE with a probability-of-agreement loss on principled grounds but does not delineate when its zero-gradient-on-uncertain-targets behavior helps versus hurts. The common thread: reviewers reject when the injected constraint reads like a **post-hoc add-on** that could explain almost any improvement, rather than the *unique* term whose removal provably collapses the targeted property.

\textbf{Examples.}
\begin{itemize}[leftmargin=12pt,topsep=2pt,itemsep=1pt]
  \item \textbf{[Oral]} \texttt{NeurIPS\_2025\_1836} \textemdash\ Cleanly exemplifies the tactic: identifies that MoE's load-balancing auxiliary loss actively suppresses expert specialization, then injects two named counter-terms (orthogonality between expert representations, variance over routing) that directly enforce the property the original auxiliary undermined.
  \\ \textcolor{gray}{\footnotesize \textit{Lesson:} MoE's load-balancing aux loss suppresses expert specialization; inject orthogonality + variance counter-terms to enforce the property the original auxiliary undermined.}
  \item \textbf{[Oral]} \texttt{ICLR\_2023\_0166} \textemdash\ Adds a structurally distinct second loss --- agreement on labeled data plus an explicit disagreement loss on unlabeled OOD --- making feature diversity an active optimization target rather than a hoped-for ensemble side-effect.
  \\ \textcolor{gray}{\footnotesize \textit{Lesson:} Agreement on labeled + explicit disagreement on unlabeled OOD --- feature diversity becomes an active optimization target, not a hoped-for ensemble side-effect.}
  \item \textbf{[Oral]} \texttt{ICML\_2023\_0414} \textemdash\ Imports inverse-propensity weighting from missing-data statistics as an explicit reweighting term, converting SSL's silent MAR assumption into a testable propensity model that augments the standard supervised loss.
  \\ \textcolor{gray}{\footnotesize \textit{Lesson:} Import inverse-propensity weighting from missing-data statistics --- convert SSL's silent MAR assumption into a testable propensity model.}
  \item \textbf{[Oral]} \texttt{ICLR\_2022\_0094} \textemdash\ Diagnoses that joint fine-tuning's silent assumption (random head can co-adapt with backbone) distorts pretrained features, then enforces 'preserve representations' as an explicit two-stage procedure (LP then FT) --- constraint injected via training schedule rather than loss term.
  \\ \textcolor{gray}{\footnotesize \textit{Lesson:} LP-then-FT staged procedure --- constraint injected via training schedule, not a loss term, to preserve pretrained features.}
  \item \textbf{[Reject]} \texttt{ICLR\_2025\_1644} \textemdash\ Imports VICReg-style variance-covariance regularization into supervised representation learning but cannot demonstrate that decorrelation is the binding constraint rather than generic capacity-side regularization.
  \\ \textcolor{gray}{\footnotesize \textit{Lesson:} VICReg-style decorrelation must be shown to be the binding constraint, not generic capacity-side regularization.}
  \item \textbf{[Reject]} \texttt{ICLR\_2024\_0767} \textemdash\ Replaces a uniform flood constant with a per-instance adaptive parameter --- exactly the constraint-injection move --- but lacks ablations isolating the uniformity-assumption violation from confounds in capacity or data difficulty.
  \\ \textcolor{gray}{\footnotesize \textit{Lesson:} Per-instance adaptive flooding is the constraint-injection move --- ablations must isolate the uniformity-assumption violation from capacity / data-difficulty confounds.}
  \item \textbf{[Reject]} \texttt{NeurIPS\_2024\_1178} \textemdash\ Substitutes standard CE with a principled collision (probability-of-agreement) objective derived from auditing how CE penalizes uncertain targets, but fails to delineate the regime where the injected behavior dominates over plain CE.
  \\ \textcolor{gray}{\footnotesize \textit{Lesson:} Probability-of-agreement objective derived from auditing CE needs to delineate the regime where injected behavior dominates over plain CE.}
\end{itemize}
\end{tcolorbox}

\begin{tcolorbox}[colback=bgskill!50, colframe=cgreen!70, title={\textbf{C34}: Generative Process Realignment via Structural Assumption Audit \textcolor{gray}{\small (O0/H5/R5, conf: \textbf{low})}}, breakable, before skip=4pt, after skip=4pt]
\textbf{Tactical pattern.} The specific move is auditing the *structural decomposition of the generative pipeline itself* --- the partition between forward and reverse processes, between prior model and observation/constraint module, between separate task-specific objectives, or between conditioning bridge and core generator --- and then either collapsing components that the field treats as necessarily separate into a single joint process, or splitting components the field treats as necessarily fused so they can be recombined at inference. Concretely: replace an uninformative terminal prior with the actual conditional input so the forward SDE's endpoint matches the test-time input (ICML\_2023\_0470); concatenate condition and target and selectively corrupt only the target portion to internalize the conditioning bridge (ICLR\_2023\_0204); train one prior unconditionally and inject task-specific likelihoods only at sampling time (ICLR\_2022\_0135); inject the constrained input mid-trajectory of an existing forward process instead of training a new conditional model (ICLR\_2022\_0133); collapse N specialized conditional objectives into a single masked-infilling objective (NeurIPS\_2023\_0758). The verbs are *unify*, *decouple*, *re-terminate*, *re-initialize* --- operating on the architecture of the generative process, not on its loss.

\textbf{Differentiation.} \textit{Sibling cluster 30 (Objective Auditing via Explicit Constraint Injection) audits the *training objective* and adds auxiliary loss terms or statistical regularizers (VICReg-style variance/covariance terms, causal invariance penalties) --- the generative or representation pipeline stays structurally intact and only the scalar loss changes. Cluster 34 instead leaves the loss type alone (it is still denoising score-matching, flow-matching, or masked prediction) and surgically modifies the *process structure* --- which endpoint the forward SDE attracts to (ICML\_2023\_0470), where conditioning enters (ICLR\_2023\_0204), whether prior and likelihood share parameters (ICLR\_2022\_0135), or whether multiple specialized objectives can be merged into one (NeurIPS\_2023\_0758). Sibling cluster 19 (Relational Structure Critique) is bound to graph/relational encodings specifically, whereas cluster 34 is bound to *probabilistic generative dynamics* (diffusion forward processes, flow trajectories, masked-conditioning sequences). Sibling cluster 39 (Geometric Targeting for Robustness) realigns objectives toward worst-case geometry; cluster 34 realigns the generative *trajectory* toward task-natural boundary conditions, which is a different axis of mismatch.}

\textbf{When to pick.} Pick this tactic when the bottleneck is provably the *structural plumbing* of an existing generative pipeline --- e.g., a hand-designed bridge module, a fixed Gaussian terminal prior, or a per-task conditional retraining loop --- rather than the loss function or the encoder. The signal is concrete: you can name a specific component of the forward/reverse/conditioning pipeline that exists for historical reasons and whose removal or merger admits a closed-form or zero-shot reformulation.

\textbf{Tactical failure.} The dominant failure mode for cluster-34 rejects is that the "audit" is post-hoc framing for what is really a *bolt-together* of two existing generative pipelines, with no demonstration that the audited structural separation was causally limiting. ICLR\_2025\_1597 unifies two peptide prediction targets into a joint generative framework but admits the gap only surfaces when measuring the neglected property directly --- competitive aggregate benchmarks made the unification cosmetically motivated rather than necessary. ICLR\_2023\_0324 trains separate generative priors for signal and structured noise and couples them at inference, but the move reads as a straightforward extension of score-based posterior sampling rather than a re-derivation revealing why prior decoupling was wrong. ICML\_2025\_1790 reframes structured missingness as conditional generation across two complementary sources, but uses a flow-matching backbone as an off-the-shelf fusion module --- reviewers read the dual-task breadth as compensating for incremental methodology. ICLR\_2025\_1516 stages contrastive alignment then generative conditioning on an EEG pipeline, but never identifies which structural assumption was violated; it only inserts a generation head where there used to be a classification head. ICLR\_2025\_1615 (TimeDiT) jointly formulates two objectives but the claimed failure mode only appears at heterogeneous large scales not exercised in standard evaluation. The shared pathology: the audit step is missing or rhetorical --- papers describe a "joint framework" without showing that the prior decomposition was load-bearing or that the unification produces a property unattainable by either component alone.

\textbf{Examples.}
\begin{itemize}[leftmargin=12pt,topsep=2pt,itemsep=1pt]
  \item \textbf{[Oral]} \texttt{ICML\_2023\_0470} \textemdash\ Most exemplary instance of re-terminating a forward generative process at the actual conditional input rather than an arbitrary uninformative prior, exposing that the terminal distribution of the SDE was a hidden design variable yielding closed-form supervision.
  \\ \textcolor{gray}{\footnotesize \textit{Lesson:} Re-terminate a forward generative process at the actual conditional input rather than an arbitrary uninformative prior --- terminal distribution was a hidden design variable yielding closed-form supervision.}
  \item \textbf{[Oral]} \texttt{ICLR\_2023\_0204} \textemdash\ Cleanly demonstrates the unification half of the tactic: dissolves the separately-trained conditioning bridge by co-encoding condition and target and corrupting only the target portion, internalizing guidance into the core diffusion process.
  \\ \textcolor{gray}{\footnotesize \textit{Lesson:} Co-encode condition and target, corrupt only the target portion --- internalize guidance into the core diffusion process, dissolve the separate conditioning bridge.}
  \item \textbf{[Oral]} \texttt{ICLR\_2022\_0135} \textemdash\ Cleanly demonstrates the decoupling half of the tactic: severs the implicit fusion of prior model and observation operator that paired-supervised training had welded together, training the prior once and injecting variant-specific constraints only at inference.
  \\ \textcolor{gray}{\footnotesize \textit{Lesson:} Sever the welded prior-model + observation-operator pairing trained-supervised pipelines created --- train the prior once, inject variant-specific constraints only at inference.}
  \item \textbf{[Oral]} \texttt{NeurIPS\_2023\_0758} \textemdash\ Audits whether explicit task labels assumed necessary across a family of speech-generation tasks are redundant given surrounding context, then collapses all specialized objectives into one masked-infilling process --- a textbook 'unify previously separate components' move.
  \\ \textcolor{gray}{\footnotesize \textit{Lesson:} Audit whether explicit task labels are redundant given context, then collapse specialized objectives into one masked-infilling process --- unify previously separate components.}
  \item \textbf{[Reject]} \texttt{ICLR\_2023\_0324} \textemdash\ The 'audit' is rhetorical --- paper proposes joint generative modeling of signal and structured noise as two priors coupled at inference, but does not establish that the prior decomposition assumption was actually load-bearing rather than a routine extension of posterior sampling.
  \\ \textcolor{gray}{\footnotesize \textit{Lesson:} Joint generative modeling of signal and structured noise as two priors needs to establish that the prior decomposition was load-bearing, not a routine extension.}
  \item \textbf{[Reject]} \texttt{ICLR\_2025\_1597} \textemdash\ Unifies two prediction targets into a joint generation framework but concedes the omitted property only matters when measured directly --- aggregate benchmarks already looked competitive, making the structural unification cosmetically rather than causally motivated.
  \\ \textcolor{gray}{\footnotesize \textit{Lesson:} Joint generation that helps only when measured directly is cosmetic --- aggregate benchmarks already looking competitive means the unification isn't causally motivated.}
  \item \textbf{[Reject]} \texttt{ICML\_2025\_1790} \textemdash\ Reuses a recent flow-matching backbone to fuse two complementary remote-sensing sources by reframing missingness as conditional generation, but the framing reads as application of an off-the-shelf generator rather than auditing a structural assumption that was actually limiting.
  \\ \textcolor{gray}{\footnotesize \textit{Lesson:} Reusing flow-matching to fuse two complementary sources by reframing missingness reads as application of an off-the-shelf generator, not auditing an actually-limiting structural assumption.}
\end{itemize}
\end{tcolorbox}

\begin{tcolorbox}[colback=bgskill!50, colframe=cgreen!70, title={\textbf{C39}: Objective Realignment for Robustness via Geometric Targeting \textcolor{gray}{\small (O9/H2/R10, conf: \textbf{high})}}, breakable, before skip=4pt, after skip=4pt]
\textbf{Tactical pattern.} Identify a specific geometric or landscape quantity (loss sharpness, Hessian/low-rank condition number, gradient-input alignment, distance to decision boundary, infinity-norm of Bellman residual) that mechanistically controls a robustness criterion, then either (a) replace the standard average-case training loss with one that matches this geometric quantity directly, or (b) add an explicit regularizer that penalizes the quantity during normal training. Concrete moves include: swap squared-error Bellman for infinity-norm Bellman to match worst-case Q-error (ICML\_2024\_1135); penalize the low-rank core's condition number to bound adversarial Lipschitz constant (NeurIPS\_2025\_1865); reweight examples by PGD-steps-to-boundary to allocate capacity to vulnerable points (ICLR\_2021\_0025); promote the perceptually-aligned-gradient byproduct of adversarial training into a standalone training objective (ICML\_2023\_0445). The signature is "drop the proxy loss, optimize the geometric quantity itself."

\textbf{Differentiation.} \textit{Versus cluster 30 (Objective Auditing via Explicit Constraint Injection), which adds auxiliary statistical/distributional constraints (variance, covariance, causal invariance) to a representation-learning objective, this cluster targets *landscape geometry* --- curvature, condition number, gradient-direction alignment, boundary distance --- and the property being enforced is specifically a robustness certificate, not a representational desideratum. Versus cluster 10 (Adaptive Dual-Guarantee Robust Objective Design), which redesigns objectives to deliver simultaneous guarantees across competing regimes (stochastic vs. adversarial, online vs. offline), this cluster commits to a single worst-case criterion and realigns the loss to *match* that criterion geometrically rather than blending guarantees. Versus cluster 9 (Structural Inductive Bias Formalization), which characterizes existing structural properties to *explain* phenomena, this cluster uses the geometric characterization *prescriptively* --- turning it into a training-time penalty.}

\textbf{When to pick.} Pick this tactic when the user can name a robustness target (worst-case error, perturbation tolerance, post-compression sensitivity) AND can connect it to a single landscape/geometric quantity that the standard training loss does not directly minimize. If the gap is between objective and criterion is geometric/curvature-based rather than distributional or constraint-based, instantiate as cluster 39 rather than cluster 30 or 10.

\textbf{Tactical failure.} The dominant failure pattern in this cluster's rejects is *naming a candidate geometric or structural target without proving it is the load-bearing one*. ICLR\_2023\_0348 collapses adversarial training's bilevel game into a single-level NTK-based objective but cannot show the kernel approximation preserves the adversarial signal that made the original work --- geometric simplification destroys the property being targeted. ICLR\_2023\_0329 hypothesizes that predictive-coding's locality "should" limit adversarial perturbation propagation but never derives a quantitative link between locality and robustness, leaving the geometric claim speculative. ICLR\_2023\_0212 recalibrates batch-norm statistics conditioned on a perturbation classifier --- a surface intervention without any link to a landscape property that controls robustness. NeurIPS\_2024\_1263 and NeurIPS\_2023\_0633 reformulate noisy-label training as DRO/Rockafellian relaxation but cannot demonstrate the chosen uncertainty geometry is tighter than existing robust losses for the noise regime they target. The shared symptom: the paper proposes the realigned objective but fails to either (i) prove a bound linking the geometric quantity to the robustness criterion, or (ii) show the new objective strictly dominates the standard one outside the favorable instances cherry-picked.

\textbf{Examples.}
\begin{itemize}[leftmargin=12pt,topsep=2pt,itemsep=1pt]
  \item \textbf{[Oral]} \texttt{ICML\_2024\_1135} \textemdash\ Canonical instance of objective-criterion realignment: explicitly diagnoses that worst-case robust Q-learning was being trained with average-case (MSE) Bellman error, and proves that swapping in the infinity-norm Bellman residual yields tighter bounds --- the tactic stripped to its essence.
  \\ \textcolor{gray}{\footnotesize \textit{Lesson:} Worst-case robust Q-learning was being trained with average-case (MSE) Bellman error; swap to infinity-norm --- objective realignment in essential form.}
  \item \textbf{[Oral]} \texttt{ICML\_2025\_1756} \textemdash\ Translates an informal safety property (resistance to fine-tuning adaptation) into the Hessian condition number and minimizes that quantity directly during pretraining, exemplifying the cluster's signature move of converting a desired behavior into a precise landscape geometric quantity.
  \\ \textcolor{gray}{\footnotesize \textit{Lesson:} Translate fine-tuning resistance into the Hessian condition number and minimize that quantity directly --- convert desired behavior into a precise landscape geometric quantity.}
  \item \textbf{[Oral]} \texttt{ICLR\_2021\_0025} \textemdash\ Uses the geometric distance from each example to the decision boundary (proxied by PGD steps) as the realigned weighting signal, allocating capacity by vulnerability --- geometric targeting at the per-example level rather than at the global loss level.
  \\ \textcolor{gray}{\footnotesize \textit{Lesson:} Geometric distance from each example to decision boundary (proxied by PGD steps) as the realigned weighting signal --- geometric targeting at per-example level.}
  \item \textbf{[Oral]} \texttt{NeurIPS\_2025\_1865} \textemdash\ Derives a perturbation bound showing the low-rank core's condition number controls adversarial Lipschitz constant, then penalizes that exact spectral quantity during dynamical low-rank training --- a textbook 'identify the geometric quantity, regularize it directly' instance.
  \\ \textcolor{gray}{\footnotesize \textit{Lesson:} Derive that the low-rank core's condition number controls adversarial Lipschitz; penalize that exact spectral quantity --- identify the geometric quantity, regularize directly.}
  \item \textbf{[Reject]} \texttt{ICLR\_2023\_0348} \textemdash\ Attempted geometric reformulation (NTK-based collapse of bilevel adversarial objective into single-level) but lost the adversarial signal in the simplification, illustrating the failure of geometric realignment that does not preserve what made the original objective work.
  \\ \textcolor{gray}{\footnotesize \textit{Lesson:} NTK-based collapse of bilevel adversarial into single-level lost the adversarial signal --- geometric realignment must preserve what made the original objective work.}
  \item \textbf{[Reject]} \texttt{ICLR\_2023\_0329} \textemdash\ Proposes predictive-coding locality as a robustness mechanism but never derives a quantitative link from local error signals to bounded perturbation propagation --- the geometric story remains a hypothesis rather than a proven realignment.
  \\ \textcolor{gray}{\footnotesize \textit{Lesson:} Predictive-coding locality without quantitative link from local error signals to bounded perturbation propagation leaves the geometric story as hypothesis, not realignment.}
  \item \textbf{[Reject]} \texttt{ICLR\_2023\_0212} \textemdash\ Recalibrates batch-norm statistics conditioned on a perturbation-type classifier without identifying any landscape or geometric property that this targets, leaving the intervention engineering rather than principled objective realignment.
  \\ \textcolor{gray}{\footnotesize \textit{Lesson:} Recalibrate batch-norm conditioned on a perturbation-type classifier without identifying any landscape property --- engineering, not principled objective realignment.}
\end{itemize}
\end{tcolorbox}

\begin{tcolorbox}[colback=bgskill!50, colframe=cgreen!70, title={\textbf{C40}: Security Reframing via Alternative Attack Surface \textcolor{gray}{\small (O14/H5/R14, conf: \textbf{high})}}, breakable, before skip=4pt, after skip=4pt]
\textbf{Tactical pattern.} The shared move is: take a security setting where attacks or defenses are conventionally gated by access to a single channel (pixel-space triggers, final output access, white-box parameters, raw training data, input-side filters), then demonstrate that an OVERLOOKED structural property --- intermediate activations (NeurIPS\_2023\_0662), the column-space geometry of logits (ICML\_2024\_1128), label-only co-annotations (ICLR\_2023\_0177), order-permutation likelihood (ICLR\_2024\_0930), parameter-space perturbation direction (ICLR\_2025\_1362), residual probability mass (ICML\_2024\_1038), prediction variance under random masking (ICLR\_2025\_1567), or the internal computational pathway itself (NeurIPS\_2024\_1235) --- is INDEPENDENTLY SUFFICIENT to mount the attack or anchor the defense. The deliverable is a constructive proof that the channel everyone watched was not necessary, plus a working pipeline that operates entirely on the alternative surface. Verbs are "shift", "substitute", "circumvent" the assumed prerequisite; objects are concrete information-bearing artifacts (logit subspace, soft-label distribution, attention scores, spectral energy) rather than reformulated objectives.

\textbf{Differentiation.} \textit{Closest sibling is cluster 39 (Objective Realignment for Robustness via Geometric Targeting): cluster 39 holds the threat model fixed and rewrites the TRAINING OBJECTIVE to better target a worst-case property the original loss missed (e.g., Bellman-infinity error, Rockafellian relaxation), whereas cluster 40 holds training fixed and rewrites WHAT COUNTS AS THE ATTACK SURFACE --- the contribution is a new channel of information leakage or intervention, not a new loss. Distinct also from cluster 12 (Structural Sensitivity Analysis for Mechanism Targeting), which stays within the canonical DP threat model and reallocates noise/budget across composition steps; cluster 40 instead enlarges or relocates the threat model itself (e.g., distillation transfers membership, ICML\_2024\_1128 shows API logits leak parameters). And distinct from cluster 30 (Objective Auditing via Explicit Constraint Injection): cluster 30 enforces a suppressed statistical property in training; cluster 40 reveals a suppressed exposure surface in deployment.}

\textbf{When to pick.} Pick this tactic when (a) the problem is adversarial --- an attack, defense, audit, or leakage analysis --- AND (b) the field's published attacks/defenses converge on a single channel (pixel triggers, final logits, gradient access, intermediate features) treated as a hard prerequisite, AND (c) you can name a concrete structural artifact in the system (a subspace, a statistic, a permutation order, a parameter direction) that carries equivalent information through a different channel.

\textbf{Tactical failure.} Rejects in this cluster fail when the proposed "alternative surface" collapses, on inspection, into an INCREMENTAL VARIATION OF THE CONVENTIONAL SURFACE rather than a categorically different channel. ICLR\_2024\_0790 adds an activation-similarity auxiliary loss to standard backdoor training --- same pixel triggers, same parameter optimization, just with a regularizer; NeurIPS\_2024\_1163 (AttnGCG) layers an attention-score loss onto vanilla GCG, presenting attention as a "new" surface while still optimizing the same suffix tokens against the same output objective; ICML\_2025\_1657 (DPOT\_L0) refines trigger optimization with an L0 budget but operates in the same input-perturbation space defenders already monitor; ICML\_2025\_1745 (relighting) widens the perturbation set to physically plausible lighting but keeps the standard input-space white-box threat model; NeurIPS\_2024\_1170 (CARSO) merely concatenates two existing defense paradigms. Reviewers read these as engineering deltas because the paper never establishes that the conventional channel was bypassable --- there is no demonstration of "X is not needed; only Y", just "X plus a tweak". The diagnostic: if the reformulation can be removed and the method still half-works through the original channel, the threat-model reframe was cosmetic.

\textbf{Examples.}
\begin{itemize}[leftmargin=12pt,topsep=2pt,itemsep=1pt]
  \item \textbf{[Oral]} \texttt{ICLR\_2023\_0177} \textemdash\ Cleanest instance of channel substitution: backdoor field had universally treated pixel-space triggers as the attack vector, and this paper shows that label co-annotations alone --- with clean images --- are a sufficient and harder-to-detect alternative surface.
  \\ \textcolor{gray}{\footnotesize \textit{Lesson:} Backdoor field universally treated pixel triggers as the vector; show that label co-annotations alone (clean images) are sufficient --- channel substitution at the categorical level.}
  \item \textbf{[Oral]} \texttt{ICML\_2024\_1128} \textemdash\ Reframes API logit outputs as a geometric artifact (column space of the unembedding matrix), making exact parameter recovery a linear-algebra exercise on a channel that prior model-extraction work treated as merely behavioral mimicry.
  \\ \textcolor{gray}{\footnotesize \textit{Lesson:} API logit outputs reframed as geometric artifact (column space of unembedding matrix) --- exact parameter recovery becomes a linear-algebra exercise, not behavioral mimicry.}
  \item \textbf{[Oral]} \texttt{ICLR\_2024\_0930} \textemdash\ Shifts the contamination-detection surface from content memorization (inaccessible behind black-box APIs) to ORDER memorization, exploiting permutation likelihood as a structural memorization fingerprint testable from output probabilities alone.
  \\ \textcolor{gray}{\footnotesize \textit{Lesson:} Shift contamination detection from content memorization (inaccessible) to ORDER memorization --- permutation likelihood as structural fingerprint testable from output probabilities.}
  \item \textbf{[Oral]} \texttt{ICLR\_2025\_1567} \textemdash\ Replaces input-space trigger detection with a behavioral surface --- prediction variance under random edge dropping --- and proves it is a universal indicator that holds across backdoor attack constructions, sidestepping the trigger-pattern arms race entirely.
  \\ \textcolor{gray}{\footnotesize \textit{Lesson:} Replace input-space trigger detection with prediction variance under random edge dropping --- universal indicator across backdoor constructions, sidesteps the trigger-pattern arms race.}
  \item \textbf{[Reject]} \texttt{ICLR\_2024\_0790} \textemdash\ Frames feature-activation similarity as a 'new' evasion surface but operationally just adds an auxiliary loss to conventional pixel-trigger backdoor training --- the original input channel is unchanged, so the reframing is cosmetic.
  \\ \textcolor{gray}{\footnotesize \textit{Lesson:} Feature-activation similarity as 'new' surface but operationally just adds an auxiliary loss to conventional pixel-trigger training --- original input channel unchanged, reframing cosmetic.}
  \item \textbf{[Reject]} \texttt{NeurIPS\_2024\_1163} \textemdash\ Claims attention scores as a manipulable attack surface but layers an attention-targeting loss onto standard GCG suffix optimization, leaving the underlying input-space threat model and optimization target intact.
  \\ \textcolor{gray}{\footnotesize \textit{Lesson:} Attention-targeting loss layered onto standard GCG suffix optimization leaves the underlying input-space threat model intact.}
  \item \textbf{[Reject]} \texttt{ICML\_2025\_1657} \textemdash\ Presents L0-bounded triggers as evading update-statistic defenses, but the surface is still adversarial input perturbation in the same federated-learning threat model --- only the perturbation budget is refined, not the channel.
  \\ \textcolor{gray}{\footnotesize \textit{Lesson:} L0-bounded triggers refined the perturbation budget but the surface is still adversarial input perturbation in the same threat model.}
\end{itemize}
\end{tcolorbox}

\subsection{Encode Invariance as Hard Constraint}
\textit{Parent methodology id: \texttt{invariance\_as\_architectural\_constraint}; 6 tactical clusters.}

\begin{tcolorbox}[colback=bgskill!50, colframe=cgreen!70, title={\textbf{C01}: Symmetry Exploitation via Equivariant Architecture Design \textcolor{gray}{\small (O17/H2/R1, conf: \textbf{low})}}, breakable, before skip=4pt, after skip=4pt]
\textbf{Tactical pattern.} The shared move is: identify the EXACT symmetry group acting on the input/parameter space (permutation of neurons, scaling of MLP weights, SE(3) on 3D points, O(n)/SO(n)/Sp(n), full transformer functional symmetry group, Clifford/multivector group), then enforce equivariance through one of three architectural moves: (a) derive a basis of equivariant linear maps via diagram algebra or weight-tying (Brauer algebra, channel-splitting, Clifford multiplication) --- see NeurIPS\_2023\_0605, ICML\_2023\_0424, ICML\_2023\_0450; (b) re-encode the input as a graph/manifold for which an off-the-shelf equivariant architecture already exists (treat NN weights as a graph, lift image features onto SO(3)) --- see ICLR\_2024\_0846, ICLR\_2023\_0250; or (c) canonicalize via alignment or frame averaging so non-equivariant components inherit equivariance --- see ICLR\_2023\_0243, ICLR\_2022\_0096, NeurIPS\_2025\_1878. The output is always a model with provable group-commuting layers, not a model trained to approximate the symmetry from augmented data.

\textbf{Differentiation.} \textit{Versus sibling 23 (Structural Inductive Bias Transfer for Generative Design), this cluster is about constructing the equivariance machinery itself or porting it to new symmetry groups/parameter-space domains, whereas 23 instantiates equivariant components inside generative pipelines whose primary contribution is the design objective (antibody binding, fragment-level molecule synthesis), not the symmetry analysis. Versus sibling 36 (Continuous Dynamics as Inductive Structure), this cluster never embeds the problem in an ODE/SDE --- it works on discrete layer algebra and group representations, not on continuous flows whose conservation laws emerge from dynamics. Versus sibling 46 (Iterative Refinement Process Redesign), the GeoDiff-style paper here (ICLR\_2022\_0097) modifies a diffusion kernel to be equivariant rather than redesigning the noise trajectory for efficiency or manifold-preservation --- the locus of the modification is the equivariance constraint, not the refinement schedule.}

\textbf{When to pick.} Pick this tactic when the problem domain admits a mathematically precise symmetry group (a named Lie group, a permutation group of a structured object, or a parameter-space symmetry like neuron-permutation/weight-scaling) AND when the success criterion includes provable equivariance, sample efficiency, or out-of-distribution generalization across symmetry orbits --- not just empirical robustness from data augmentation.

\textbf{Tactical failure.} With only one reject (NeurIPS\_2023\_0721) the failure mode is necessarily speculative, but the reject's profile is informative: it proposes adding bilinear tensor product layers plus a local SO(2) gauge-equivariant operation on top of the already-mature SO(3)-equivariant family, marketed as "a richer class of equivariant maps." The likely tactical failure is incremental-extension-of-an-established-equivariant-toolbox: the symmetry group is the same one (SO(3)) the community has already saturated with linear-equivariant constructions, the algebraic addition (bilinear with Clebsch--Gordan contraction) is a known move in physics, and the new structure (local SO(2) gauge) is bolted on for one application domain (jet physics) rather than identifying a previously-uncharacterized symmetry of a new domain (as ICML\_2023\_0450 did for weight spaces or NeurIPS\_2024\_1303 did for scaling symmetry). The lesson is that equivariance papers reward NEW symmetries on NEW domains or canonical algebraic foundations, not richer linear-algebraic gadgets over an already well-served group.

\textbf{Examples.}
\begin{itemize}[leftmargin=12pt,topsep=2pt,itemsep=1pt]
  \item \textbf{[Oral]} \texttt{ICML\_2023\_0450} \textemdash\ Canonical example of identifying a previously-unstudied symmetry group (joint row-column neuron permutations of weight matrices) and deriving equivariant layers commuting with it from scratch.
  \\ \textcolor{gray}{\footnotesize \textit{Lesson:} Identify a previously-unstudied symmetry group first (joint row-column neuron permutations), then derive equivariant layers from scratch --- finding the right symmetry IS the contribution.}
  \item \textbf{[Oral]} \texttt{ICLR\_2024\_0846} \textemdash\ Canonical example of the data-representation-swap variant: rather than build new equivariant layers, recast NN weights as a graph so off-the-shelf permutation-equivariant GNNs become applicable for free.
  \\ \textcolor{gray}{\footnotesize \textit{Lesson:} Data-representation swap (NN weights $\to$ graph) makes off-the-shelf permutation-equivariant GNNs applicable for free --- sometimes the cleanest move is not new layers but a new view of the data.}
  \item \textbf{[Oral]} \texttt{NeurIPS\_2023\_0605} \textemdash\ Canonical example of importing an external algebraic structure (Clifford algebra) as the substrate from which equivariance falls out automatically across many groups, rather than per-group representation engineering.
  \\ \textcolor{gray}{\footnotesize \textit{Lesson:} Importing an external algebraic substrate (Clifford algebra) makes equivariance fall out automatically across many groups --- borrow at the right structural level instead of per-group engineering.}
  \item \textbf{[Oral]} \texttt{ICLR\_2022\_0096} \textemdash\ Canonical canonicalization variant: prove that averaging over a small 'frame' subset of group elements yields exact equivariance, replacing intractable group-wide averaging with a provably-sufficient representative set.
  \\ \textcolor{gray}{\footnotesize \textit{Lesson:} Averaging over a small representative 'frame' subset can give exact equivariance --- a tractable substitute for intractable group-wide averaging is more valuable than a new architecture.}
  \item \textbf{[Reject]} \texttt{NeurIPS\_2023\_0721} \textemdash\ Adds bilinear tensor product layers and a local SO(2) gauge to existing SO(3)-equivariant networks --- an incremental algebraic extension on a saturated group rather than identification of a new symmetry domain or a new algebraic foundation.
  \\ \textcolor{gray}{\footnotesize \textit{Lesson:} Adding bilinear tensor products on a saturated group (SO(3)) is incremental --- equivariant work needs either a new symmetry domain OR a new algebraic foundation.}
\end{itemize}
\end{tcolorbox}

\begin{tcolorbox}[colback=bgskill!50, colframe=cgreen!70, title={\textbf{C23}: Structural Inductive Bias Transfer for Generative Design \textcolor{gray}{\small (O8/H1/R13, conf: \textbf{high})}}, breakable, before skip=4pt, after skip=4pt]
\textbf{Tactical pattern.} Bake a domain-derived structural primitive --- E(3)/SE(n) symmetry over atomic coordinates, chemically meaningful fragment vocabularies, physical interaction fields (shape/electrostatics/pharmacophores), or preference orderings over property scores --- directly into a *generative* objective (diffusion, flow matching, conditional language model, or DPO-style finetuning) for molecular or protein design. Concretely the move is: replace the default unconstrained sampling target (atom-by-atom coordinates, raw sequence tokens, scalar reward) with a target whose support is already filtered by the prior --- joint coordinate+property denoising (ICLR\_2025\_1582), fragment-level discrete flow (NeurIPS\_2025\_1872), ellipsoidal layout conditioning (ICLR\_2025\_1546), or pairwise energy-rank loss (NeurIPS\_2024\_1208, ICML\_2024\_1100) --- so that property optimization, validity, and expressivity are co-resolved at training time rather than via post-hoc filtering. The pivotal claim is that an architectural+objective prior tailored to the design domain breaks an apparent expressivity$\leftrightarrow$generalization or coverage$\leftrightarrow$validity tradeoff (ICML\_2024\_1033, ICML\_2024\_1002).

\textbf{Differentiation.} \textit{Where Sibling Cluster 1 (Equivariant Architectures) enforces symmetry on a *processing* task --- consuming existing structured inputs (weight spaces, graphs, point clouds) and producing predictions that commute with the symmetry group --- this cluster lifts that same constraint logic into the *generative* setting for biomolecular design: the symmetry/fragment/property prior shapes the noise process, denoising network, or preference loss so the *sampled* artifact is born domain-valid (e.g., ICLR\_2023\_0182 generates CDR coordinates+sequences via E(3)-equivariant translation, not just classifies them). It also differs from Sibling Cluster 46 (Iterative Refinement Redesign), which redesigns the noise trajectory for general signal-restoration efficiency --- here the redesign is keyed specifically to molecular/protein domain primitives (fragments, pharmacophores, ellipsoidal scaffolds, synthesizability rewards), and the validation metric is a chemistry/biology property, not perceptual quality or sample efficiency.}

\textbf{When to pick.} Pick this tactic when the design target is a discrete-continuous biomolecular object (small molecule, protein, antibody, crystal) whose *validity* and *desired property* are governed by a known domain primitive --- a symmetry group, fragment library, physical field, or rankable score --- that current generative pipelines treat as a post-hoc filter. Choose another sibling instead if the task is recognition/processing of given structures (Sibling 1) or generic restoration where domain doesn't supply natural primitives (Sibling 46).

\textbf{Tactical failure.} Rejected papers in this cluster overwhelmingly suffer from "correct prior, insufficient leverage": the proposed structural bias is principled but is either a straightforward composition with an off-the-shelf generator (NeurIPS\_2025\_1872 fragment-level flow matching, ICLR\_2024\_0849 harmonic prior + flow matching for multi-ligand docking, NeurIPS\_2024\_1208 RLHF$\to$energy-rank port to molecules), narrowly targets a niche subpopulation (ICLR\_2025\_1341 matrix VAE for pharmacogenes, NeurIPS\_2023\_0679 linker prompt-tuning for ESMFold heterodimers), or contradicts the field's accepted prior without a strong empirical case (NeurIPS\_2024\_1324 argues *against* enforced invariance for crystals, ICML\_2025\_1692 produces a complete SE(n) invariant but mostly mathematical). Reject reviewers consistently flag (a) limited generative scope (ICLR\_2023\_0189's fragment views are used only for contrastive pretraining, not generation), (b) the prior reduces to a clever loss/feature with no architectural depth (ICML\_2025\_1732 Tanimoto-as-dense-reward, NeurIPS\_2023\_0642 predictor-as-oracle, ICML\_2025\_1746 PLM scoring functions), or (c) the unlocked capability isn't benchmarked against a domain-standard alternative (NeurIPS\_2023\_0700 bi-level Gibbs vs. directed evolution baselines). Orals avoid this by tying the prior to a *joint* expressivity-generalization or validity-property tradeoff that no prior method resolved.

\textbf{Examples.}
\begin{itemize}[leftmargin=12pt,topsep=2pt,itemsep=1pt]
  \item \textbf{[Oral]} \texttt{ICLR\_2023\_0182} \textemdash\ Casts antibody-antigen complex as a multi-component E(3)-equivariant graph and generates CDR coordinates+residues *jointly* in parallel --- the canonical 'symmetry-as-generative-prior for biomolecular design' move.
  \\ \textcolor{gray}{\footnotesize \textit{Lesson:} Cast antibody-antigen complex as multi-component E(3)-equivariant graph and generate CDR coordinates+residues JOINTLY in parallel --- symmetry-as-generative-prior in canonical form.}
  \item \textbf{[Oral]} \texttt{ICLR\_2025\_1582} \textemdash\ ShEPhERD diffuses shape, electrostatics, and pharmacophores as *joint* generative targets, exemplifying this cluster's signature of folding interaction-property priors into the denoising objective rather than post-filtering.
  \\ \textcolor{gray}{\footnotesize \textit{Lesson:} Diffuse shape, electrostatics, and pharmacophores as JOINT generative targets --- fold interaction-property priors into the denoising objective rather than post-filtering.}
  \item \textbf{[Oral]} \texttt{ICML\_2024\_1033} \textemdash\ Explicitly frames the tactic: chemically motivated fragment biases as architectural inductive bias resolve the expressivity-generalization tradeoff in molecular GNNs, articulating the cluster's core hypothesis.
  \\ \textcolor{gray}{\footnotesize \textit{Lesson:} Chemically motivated fragment biases as architectural inductive bias resolve the expressivity-generalization tradeoff --- articulate the cluster's core hypothesis explicitly.}
  \item \textbf{[Oral]} \texttt{ICML\_2024\_1100} \textemdash\ Imports preference optimization (DPO) into conditional residue flow matching for molecules, instantiating the cluster's 'preference ordering as structural prior' variant for property-targeted generation.
  \\ \textcolor{gray}{\footnotesize \textit{Lesson:} Import preference optimization (DPO) into conditional residue flow matching --- preference ordering as structural prior for property-targeted generation.}
  \item \textbf{[Reject]} \texttt{NeurIPS\_2024\_1208} \textemdash\ Energy-rank alignment ports RLHF preference loss to chemical energy rankings --- a direct transfer with the right prior but insufficient architectural commitment, illustrating the 'clever loss without depth' failure.
  \\ \textcolor{gray}{\footnotesize \textit{Lesson:} Direct transfer of RLHF preference loss to chemical energy rankings is the right prior with insufficient architectural commitment --- clever loss without depth.}
  \item \textbf{[Reject]} \texttt{NeurIPS\_2024\_1324} \textemdash\ Inverts the cluster's logic by *removing* the invariance prior for crystal generation; the principled claim that learned invariances substitute for enforced ones lacked the empirical leverage to overturn the domain's standing assumption.
  \\ \textcolor{gray}{\footnotesize \textit{Lesson:} Removing the invariance prior for crystal generation needs the empirical leverage to overturn the domain's standing assumption --- the principled claim alone isn't enough.}
  \item \textbf{[Reject]} \texttt{ICLR\_2023\_0189} \textemdash\ Fragment-based contrastive views are a domain-correct prior but used only for representation pretraining, not the generative-design payoff this cluster's accepted papers deliver.
  \\ \textcolor{gray}{\footnotesize \textit{Lesson:} Fragment-based contrastive views used only for representation pretraining miss the generative-design payoff this cluster's accepted papers deliver.}
\end{itemize}
\end{tcolorbox}

\begin{tcolorbox}[colback=bgskill!50, colframe=cgreen!70, title={\textbf{C26}: Geometric Structure Exploitation for Representation Learning \textcolor{gray}{\small (O8/H0/R4, conf: \textbf{medium})}}, breakable, before skip=4pt, after skip=4pt]
\textbf{Tactical pattern.} Pick a specific geometric object --- a linear subspace, hyperbolic ball, polytope, anchor-defined similarity space, topology-preserving grid, or smoothness-regularized manifold --- and force the learned representation to live in (or be readable through) that object so that semantic relationships become geometric ones. The verb-object signature is concrete: project token embeddings through existing weight matrices to expose linear concept axes (ICML\_2025\_1688), define each instance by its similarity to fixed anchors (ICLR\_2023\_0323), add an entailment-containment loss in the Poincaré ball (ICLR\_2025\_1379), generalize binary direction-vectors to k-simplex polytopes and nested complexes (ICLR\_2025\_1610), partition features by Laplacian smoothness before embedding (ICML\_2025\_1776), or add a 2D-spatial-smoothness penalty over attention units (ICLR\_2025\_1617). The shared move is structural: the geometry of the representation *space* (not the operations acting on it) is engineered to mirror the relational structure of the semantics (hierarchy, categoriality, multi-manifold, topographic adjacency).

\textbf{Differentiation.} \textit{Versus sibling cluster 1 (Symmetry Exploitation via Equivariant Architecture Design), this cluster does not constrain the *operations* (layers, message passing) to commute with a group action; instead it constrains the *embedding space itself* to have a target shape --- hierarchy lives because the metric is hyperbolic, categoriality lives because the locus is a polytope, gauge-arbitrariness vanishes because coordinates are replaced by anchor-relations. Equivariant networks would still produce arbitrary Euclidean embeddings; this cluster picks a non-trivial geometry for those embeddings to inhabit. Versus sibling 36 (Continuous Dynamics as Inductive Structure), the invariance is static-spatial (where points sit) rather than temporal-dynamical (how trajectories evolve). Versus sibling 23 (Structural Inductive Bias Transfer for Generative Design), the target is interpretable/transferable representation rather than constrained sample generation in a specific scientific domain.}

\textbf{When to pick.} Pick this when the invariance you care about is a property of the *embedding manifold* --- hierarchy, multi-class categoriality, multi-modal alignment, topographic neighborhood --- rather than a symmetry of the input domain or the dynamics of a process. Concrete signal: you can name the target geometric object (Poincaré ball, polytope, anchor-similarity vector, 2D smoothness grid) before you choose the architecture.

\textbf{Tactical failure.} Rejects in this cluster fail when the imposed geometric structure looks bolted-on rather than load-bearing: the geometry is introduced but never shown to be necessary versus a strong vanilla-representation baseline, or its semantic payoff is asserted by analogy rather than by a clean ablation. ICLR\_2024\_0789 disentangles centroid/orientation into separate channels but cannot demonstrate that a standard end-to-end encoder couldn't recover the same precision, leaving the explicit geometric channels looking like a stylistic choice. NeurIPS\_2023\_0722's retinotopy-inspired learned grids substitute for fixed ROIs, but the biological-inductive-bias framing carries the paper farther than the empirical gains justify, and the spatial parameterization reads as one of many viable flexibilizations. NeurIPS\_2024\_1241 shows a linear bridge between LM and collaborative-filtering spaces --- exactly the kind of "linear-accessibility" finding ICML\_2025\_1688 turned into an Oral --- but stops at a probing demonstration without converting the geometric observation into a method that beats CF baselines. ICLR\_2023\_0350 co-trains a topology-preserving SOM grid with contrastive learning but the SOM layer never becomes essential to the headline metric, so the geometric structure looks decorative. The pattern: identifying or imposing a geometry is cheap; reviewers reward papers only when the geometry unlocks something (zero-shot stitching, training-cost reduction, hierarchical generalization, mechanistic interpretation) that no flat-Euclidean alternative achieves.

\textbf{Examples.}
\begin{itemize}[leftmargin=12pt,topsep=2pt,itemsep=1pt]
  \item \textbf{[Oral]} \texttt{ICLR\_2023\_0323} \textemdash\ Archetypal example of replacing absolute-coordinate embeddings with a relational geometry (similarity-to-anchors), turning an arbitrary gauge into a structural invariance that yields zero-shot latent-space stitching.
  \\ \textcolor{gray}{\footnotesize \textit{Lesson:} Replace absolute-coordinate embeddings with similarity-to-anchors --- turn an arbitrary gauge into a structural invariance yielding zero-shot latent-space stitching.}
  \item \textbf{[Oral]} \texttt{ICLR\_2025\_1610} \textemdash\ Pure instance of picking the right geometric object (vector $\to$ polytope $\to$ nested simplicial complex) so that categorical and hierarchical concept structure becomes literally readable off the unembedding space.
  \\ \textcolor{gray}{\footnotesize \textit{Lesson:} Vector $\to$ polytope $\to$ nested simplicial complex --- categorical and hierarchical concept structure becomes literally readable off the unembedding space.}
  \item \textbf{[Oral]} \texttt{ICLR\_2025\_1379} \textemdash\ Cleanly pairs a curved geometry (hyperbolic ball) with an explicit containment loss, demonstrating the cluster's tenet that a target geometry must be enforced by a structural objective rather than assumed to emerge.
  \\ \textcolor{gray}{\footnotesize \textit{Lesson:} Pair curved geometry (hyperbolic ball) with explicit containment loss --- target geometry must be enforced by a structural objective, not assumed to emerge.}
  \item \textbf{[Oral]} \texttt{ICML\_2025\_1688} \textemdash\ Discovers that a generative model's representation space already contains a linearly-projectable geometric structure for concepts, exemplifying the cluster's reverse direction --- identify, rather than impose, the exploitable geometry.
  \\ \textcolor{gray}{\footnotesize \textit{Lesson:} Generative model's representation space already contains a linearly-projectable geometric structure --- sometimes the move is to identify, not impose, the exploitable geometry.}
  \item \textbf{[Reject]} \texttt{ICLR\_2024\_0789} \textemdash\ Imposes disentangled geometric channels (centroid, orientation) without showing that the explicit geometric factorization unlocks performance unreachable by standard end-to-end training.
  \\ \textcolor{gray}{\footnotesize \textit{Lesson:} Disentangled centroid + orientation channels need to unlock performance unreachable by standard end-to-end training --- geometric factorization must be load-bearing.}
  \item \textbf{[Reject]} \texttt{NeurIPS\_2024\_1241} \textemdash\ Demonstrates the geometric observation (LM$\to$CF space is linearly mappable) but stops at a probe rather than converting the linear-geometry insight into a method that beats baselines.
  \\ \textcolor{gray}{\footnotesize \textit{Lesson:} LM$\to$CF linear mappability is a probe observation, not a method that beats baselines --- convert linear-geometry insight into operational capability.}
  \item \textbf{[Reject]} \texttt{ICLR\_2023\_0350} \textemdash\ Joint contrastive+topology-preserving training imposes an SOM grid geometry, but the structured grid never becomes load-bearing for the headline result, leaving the geometric layer decorative.
  \\ \textcolor{gray}{\footnotesize \textit{Lesson:} Joint contrastive + topology-preserving SOM grid never becomes load-bearing for the headline result --- the geometric layer stays decorative.}
\end{itemize}
\end{tcolorbox}

\begin{tcolorbox}[colback=bgskill!50, colframe=cgreen!70, title={\textbf{C27}: 2D-to-3D Semantic Lifting with Structural Priors \textcolor{gray}{\small (O10/H0/R5, conf: \textbf{high})}}, breakable, before skip=4pt, after skip=4pt]
\textbf{Tactical pattern.} The shared move is to project 2D signals (foundation-model features, planar predictions, photometric losses) into a 3D-native representation (sparse voxels, 3D Gaussians, neural plane fields, SDF patches) and bolt on auxiliary geometric supervision (normal maps, SDF loss, multi-view consistency, depth-based projection) so the lifting itself does the heavy lifting instead of the network. Concretely: pick a representation whose primitives carry explicit 3D coordinates, define a deterministic 2D$\to$3D projection rule (depth unprojection, view-frustum splatting, plane-field distillation), and let either a photometric/SDF objective or cross-view consistency replace the missing 3D ground truth. The structural prior is the geometry of the lifting operator (it commutes with view changes / preserves spatial coherence), not a learned symmetry. This converts an under-constrained 2D-supervised reconstruction task into a well-posed one without paying for full 3D supervision.

\textbf{Differentiation.} \textit{Versus sibling cluster 1 (Symmetry Exploitation via Equivariant Architecture Design), which encodes a continuous group symmetry by making layers commute with group actions, this cluster encodes a *coordinate-system* prior: it commits to an explicit 3D scaffold (voxels, Gaussians, neural fields) and uses the projective geometry between 2D and 3D as the fixed structure. There is no group-equivariance proof --- the invariance enforced is multi-view photometric/feature consistency, mediated by depth and camera intrinsics. Versus sibling 26 (Geometric Structure Exploitation for Representation Learning), which imposes geometry on *latent* representation spaces (linear projections, hyperbolic anchors, topology-preserving grids), this cluster imposes geometry on the *physical-world* representation --- the 3D scaffold corresponds to actual scene coordinates, not abstract semantic axes. The differentiator is therefore: structural prior = real Euclidean 3-space + camera projection, not algebraic symmetry and not latent-space geometry.}

\textbf{When to pick.} Pick this instantiation when (a) the target output lives in physical 3D scene coordinates and (b) the supervision you actually have is 2D (foundation-model masks/features, posed/unposed images, planar predictions) --- i.e., a 2D$\to$3D supervision gap exists. The signal that this is the right cluster (rather than equivariant layers or latent-geometry priors) is that you can write down a closed-form projection from your supervision domain to your output representation using camera/depth.

\textbf{Tactical failure.} Reject papers in this cluster fail by either (i) keeping the lifting move shallow --- ICLR\_2025\_1453 (InfiniteMesh) inserts a 2D-inpainting view-interpolation stage *before* a black-box 3D reconstructor, so the structural prior never reaches the 3D representation and reviewers read it as a "preprocessing trick"; (ii) targeting a 2D/loss-side fix that does not actually instantiate a 3D scaffold --- ICLR\_2024\_0795 (Completion Consistency) modifies the supervision objective for one-to-many completions without injecting a stronger 3D inductive bias, leaving the architectural move incremental; (iii) over-pruning the geometric prior --- NeurIPS\_2024\_1332 (Cortical Surfaces) uses ribbon segmentation as the only supervision but the lifting from volumetric mask to surface relies on regularizers rather than a representation that bakes surface structure in, so the well-posedness story is fragile; (iv) substituting the lift with a generic latent-Transformer (ICLR\_2025\_1536 POC-SLT) or scalar-regression shortcut (ICLR\_2023\_0372 Eigen-Lengths) that bypasses explicit 3D structure entirely. The unifying pattern: rejects pay for the *language* of 3D priors but not the *commitment* --- they don't replace the structure-agnostic intermediate (triplane, latent code, scalar) with a coordinate-bearing 3D representation, so the geometric inductive bias is not actually load-bearing.

\textbf{Examples.}
\begin{itemize}[leftmargin=12pt,topsep=2pt,itemsep=1pt]
  \item \textbf{[Oral]} \texttt{NeurIPS\_2024\_1258} \textemdash\ Canonical instance: replaces triplane intermediate with sparse 3D voxels + 3D convs and adds normal-map and SDF auxiliary supervision, so the 3D scaffold and the structural prior are simultaneously committed to.
  \\ \textcolor{gray}{\footnotesize \textit{Lesson:} Sparse 3D voxels + 3D convs + normal-map and SDF auxiliary supervision --- commit to 3D scaffold AND structural prior simultaneously, not just one.}
  \item \textbf{[Oral]} \texttt{ICLR\_2025\_1416} \textemdash\ Exemplifies the 2D-foundation-model$\to$3D lift: takes per-view SAM outputs and projects them into 3D via a geometry-aware aggregation module, inheriting 2D pretraining while making the 3D coordinate system the carrier.
  \\ \textcolor{gray}{\footnotesize \textit{Lesson:} Project per-view SAM outputs into 3D via geometry-aware aggregation --- inherit 2D pretraining while making the 3D coordinate system the carrier.}
  \item \textbf{[Oral]} \texttt{ICLR\_2025\_1520} \textemdash\ Uses a neural plane field as a globally consistent 3D scaffold that distills inconsistent multi-view 2D plane predictions, letting the projective consistency of the field resolve the supervision conflicts.
  \\ \textcolor{gray}{\footnotesize \textit{Lesson:} Neural plane field as globally consistent 3D scaffold distilling inconsistent multi-view 2D plane predictions --- projective consistency resolves supervision conflicts.}
  \item \textbf{[Oral]} \texttt{NeurIPS\_2025\_1863} \textemdash\ Lifts CLIP features into a hierarchical voxel/instance/zone 3D structure via depth projection, using the 3D representation as the substrate for multi-scale spatial reasoning rather than as a rendering target.
  \\ \textcolor{gray}{\footnotesize \textit{Lesson:} Lift CLIP features into hierarchical voxel/instance/zone 3D --- 3D representation becomes the substrate for multi-scale reasoning, not just a rendering target.}
  \item \textbf{[Reject]} \texttt{ICLR\_2025\_1453} \textemdash\ The lift is bolted on as a 2D inpainting interpolation between sparse diffusion views and fed to a generic reconstructor, so the 3D structural prior is never actually injected into a coordinate-bearing representation.
  \\ \textcolor{gray}{\footnotesize \textit{Lesson:} 2D inpainting interpolation between sparse views fed to a generic reconstructor never injects 3D structural prior into a coordinate-bearing representation.}
  \item \textbf{[Reject]} \texttt{ICLR\_2024\_0795} \textemdash\ Modifies the completion objective to handle one-to-many supervision but does not introduce any 3D-native scaffold or geometric prior --- the move is loss-side, not lifting-side.
  \\ \textcolor{gray}{\footnotesize \textit{Lesson:} Modifying the completion objective for one-to-many supervision is loss-side, not lifting-side --- no 3D-native scaffold introduced.}
  \item \textbf{[Reject]} \texttt{NeurIPS\_2024\_1332} \textemdash\ Replaces surface ground truth with volumetric ribbon-segmentation supervision but relies on soft regularizers rather than a representation that intrinsically encodes surface structure, leaving the well-posedness argument under-supported.
  \\ \textcolor{gray}{\footnotesize \textit{Lesson:} Volumetric ribbon-segmentation supervision with soft regularizers leaves the well-posedness argument under-supported --- needs a representation that intrinsically encodes surface structure.}
\end{itemize}
\end{tcolorbox}

\begin{tcolorbox}[colback=bgskill!50, colframe=cgreen!70, title={\textbf{C36}: Continuous Dynamics as Inductive Structure \textcolor{gray}{\small (O10/H0/R6, conf: \textbf{medium})}}, breakable, before skip=4pt, after skip=4pt]
\textbf{Tactical pattern.} Take a problem not originally framed as continuous-time evolution and re-cast it as a trajectory in an ODE/SDE/Lagrangian/Hamiltonian/optimal-transport system, so that the mathematical machinery attached to that framework --- adjoint sensitivity, time-reversal symmetry, saddle-point dualities, conservation laws, variational principles, Feynman-Kac formulas --- becomes a usable computational resource. The verbs are concrete: parameterize the drift of an SDE with a neural net (ICLR\_2024\_0877, ICLR\_2025\_1482), embed known PDE constraints inside the ODE field (ICLR\_2024\_0792), derive saddle-point/adjoint equations for transport (NeurIPS\_2023\_0628, NeurIPS\_2025\_1835), or rewrite parameter updates as Euler-Lagrange equations (ICML\_2023\_0448). The payoff is a capability impossible in the discrete formulation --- exact likelihoods, gradient-free backprop, single-pass dual optimization, hard conservation guarantees, or extrapolation between sparse snapshots.

\textbf{Differentiation.} \textit{The hard contrast is with sibling cluster 46 (Iterative Refinement Process Redesign), which stays inside the diffusion/iterative-refinement paradigm and tweaks the noise schedule, trajectory endpoints, or guidance interval of an already-continuous process. Cluster 36 instead targets problems that were NOT originally continuous (parameter learning, climate state, distribution-shift forecasting, snapshot-based trajectory inference, interaction discovery) and lifts them into a continuous-time dynamical system specifically to inherit its analytic structure. Versus sibling 1 (equivariant architectures), the invariance being exploited is dynamical (time-reversal, conservation, variational extremality) rather than group-theoretic; versus siblings 23, 26, 27, the structural prior is a differential equation governing evolution, not a static geometric/topological constraint on representations.}

\textbf{When to pick.} Pick this cluster when the problem has a natural notion of evolution or transport between two states (snapshots over time, source-to-target distributions, layer-by-layer features, iterative refinement steps) AND there exists a body of continuous-time mathematics (SDE theory, optimal control, Hamiltonian/Lagrangian mechanics, OT) whose closed-form identities (adjoints, reverse-time SDEs, saddle-point duals, conservation laws) you can directly import. If the invariance is group-symmetric on a static input rather than dynamical, prefer sibling 1 instead.

\textbf{Tactical failure.} The dominant reject failure mode is using the continuous reformulation as a descriptive wrapper rather than as a lever that unlocks new computation. NeurIPS\_2024\_1228 maps ResNet feature learning to Hamiltonian geodesics but the analogy explains an observed bottleneck rather than enabling a new training procedure --- the dynamical lens is interpretive, not operational. NeurIPS\_2024\_1176 (CLODE) and ICLR\_2025\_1534 (CO2 autoencoder) drop a Neural ODE / SDE term into a domain-specific pipeline (low-light enhancement, peatland gap-filling) where the continuous framing does not deliver a structural advantage absent in the discrete baseline --- reviewers read these as routine applications of off-the-shelf continuous-time blocks. ICLR\_2024\_0848 stacks a mixture-of-experts on top of a Graph ODE without showing the ODE structure is what enables the heterogeneous-interaction gain. ICLR\_2025\_1519 and NeurIPS\_2023\_0686 propose unifying/alternative formulations (superposition of basis functions, Monte Carlo PDE) but stay close to physics-informed-network territory without demonstrating that the new framing unlocks a capability the existing PINN literature could not reach. The common thread: the continuous embedding is present but does not earn its keep through a derivation, identity, or guarantee that the discrete approach demonstrably could not produce.

\textbf{Examples.}
\begin{itemize}[leftmargin=12pt,topsep=2pt,itemsep=1pt]
  \item \textbf{[Oral]} \texttt{ICLR\_2021\_0063} \textemdash\ Canonical instance of the tactic: recognizing that score-matching and denoising-diffusion are both discretizations of one Itô SDE unlocks reverse-time SDE, exact ODE likelihoods, and predictor-corrector --- capabilities unreachable from either discrete variant.
  \\ \textcolor{gray}{\footnotesize \textit{Lesson:} Recognize score-matching and denoising-diffusion as discretizations of one Itô SDE --- unlocks reverse-time SDE, exact ODE likelihoods, and predictor-corrector unreachable from either discrete variant.}
  \item \textbf{[Oral]} \texttt{NeurIPS\_2025\_1903} \textemdash\ Embeds training itself in a conservative dynamical system so that time-reversal symmetry replaces backpropagation with a forward re-simulation --- the continuous framework's symmetry is converted into a concrete computational resource.
  \\ \textcolor{gray}{\footnotesize \textit{Lesson:} Embed training in a conservative dynamical system --- time-reversal symmetry replaces backpropagation with forward re-simulation, continuous symmetry as concrete computational resource.}
  \item \textbf{[Oral]} \texttt{NeurIPS\_2023\_0628} \textemdash\ Connects static entropic OT to the Schrödinger Bridge dynamic saddle-point, importing the bridge's continuous-time duality to collapse iterative EOT into single-pass end-to-end training.
  \\ \textcolor{gray}{\footnotesize \textit{Lesson:} Connect static entropic OT to Schrödinger Bridge dynamic saddle-point --- import bridge's continuous-time duality to collapse iterative EOT into single-pass training.}
  \item \textbf{[Oral]} \texttt{ICLR\_2024\_0792} \textemdash\ Hard-codes the continuity equation into the drift of a Neural ODE so that conservation is satisfied by construction --- the differential-equation embedding directly delivers a guarantee a black-box forecaster cannot.
  \\ \textcolor{gray}{\footnotesize \textit{Lesson:} Hard-code the continuity equation into the drift of a Neural ODE --- conservation satisfied by construction, a guarantee a black-box forecaster cannot deliver.}
  \item \textbf{[Reject]} \texttt{NeurIPS\_2024\_1228} \textemdash\ Hamiltonian-geodesic mapping is treated as an explanatory lens for ResNet bottleneck structure rather than producing a new training algorithm or guarantee, so the continuous embedding remains descriptive.
  \\ \textcolor{gray}{\footnotesize \textit{Lesson:} Hamiltonian-geodesic mapping treated as explanatory lens, not producing a new training algorithm or guarantee --- descriptive embedding without payoff.}
  \item \textbf{[Reject]} \texttt{NeurIPS\_2024\_1176} \textemdash\ Wraps low-light curve adjustment in a Neural ODE for adaptive step count, but the continuous framing yields no structural property (conservation, adjoint identity, transport duality) beyond hyperparameter convenience.
  \\ \textcolor{gray}{\footnotesize \textit{Lesson:} Neural ODE for adaptive step count yields hyperparameter convenience but no structural property (conservation, adjoint identity, transport duality).}
  \item \textbf{[Reject]} \texttt{ICLR\_2025\_1534} \textemdash\ Adds an SDE-based loss term to an autoencoder for CO2 gap-filling --- the continuous-dynamics ingredient is a soft regularizer on a narrow application rather than a framework whose machinery is exploited.
  \\ \textcolor{gray}{\footnotesize \textit{Lesson:} SDE-based loss term as soft regularizer is a narrow application, not a framework whose machinery is exploited.}
\end{itemize}
\end{tcolorbox}

\begin{tcolorbox}[colback=bgskill!50, colframe=cgreen!70, title={\textbf{C46}: Iterative Refinement Process Redesign \textcolor{gray}{\small (O9/H2/R7, conf: \textbf{medium})}}, breakable, before skip=4pt, after skip=4pt]
\textbf{Tactical pattern.} Take the diffusion / iterative-refinement framework as a given primitive and surgically swap out one of its three canonical components --- (a) the corruption process (replace Gaussian noise with domain-structured corruption: LR-HR residual shifts in ResShift, random transpositions on the symmetric group in SymmetricDiffusers, graph-coupled transition matrices in Graph Diffusion Transformers, heat-kernel diffusion on directed Laplacians), (b) the sampling trajectory (Picard iteration to parallelize sequential denoising in Parallel Sampling, time-restricted classifier-free guidance, midpoint-evaluated DPS approximations), or (c) the prediction target / objective (define the diffusion process directly over the discriminative output space such as correspondence maps in Dense Matching, add a co-prediction objective like flow in VideoJAM). The unifying verb-object is "rederive the noise/trajectory/objective so it commutes with a domain constraint or efficiency budget" --- never "design a new architecture."

\textbf{Differentiation.} \textit{Closest sibling is cluster 36 (Continuous Dynamics as Inductive Structure): cluster 36 papers REACH FOR continuous-time dynamical formulations (ODEs/SDEs/OT flows) as the inductive primitive --- Score-Based SDE (ICLR\_2021\_0063) literally INTRODUCES the continuous diffusion framework. Cluster 46 starts AFTER that framework already exists and modifies its internals (noise distribution, sampling order, guidance schedule, prediction head). Concretely: cluster 36 papers prove "this problem IS a continuous dynamical system"; cluster 46 papers prove "the canonical diffusion process is the wrong instantiation for this domain --- here is the structurally correct corruption/trajectory." Versus cluster 23 (Structural Inductive Bias for Generative Design), cluster 23 papers redesign generator architectures and conditioning contexts at large for molecular/biological generation, while cluster 46 leaves the architecture intact and only edits the diffusion process itself (e.g., FragFM in cluster 23 redesigns the unit of generation to fragments, whereas Graph Diffusion Transformers in cluster 46 keeps node-level generation but redesigns the transition matrix).}

\textbf{When to pick.} Pick cluster 46 when you already plan to use diffusion / iterative refinement and have identified a SPECIFIC mismatch between standard Gaussian-IID noise (or uniform-step sampling, or generic L2 prediction) and your domain --- e.g., the source distribution is already close to the target (super-resolution), the data lives on a non-Euclidean structure (groups, graphs, directed Laplacians), or sequential sampling cost is the bottleneck. Pick cluster 36 instead when you must FIRST argue that a continuous dynamical formulation is appropriate at all.

\textbf{Tactical failure.} Reject papers in this cluster fall into three recurring traps. First, "drop-in component swap without re-derivation": ICLR\_2025\_1382 substitutes a consistency-model x0 prediction for the DDPM posterior mean inside an existing solver and ICLR\_2025\_1373 (ClearSR) injects VAE-encoder LR embeddings into a frozen latent diffusion model --- both swap a single block without proving the substitution is structurally necessary or analyzing how it changes the underlying SDE/likelihood, so reviewers read them as engineering tweaks rather than process redesigns. Second, "redesign attached to a niche or post-hoc use case": ICLR\_2025\_1651 (WMAdapter) bolts an adapter onto a pretrained diffusion model for watermarking, and NeurIPS\_2025\_1871 (NCDPO) reframes diffusion policies as noise-conditioned deterministic for RL fine-tuning --- the redesign is real, but the stakes (watermarking, one RL setting) feel local and the rederivation is not connected to a broader principle. Third, "principled noise process but missing the consequence": ICLR\_2024\_0811 derives closed-form heat-kernel corruptions for directed-graph Laplacians and ICML\_2025\_1822 introduces "generation-rate" manifold geometry, but neither demonstrates that the new corruption/lens unlocks a capability the canonical diffusion process cannot match --- the math is correct, the payoff is not legible. ICLR\_2024\_0991 (Diffusion Inversion) similarly optimizes initial noise rather than latents, but the gain over existing latent-edit baselines is presented as preservation rather than new capability. The common deficit: orals like ResShift and SymmetricDiffusers re-derive the full Markov chain or denoising objective from the new corruption AND show that no canonical Gaussian variant can cover the use case; rejects show only the swap.

\textbf{Examples.}
\begin{itemize}[leftmargin=12pt,topsep=2pt,itemsep=1pt]
  \item \textbf{[Oral]} \texttt{NeurIPS\_2023\_0719} \textemdash\ Canonical exemplar of corruption-process redesign: replaces the universal Gaussian-noise prior with an LR-HR residual-shift Markov chain that exploits the source distribution already being close to the target, collapsing hundreds of denoising steps to a handful.
  \\ \textcolor{gray}{\footnotesize \textit{Lesson:} Replace universal Gaussian-noise prior with LR-HR residual-shift Markov chain exploiting source $\approx$ target --- collapse hundreds of denoising steps to a handful by redesigning corruption.}
  \item \textbf{[Oral]} \texttt{ICLR\_2025\_1600} \textemdash\ Defines diffusion noise as random transpositions on the symmetric group and rederives the denoising objective via group Fourier theory --- the textbook case of redesigning the corruption to commute with the domain's algebraic structure rather than embedding into Euclidean space.
  \\ \textcolor{gray}{\footnotesize \textit{Lesson:} Define diffusion noise as random transpositions on the symmetric group; rederive the denoising objective via group Fourier theory --- corruption commutes with domain's algebraic structure.}
  \item \textbf{[Oral]} \texttt{NeurIPS\_2024\_1225} \textemdash\ Replaces independent node/edge noise with a graph-coupled transition matrix, mathematically reworking the forward process so reverse transitions remain tractable while respecting molecular relational structure --- pure noise-process redesign with architecture left intact.
  \\ \textcolor{gray}{\footnotesize \textit{Lesson:} Replace independent node/edge noise with a graph-coupled transition matrix --- rework the forward process so reverse transitions remain tractable while respecting molecular structure.}
  \item \textbf{[Oral]} \texttt{ICML\_2023\_0504} \textemdash\ Trajectory-side instance of the tactic: keeps the diffusion model unchanged but replaces the canonical sequential sampler with Picard fixed-point iteration that parallelizes denoising steps, trading depth for compute.
  \\ \textcolor{gray}{\footnotesize \textit{Lesson:} Keep the diffusion model unchanged but replace sequential sampler with Picard fixed-point iteration --- parallelize denoising steps, trade depth for compute.}
  \item \textbf{[Reject]} \texttt{ICLR\_2025\_1382} \textemdash\ Drop-in swap of DDPM posterior-mean approximation for a consistency-model prediction inside an inverse solver, without rederiving why this is principled for nonlinear measurement operators --- the canonical 'redesign-as-substitution' failure.
  \\ \textcolor{gray}{\footnotesize \textit{Lesson:} Drop-in DDPM $\to$ consistency-model swap inside an inverse solver without rederiving why this is principled for nonlinear measurement operators is the canonical 'redesign-as-substitution' failure.}
  \item \textbf{[Reject]} \texttt{ICLR\_2024\_0811} \textemdash\ Derives a structurally correct directed-Laplacian heat-kernel corruption process but fails to show the new noise process unlocks a generative capability the canonical Gaussian diffusion cannot reach --- math without legible payoff.
  \\ \textcolor{gray}{\footnotesize \textit{Lesson:} Structurally correct directed-Laplacian heat-kernel corruption that fails to show new generative capability is math without legible payoff.}
  \item \textbf{[Reject]} \texttt{ICLR\_2025\_1651} \textemdash\ Conditions the denoising process on a watermark signal via an adapter rather than redesigning the diffusion process itself; reviewers read it as a post-hoc plug-in to a frozen base model rather than a principled process redesign.
  \\ \textcolor{gray}{\footnotesize \textit{Lesson:} Conditioning denoising on a watermark via an adapter is a post-hoc plug-in to a frozen base model --- the cluster requires redesigning the diffusion process itself, not adding a side channel.}
\end{itemize}
\end{tcolorbox}

\subsection{Reframe via Alternative Formalism}
\textit{Parent methodology id: \texttt{formal\_reframing\_via\_equivalence}; 9 tactical clusters.}

\begin{tcolorbox}[colback=bgskill!50, colframe=cgreen!70, title={\textbf{C06}: Game-Theoretic Equilibrium Reformulation \textcolor{gray}{\small (O19/H0/R12, conf: \textbf{high})}}, breakable, before skip=4pt, after skip=4pt]
\textbf{Tactical pattern.} Establish a precise correspondence between a target problem (preference alignment, calibration, eigendecomposition, recommendation, multicalibration, variational inequalities) and a game-theoretic solution concept (Nash equilibrium, correlated equilibrium, coarse correlated equilibrium, swap regret minimum), then import an existing no-regret algorithm (FTRL, OMD, regret matching, swap-regret minimization, no-regret dynamics) so that its convergence/sample-complexity guarantees transfer onto the target problem unchanged. The verbs are concrete: define utility functions whose Nash equilibrium recovers the optimum (EigenGame, ICLR\_2021\_0018), prove an equivalence between a target deviation structure and swap regret (ICML\_2023\_0491, NeurIPS\_2025\_1884), restrict players to a tractable behavioral class (CVaR + quantal response in ICLR\_2025\_1628), or relax a pointwise constraint to its distributional analogue mirroring correlated equilibrium (ICML\_2025\_1706). The objects are shared across the cluster: utility/regret functions, equilibrium fixed points, and no-regret update rules --- distinguishing this cluster from siblings that import other formalisms.

\textbf{Differentiation.} \textit{Whereas sibling cluster 28 (Principled Equivalence Recovery via Objective Reformulation) reveals that an *already-existing* heuristic is equivalent to a classical objective via duality/representer theorems --- i.e., the equivalence is retrospective and the algorithm is unchanged --- cluster 6 actively *casts a non-game problem AS a game* and then runs a no-regret learner that did not previously apply (e.g., EigenGame turning PCA into a competitive game with new parallel updates, or ICLR\_2025\_1460 turning reward-maximization into a two-player zero-sum equilibrium search). Versus sibling 22 (Generative Reframing of Sequential Decision Problems), which imports density-estimation machinery to replace optimization, cluster 6 imports multi-agent equilibrium and regret-minimization machinery --- the deliverable is a fixed point or a regret bound, not a sampling distribution. Versus sibling 43 (Combinatorial Search via Sequential Decision Decomposition), which decomposes a single search problem into one-player MDP-style steps, cluster 6 explicitly introduces *multiple* interacting players or self-referential conditioning events (e.g., ICML\_2025\_1720 treats conditioning events as players) so that the right primitive is swap/external regret across players rather than reward in a single-agent process.}

\textbf{When to pick.} Pick this tactic when the target problem has a self-referential, multi-party, or fixed-point structure whose "solution" is naturally a Nash/correlated/coarse-correlated equilibrium --- e.g., predictions must be consistent with distributions conditioned on themselves, components are defined by mutual orthogonality, or multiple agents optimize against a shared system --- AND when a known no-regret primitive (regret matching, swap regret, FTRL with information-set regularizers) can be mapped onto the problem's deviation/coupling structure. If your problem instead has a single objective whose optimum is already characterized but expensive (use sibling 16 or 22), or whose optimum needs new expressivity bounds (sibling 8), this is the wrong choice.

\textbf{Tactical failure.} Reject papers in this cluster fail in three sharp ways. First, they bolt a learned/heuristic component onto a provable no-regret algorithm in a way that erodes the guarantee that justified the game-theoretic framing in the first place: ICML\_2025\_1675 parameterizes regret-minimization updates with a meta-learned model to bridge the CCE-Nash gap, and ICML\_2025\_1754 replaces local regret decomposition with a learned global aggregator --- both lose the clean fixed-point characterization without a matching new theorem. Second, they extend a known regret/equilibrium bound to a new setting via additional technical machinery without isolating a new structural primitive: ICLR\_2023\_0173 adds a uniqueness condition for many-to-one matching bandits, ICLR\_2024\_0975 builds composite features for linear function approximation in extensive-form games, and NeurIPS\_2023\_0591 fixes biased gradient estimation for NE --- each is a competent extension but does not reveal a new equilibrium structure that reviewers can point to. Third, they introduce a new game variant or abstraction (ICLR\_2024\_0988 Welfare Diplomacy, ICLR\_2023\_0236 formalized feints, NeurIPS\_2024\_1240 K-recall abstractions, NeurIPS\_2025\_1883 topological satisficing paths) without delivering the matching algorithm-with-guarantee that the cluster's accepted papers always provide. ICML\_2025\_1671 (an O(n\textasciicircum{}2) implementation paper) sits at the engineering end of the same failure mode: machinery without a new conceptual primitive.

\textbf{Examples.}
\begin{itemize}[leftmargin=12pt,topsep=2pt,itemsep=1pt]
  \item \textbf{[Oral]} \texttt{ICLR\_2021\_0018} \textemdash\ Canonical instance of casting a non-game problem (PCA) AS a multi-player game by hand-designing per-eigenvector utilities whose Nash equilibrium exactly recovers the optimization solution --- the cluster's prototype move.
  \\ \textcolor{gray}{\footnotesize \textit{Lesson:} Cast a non-game problem (PCA) AS a game by hand-designing per-eigenvector utilities whose Nash equilibrium recovers the optimum --- the cluster's prototype move.}
  \item \textbf{[Oral]} \texttt{ICLR\_2025\_1460} \textemdash\ Imports no-regret learning into LLM alignment by reframing reward-maximization as a two-player zero-sum game over preferences, exemplifying the cluster's tactic of inheriting equilibrium guarantees in a domain that previously used scalar-reward gradient methods.
  \\ \textcolor{gray}{\footnotesize \textit{Lesson:} Reframe reward maximization as two-player zero-sum over preferences and inherit no-regret machinery --- equilibrium guarantees come for free in domains using scalar-reward gradients.}
  \item \textbf{[Oral]} \texttt{NeurIPS\_2025\_1884} \textemdash\ Pure form of the cluster's tactic: prove that multi-dimensional calibration error IS swap regret in online linear optimization, then transfer existing OLO swap-regret algorithms unchanged --- no new algorithmic machinery, just the equivalence.
  \\ \textcolor{gray}{\footnotesize \textit{Lesson:} Prove the equivalence (multi-dim calibration error IS swap regret in OLO) and transfer existing algorithms unchanged --- no new machinery, just the equivalence as the contribution.}
  \item \textbf{[Oral]} \texttt{ICML\_2025\_1706} \textemdash\ Mirrors the correlated-equilibrium relaxation move at a more general level: relaxes a pointwise variational-inequality constraint to its distributional analogue, importing CE algorithms to obtain polynomial-time tractability for nonmonotone operators.
  \\ \textcolor{gray}{\footnotesize \textit{Lesson:} Relax pointwise variational-inequality constraint to its distributional analogue --- mirror correlated-equilibrium relaxations at higher generality to inherit polynomial-time tractability.}
  \item \textbf{[Reject]} \texttt{ICML\_2025\_1675} \textemdash\ Bolts meta-learning onto regret-minimization updates to bridge the CCE-Nash gap, eroding the provable convergence that motivated the game-theoretic framing without supplying a matching new equilibrium-structure theorem.
  \\ \textcolor{gray}{\footnotesize \textit{Lesson:} Bolting meta-learning onto regret-minimization erodes the provable convergence --- supplying a learned heuristic without a matching equilibrium-structure theorem fails the cluster.}
  \item \textbf{[Reject]} \texttt{ICML\_2025\_1754} \textemdash\ Replaces local regret decomposition with a learned global aggregator in self-play --- the same failure shape as 1675: trades a clean cluster primitive (local regret bounds) for a learned heuristic without proving new global guarantees.
  \\ \textcolor{gray}{\footnotesize \textit{Lesson:} Replacing local regret decomposition with a learned global aggregator trades a clean primitive for an unconvinced heuristic --- same failure shape as 1675.}
  \item \textbf{[Reject]} \texttt{ICLR\_2023\_0173} \textemdash\ Extends bandit-learning regret bounds from one-to-one to many-to-one matching markets via an added uniqueness condition, but the contribution reads as a competent setting-extension rather than the new equilibrium-structure primitive that accepted cluster papers exhibit.
  \\ \textcolor{gray}{\footnotesize \textit{Lesson:} Setting-extension of bandit-learning regret bounds to many-to-one matching is a competent extension but lacks the new equilibrium-structure primitive accepted papers exhibit.}
\end{itemize}
\end{tcolorbox}

\begin{tcolorbox}[colback=bgskill!50, colframe=cgreen!70, title={\textbf{C08}: Formal Abstraction for Tight Expressivity Bounds \textcolor{gray}{\small (O14/H1/R1, conf: \textbf{low})}}, breakable, before skip=4pt, after skip=4pt]
\textbf{Tactical pattern.} Map a neural architecture's forward computation onto a pre-existing formal-mathematical object whose distinguishing power is already fully characterized --- transformation semigroups (ICLR\_2023\_0380), homomorphism-count vectors over relational structures (ICLR\_2024\_0779, NeurIPS\_2024\_1333, ICLR\_2025\_1446), circuit-complexity classes (NeurIPS\_2023\_0749), tensor first-order logic with counting (ICLR\_2022\_0091, ICLR\_2025\_1623), polynomial bases over equivariant tensors (ICML\_2023\_0451), or polyhedral lattices (NeurIPS\_2025\_1856) --- then import that formalism's native separation theorems to derive tight upper bounds, explicit indistinguishability counterexamples, and strict hierarchy results. The signature verbs are "encode-as," "count-exactly," "prove-equivalent-to," "construct-counterexample-via," with the architecture serving as the system whose expressivity is *measured by* the formal object, not merely *inspired by* it. Outputs are typically of the form "architecture A computes exactly the homomorphism counts of pattern class P" or "architecture A requires depth Ω(f(n)) to simulate B" --- quantitative or constructive, not qualitative.

\textbf{Differentiation.} \textit{Whereas sibling 14 (Structural Reduction to Provably Benign Loss Landscapes) reframes architecture analysis to study *optimization dynamics* (whether SGD reaches a global min), this cluster reframes architecture analysis to study *input-output expressivity* (which functions/separations are reachable in principle). And whereas sibling 28 (Principled Equivalence Recovery via Objective Reformulation) maps a *heuristic objective* to a classical principled framework to justify it post-hoc, cluster 8 maps an *architectural primitive* to a formal language whose separation power is already classified, in order to derive impossibility/tightness results rather than vindicate a method. Sibling 16 (Mechanism Reframing via Mathematical Equivalence) is the closest cousin but targets cheaper *implementations* of an existing operation; cluster 8 is uninterested in compute cost and instead extracts capability ceilings via an external formalism's catalog of separations.}

\textbf{When to pick.} Pick this cluster's tactic when the open question is "what can this architecture provably compute / not compute?" and there exists a mature formal catalog (graph homomorphism theory, semigroup theory, circuit complexity, FO logic with counting, polynomial invariant theory) whose objects naturally align with the architecture's update rule. Pick a sibling instead if the question is about training convergence, runtime cost, or post-hoc justification of an empirical objective.

\textbf{Tactical failure.} With only one Reject in this cluster, the failure mode is necessarily speculative. The single rejected paper, ICML\_2025\_1808, attempts the same tactic --- encoding a transformer's fixed-precision arithmetic as a structural constraint and proving the model collapses to recognizing only finitely or co-finitely many strings --- but appears to invoke a strong assumption on query-key parameter structure to obtain its tight bound, narrowing the result to a contrived precision regime rather than a natural language-class boundary. The likely failure pattern is *bound brittleness*: the formal-abstraction tactic collapses to a reject when the formalism's match to the architecture requires non-standard assumptions (artificial precision models, tailored parameter constraints, ad-hoc graph features) so the resulting "tight" bound characterizes the assumption rather than the architecture, and reviewers cannot import the result into the broader expressivity literature. Confidence is low --- one reject does not establish a pattern.

\textbf{Examples.}
\begin{itemize}[leftmargin=12pt,topsep=2pt,itemsep=1pt]
  \item \textbf{[Oral]} \texttt{ICLR\_2023\_0380} \textemdash\ Translates transformer depth analysis into transformation-semigroup decomposition and imports the classical solvable-vs-non-solvable group dichotomy to predict O(log T) vs O(1) depth tightly --- a textbook instance of borrowing a fully classified algebraic catalog to obtain expressivity bounds.
  \\ \textcolor{gray}{\footnotesize \textit{Lesson:} Translate transformer depth into transformation-semigroup decomposition; import the classical solvable-vs-non-solvable dichotomy to predict O(log T) vs O(1) tightly.}
  \item \textbf{[Oral]} \texttt{ICLR\_2024\_0779} \textemdash\ Replaces the binary WL hierarchy with the homomorphism-count vector --- a quantitative algebraic object from finite model theory --- to produce a strictly finer expressivity ordering for GNN classes.
  \\ \textcolor{gray}{\footnotesize \textit{Lesson:} Replace binary WL hierarchy with homomorphism-count vector --- a quantitative algebraic object from finite model theory produces a strictly finer expressivity ordering.}
  \item \textbf{[Oral]} \texttt{NeurIPS\_2023\_0749} \textemdash\ Imports circuit-complexity lower bounds (TC0/AC0-style) directly to prove transformer impossibility for certain tasks and then constructively shows chain-of-thought converts sequential generation into effective depth --- a clean architecture-to-formal-class encoding.
  \\ \textcolor{gray}{\footnotesize \textit{Lesson:} Import circuit-complexity lower bounds (TC0/AC0) directly to prove transformer impossibility, then constructively show CoT converts sequential generation into effective depth.}
  \item \textbf{[Oral]} \texttt{ICLR\_2022\_0091} \textemdash\ Encodes GNN computations into a tensor first-order language so that the WL hierarchy collapses to two syntactic counts (index number, summation depth), turning expressivity analysis into a purely syntactic check.
  \\ \textcolor{gray}{\footnotesize \textit{Lesson:} Encode GNN computations into a tensor first-order language --- WL hierarchy collapses to two syntactic counts (index number, summation depth), expressivity becomes a syntactic check.}
  \item \textbf{[Reject]} \texttt{ICML\_2025\_1808} \textemdash\ Pursues the same tactic --- encode fixed-precision transformers in a formal language-theoretic frame to bound expressivity --- but obtains its tight collapse to finite/co-finite languages only by imposing a structural assumption on query-key parameters, so the bound characterizes the assumption rather than transformers as practitioners use them.
  \\ \textcolor{gray}{\footnotesize \textit{Lesson:} Imposing a structural assumption on query-key parameters to obtain tight bounds characterizes the assumption, not transformers as practitioners use them.}
\end{itemize}
\end{tcolorbox}

\begin{tcolorbox}[colback=bgskill!50, colframe=cgreen!70, title={\textbf{C14}: Structural Reduction to Provably Benign Loss Landscapes \textcolor{gray}{\small (O16/H1/R8, conf: \textbf{high})}}, breakable, before skip=4pt, after skip=4pt]
\textbf{Tactical pattern.} Identify a concrete structural lever --- NTK linearization (ICML\_2024\_1059, ICLR\_2025\_1649), low-rank/invariant subspace (ICML\_2024\_1014, ICML\_2025\_1726), exact convex lift (ICLR\_2022\_0141, ICLR\_2025\_1425), proportional asymptotic regime (ICML\_2023\_0458), continuous-time mean-field limit (ICLR\_2021\_0027, ICLR\_2025\_1441), smooth surjective parameterization of a constraint set (ICML\_2023\_0444), or dimensional lifting (ICML\_2023\_0500) --- and PROVE that under this lever every critical point of the loss is either global, escapable (strict saddle), or aligned with a known optimal structure. The verbs are "characterize all critical points", "prove no spurious local minima", "establish global convergence of gradient flow", "convert spurious minima into strict saddles". The object is always the loss landscape geometry of a specific architecture (LoRA, two-layer ReLU, autoencoders, Neural ODEs, tensor factorization), and the deliverable is a landscape-level guarantee, not an objective rewrite or a faster algorithm.

\textbf{Differentiation.} \textit{Sibling cluster 28 (Principled Equivalence Recovery) also reformulates objectives via duality/decomposition, but its payoff is showing a heuristic objective EQUALS a classical one (e.g., SAM = Bayes relaxation, diffusion loss = ELBO) --- justifying the loss function. Cluster 14's payoff is geometric: under the reframing, the LANDSCAPE of an unchanged training objective becomes provably benign so first-order methods reach a global solution. Sibling 8 (Formal Abstraction for Expressivity) similarly applies algebraic frameworks to neural nets, but characterizes the static function class (what the architecture CAN compute), whereas cluster 14 characterizes the dynamic optimization trajectory and stationary-point structure. The signature distinction: this cluster's theorems are statements about gradient dynamics and Hessian structure, not about objective equivalence (28) or representational capacity (8).}

\textbf{When to pick.} Pick this instantiation when the empirical phenomenon you must explain is "first-order training reliably converges to a good solution despite non-convexity" and the architecture admits a known structural restriction (low-rank factorization, wide/infinite-width limit, piecewise-linear activation, algebraic tensor product). If the puzzle is instead "why does this objective work" ($\to$ cluster 28) or "what can this architecture express" ($\to$ cluster 8), pick those instead.

\textbf{Tactical failure.} Reject papers in this cluster fail by stopping one step short of the landscape-level theorem. NeurIPS\_2023\_0681 derives rigorous rank upper bounds on Jacobians from architectural choices but never converts these bounds into a no-spurious-minima or global-convergence statement, leaving the result as an architectural observation rather than an optimization guarantee. NeurIPS\_2025\_1889 proposes a UDV-decomposition architecture to make implicit bias explicit but offers convergence claims tied to design intuition rather than a rigorous critical-point characterization comparable to the Oral papers. ICLR\_2025\_1592 bounds the generalization gap at approximate stationary points but treats the non-convex landscape as a black box rather than collapsing it via a structural lever, so reviewers see a statistical extension rather than a landscape-reduction result. ICLR\_2024\_0946 argues by constructed examples that simplicity bias comes from architecture/initialization, but examples cannot substitute for the all-critical-points theorem this cluster's Orals deliver. The shared pathology: the structural insight is identified but the "provably benign landscape" payoff is asserted, illustrated, or partially bounded rather than fully proved end-to-end.

\textbf{Examples.}
\begin{itemize}[leftmargin=12pt,topsep=2pt,itemsep=1pt]
  \item \textbf{[Oral]} \texttt{ICML\_2024\_1059} \textemdash\ Canonical instance of the tactic: linearizes LoRA via NTK, then proves the low-rank factorization preserves the benign landscape so all local minima of the linearized loss are global.
  \\ \textcolor{gray}{\footnotesize \textit{Lesson:} Linearize LoRA via NTK and prove low-rank factorization preserves benign landscape --- all local minima of the linearized loss become global, the textbook landscape-collapse instance.}
  \item \textbf{[Oral]} \texttt{ICLR\_2022\_0141} \textemdash\ Constructs an EXACT (not approximate) convex program whose solutions are in bijection with all globally optimal two-layer ReLU networks --- the strongest possible form of structural collapse to a benign landscape.
  \\ \textcolor{gray}{\footnotesize \textit{Lesson:} EXACT (not approximate) convex program in bijection with all globally optimal two-layer ReLU networks --- the strongest possible structural collapse.}
  \item \textbf{[Oral]} \texttt{ICML\_2023\_0500} \textemdash\ Shows that lifting to higher dimensions does not eliminate spurious minima but converts them into strict saddles with negative curvature --- a distinctive landscape-geometry move within the cluster.
  \\ \textcolor{gray}{\footnotesize \textit{Lesson:} Lifting to higher dimensions converts spurious minima into strict saddles with negative curvature --- distinctive landscape-geometry move within the cluster.}
  \item \textbf{[Oral]} \texttt{ICML\_2024\_1014} \textemdash\ Proves an invariant-subspace theorem for the gradient flow of deep matrix factorization, reducing overparameterized training to a low-dimensional dynamical system that explains efficient convergence.
  \\ \textcolor{gray}{\footnotesize \textit{Lesson:} Invariant-subspace theorem for gradient flow of deep matrix factorization --- overparameterized training reduces to a low-dimensional dynamical system.}
  \item \textbf{[Reject]} \texttt{NeurIPS\_2023\_0681} \textemdash\ Identifies the structural lever (architecture-induced gradient-rank bounds) but stops at rank inequalities and never delivers the no-spurious-minima or global-convergence theorem the cluster's Orals provide.
  \\ \textcolor{gray}{\footnotesize \textit{Lesson:} Architecture-induced gradient-rank bounds without the no-spurious-minima or global-convergence theorem --- stops one step short of what the cluster requires.}
  \item \textbf{[Reject]} \texttt{NeurIPS\_2025\_1889} \textemdash\ Proposes a constrained UDV architecture to realize implicit low-rank bias explicitly, but the convergence claim is tied to design intuition rather than a rigorous critical-point characterization of the new landscape.
  \\ \textcolor{gray}{\footnotesize \textit{Lesson:} Constrained UDV architecture for implicit low-rank bias is design intuition, not a rigorous critical-point characterization of the new landscape.}
  \item \textbf{[Reject]} \texttt{ICLR\_2025\_1592} \textemdash\ Extends statistical guarantees to approximate stationary points but does not collapse the non-convex landscape via a structural lever, so it reads as an analytic extension rather than a landscape-reduction proof.
  \\ \textcolor{gray}{\footnotesize \textit{Lesson:} Statistical guarantees on approximate stationary points are analytic extension, not a structural lever collapsing the non-convex landscape.}
\end{itemize}
\end{tcolorbox}

\begin{tcolorbox}[colback=bgskill!50, colframe=cgreen!70, title={\textbf{C16}: Mechanism Reframing via Mathematical Equivalence \textcolor{gray}{\small (O13/H9/R18, conf: \textbf{high})}}, breakable, before skip=4pt, after skip=4pt]
\textbf{Tactical pattern.} Take a bottleneck operator (usually softmax attention, but also pairwise graph kernels, dense convolutions, or unrolled fixed-point iterations) and re-express it through a specific mathematical equivalence --- kernel decomposition (Performers, Random Feature Attention, Taming graph kernels), structured-matrix duality (Mamba/SSM-Attention duality), spectral/Fourier basis change (DiJiang, Fourmer, Spherical FNO, Embedding Fourier), implicit parameterization with hierarchical aggregation (Hyena, Pyraformer), or learned sparse sampling at carefully chosen reference points (Deformable DETR) --- so that the rewritten form preserves the operator's input-output behavior while collapsing its complexity. The verbs are concrete: factor a kernel into inner products of random features so an N$^2$ sum becomes a linear recurrence; swap a flat-space DFT for spherical harmonics on a curved domain; precompute neighborhood sketches that supply the same statistic a subgraph-GNN extracts expensively. The tactic is operator-level rewriting backed by an equivalence proof, not architectural substitution justified empirically.

\textbf{Differentiation.} \textit{Cluster 16 rewrites the COMPUTATIONAL OPERATOR (the actual matmul, softmax, or sum being run inside a layer) using a mathematical equivalence that delivers complexity wins --- distinct from sibling cluster 28 (Principled Equivalence Recovery via Objective Reformulation), which reframes the LOSS/OBJECTIVE to recover a classical principled framework (diffusion ELBO, SAM as Bayes relaxation) without changing how any inner operation is computed. It also differs from sibling cluster 8 (Formal Abstraction for Tight Expressivity Bounds), where the formal translation (algebraic structures, WL hierarchies) is in service of proving what the architecture CAN OR CANNOT EXPRESS, not of producing a cheaper implementation. In short: cluster 28 = "this objective is secretly that classical objective"; cluster 8 = "this architecture computes this formal class"; cluster 16 = "this operator equals that cheaper operator on the same inputs."}

\textbf{When to pick.} Pick this cluster when the target problem already has a working but expensive operator (quadratic attention, full subgraph extraction, dense convolution over a structured domain) AND you can identify a latent mathematical object inside it --- a kernel, a circulant/structured matrix, a transform-domain decomposition, a fixed-point --- that admits a known efficient form. Skip it when your candidate change is an architectural swap (replace module A with module B) without an equivalence between the two computations, or when the bottleneck is in the objective rather than the operator.

\textbf{Tactical failure.} Reject papers in this cluster fall into three recurrent traps. First, additive combination dressed as reframing: NeurIPS\_2025\_1843 (Fourier + wavelet transformer) and ICLR\_2023\_0214 (Fourier + dynamic edges) glue two known mechanisms together without identifying the unified mathematical object that makes one reduce to the other --- the "reframing" is just a wider architecture. Second, mechanism substitution without an equivalence proof: ICLR\_2023\_0296 (attention-as-memory), NeurIPS\_2025\_1849 (cross-attn for self-attn in speculative decoding), ICLR\_2025\_1483 (skeleton-sparsified attention), and NeurIPS\_2025\_1920 (recurrent memory bolted onto Decision Transformer) re-label the operator but never establish that the new operator computes the same function (or a controlled approximation) as the old one --- they rely on benchmark parity instead. Third, narrow algebraic improvements with no broader generalization: ICLR\_2023\_0180 (Hadamard $\to$ general matrix in bilinear pooling), ICLR\_2025\_1624 (reorder linear-attention matmuls for training memory), and ICLR\_2024\_0916 (quaternion-circulant $\leftrightarrow$ DFT) prove a real equivalence but on so narrow a slice that reviewers see a competent algebra exercise rather than a methodological move. The signature failure: the equivalence does no work --- it is decoration on what is otherwise an architectural or engineering tweak.

\textbf{Examples.}
\begin{itemize}[leftmargin=12pt,topsep=2pt,itemsep=1pt]
  \item \textbf{[Oral]} \texttt{ICML\_2023\_0469} \textemdash\ Replaces softmax attention with implicitly parameterized long convolutions plus element-wise gating, isolating attention's two essential properties (long-range mixing, data-adaptive weighting) and providing each via a separate sublinear primitive --- the canonical operator-level rewrite of this cluster.
  \\ \textcolor{gray}{\footnotesize \textit{Lesson:} Replace softmax attention with implicit long convolutions + element-wise gating --- isolate attention's two essential properties (long-range mixing, data-adaptive weighting) into separate sublinear primitives.}
  \item \textbf{[Oral]} \texttt{ICLR\_2021\_0057} \textemdash\ Constructs strictly positive orthogonal random features for the softmax kernel, factoring the N$^2$ attention sum into a product of linear operations with proven unbiasedness --- the textbook kernel-decomposition instantiation.
  \\ \textcolor{gray}{\footnotesize \textit{Lesson:} Strictly positive orthogonal random features factor N$^2$ attention into a product of linear operations with proven unbiasedness --- textbook kernel-decomposition reframing.}
  \item \textbf{[Oral]} \texttt{ICLR\_2021\_0014} \textemdash\ Reframes dense pairwise attention as learned sparse sampling at K data-adaptive reference offsets, exploiting that the same semantic role (localized aggregation) can be filled by a structurally different operator.
  \\ \textcolor{gray}{\footnotesize \textit{Lesson:} Reframe dense pairwise attention as learned sparse sampling at K data-adaptive reference offsets --- same semantic role, structurally different operator.}
  \item \textbf{[Oral]} \texttt{ICML\_2023\_0540} \textemdash\ Replaces the DFT inside a Fourier neural operator with spherical harmonics --- the basis change is a precise mathematical equivalence demanded by the curved domain, eliminating systematic pole/date-line artifacts.
  \\ \textcolor{gray}{\footnotesize \textit{Lesson:} Replace DFT with spherical harmonics inside Fourier neural operators --- basis change is a precise mathematical equivalence demanded by the curved domain.}
  \item \textbf{[Reject]} \texttt{ICML\_2025\_1669} \textemdash\ Stacks local adaptive weighting + global stationary propagation + implicit differentiation: each component is a valid reframing primitive but the combination yields no unifying equivalence beyond additive engineering gains.
  \\ \textcolor{gray}{\footnotesize \textit{Lesson:} Stacking local adaptive weighting + global stationary propagation + implicit differentiation yields no unifying equivalence beyond additive engineering gains.}
  \item \textbf{[Reject]} \texttt{NeurIPS\_2025\_1843} \textemdash\ Concatenates Fourier and wavelet transforms into a 'subquadratic transformer' without showing the union is mathematically equivalent to or a controlled approximation of self-attention --- the reframing is architectural, not operator-level.
  \\ \textcolor{gray}{\footnotesize \textit{Lesson:} Concatenating Fourier and wavelet transforms into a 'subquadratic transformer' without showing mathematical equivalence to attention is architectural reframing, not operator-level.}
  \item \textbf{[Reject]} \texttt{ICLR\_2025\_1483} \textemdash\ Sparsifies attention along skeletal-graph topology to cut cost, but provides no equivalence between the sparsified operator and full attention --- relies on benchmark parity, the canonical 'mechanism substitution without proof' failure.
  \\ \textcolor{gray}{\footnotesize \textit{Lesson:} Sparsifying attention along skeletal-graph topology without an equivalence proof is mechanism substitution relying on benchmark parity.}
\end{itemize}
\end{tcolorbox}

\begin{tcolorbox}[colback=bgskill!50, colframe=cgreen!70, title={\textbf{C22}: Generative Reframing of Sequential Decision Problems \textcolor{gray}{\small (O44/H5/R23, conf: \textbf{high})}}, breakable, before skip=4pt, after skip=4pt]
\textbf{Tactical pattern.} Take a problem traditionally solved by an optimizer, planner, or RL update rule (policy search, trajectory optimization, world-model rollout, exploration, IRL) and replace it with sampling from a learned conditional generative model --- typically a diffusion model, autoregressive Transformer, or latent-variable model --- over trajectories, actions, future states, or rewards. The verbs are concrete: train p(trajectory | return) and sample at high return (ICLR\_2023\_0256); tokenize state-action sequences and autoregress with a Transformer instead of recurrent dynamics (ICLR\_2023\_0379); use diffusion guidance gradients in place of value-function backups for plan refinement (ICML\_2023\_0409, ICLR\_2025\_1401); derive a variational lower bound so a latent generative model becomes a multimodal policy (ICML\_2023\_0519); reinterpret diffusion training as MaxEnt-IRL (NeurIPS\_2024\_1256). The optimizer is not augmented --- it is removed and its role taken over by conditional density estimation.

\textbf{Differentiation.} \textit{Versus sibling cluster 43 (Combinatorial Search via Sequential Decision Decomposition), which moves in the opposite direction --- taking discrete search and *adding* MDP structure plus tree/RL search --- this cluster *removes* the explicit search/optimization layer and replaces it with a generative sampler conditioned on the desired outcome. Versus sibling cluster 45 (Generative Framework Generalization via Mathematical Abstraction), which stays inside generative modeling and extends its mathematical formalism (e.g., kinetic-optimal flow matching, data-dependent stochastic interpolants), this cluster keeps generative models off-the-shelf and ports them into a foreign problem domain (RL/control/IRL). Versus sibling cluster 28 (Principled Equivalence Recovery via Objective Reformulation), which exposes a hidden equivalence between an existing heuristic and a classical objective, papers here are not retrospectively justifying an existing trick --- they are actively swapping out a Bellman-style or argmax-style solver for a learned sampler.}

\textbf{When to pick.} Pick this instantiation when the problem has rich offline trajectory or demonstration data with multimodal/multi-solution structure that a unimodal policy or value function flattens, AND when desired outcomes (return, goal, constraint, skill) can be expressed as a conditioning input rather than a scalar reward to maximize. If the bottleneck is search over a discrete combinatorial space with no demonstration data, instantiate as cluster 43 instead.

\textbf{Tactical failure.} Reject papers in this cluster overwhelmingly fail by bolting a generative model onto an existing RL pipeline as an auxiliary component rather than committing to it as the load-bearing solution concept --- the optimization framing survives, and the generative model is just a fancier prior, scorer, or data augmentation. NeurIPS\_2023\_0616 keeps behavioral cloning as the policy and uses diffusion only as a side scorer at inference, never replacing the conditional model. ICLR\_2023\_0298 swaps in a VQ-VAE codebook over actions but then runs vanilla DQN on top, treating the generative codebook as a preprocessing layer rather than the policy itself. NeurIPS\_2024\_1210 layers diffusion-generated performance trajectories on top of an unchanged hierarchical UED loop; ICML\_2025\_1788 and ICLR\_2023\_0220 use a learned generative/dynamics model to pre-screen or self-bootstrap candidates while leaving the surrounding QD or control objective untouched. ICLR\_2025\_1530 and NeurIPS\_2024\_1152 frame the contribution as iterative trajectory selection or sequence chunking rather than as a genuine generative reframing of the decision objective. The diagnostic signal: if removing the generative model leaves a recognizable RL/optimization algorithm intact, reviewers read it as incremental and the paper is rejected.

\textbf{Examples.}
\begin{itemize}[leftmargin=12pt,topsep=2pt,itemsep=1pt]
  \item \textbf{[Oral]} \texttt{ICLR\_2023\_0256} \textemdash\ Prototype paper for the cluster --- explicitly argues conditional generative modeling can entirely replace the optimization framing of decision-making and trains a diffusion model on return-conditioned trajectories instead of running policy improvement.
  \\ \textcolor{gray}{\footnotesize \textit{Lesson:} Argue conditional generative modeling can entirely REPLACE the optimization framing of decision-making --- train a diffusion model on return-conditioned trajectories instead of running policy improvement.}
  \item \textbf{[Oral]} \texttt{ICLR\_2023\_0379} \textemdash\ Replaces dynamical-system world-model rollouts with autoregressive next-token prediction over discretized state-action sequences, importing the generative-modeling paradigm wholesale into model-based RL rather than augmenting an existing planner.
  \\ \textcolor{gray}{\footnotesize \textit{Lesson:} Replace dynamical-system world-model rollouts with autoregressive next-token prediction over discretized state-action sequences --- import generative modeling wholesale into model-based RL.}
  \item \textbf{[Oral]} \texttt{ICML\_2023\_0409} \textemdash\ Recasts offline planning as guided diffusion sampling where reward gradients steer denoising toward high-return trajectories, with the diffusion sampler functioning as the planner itself rather than as a sample source for a downstream optimizer.
  \\ \textcolor{gray}{\footnotesize \textit{Lesson:} Recast offline planning as guided diffusion sampling where reward gradients steer denoising --- the diffusion sampler IS the planner, not a sample source for downstream optimization.}
  \item \textbf{[Oral]} \texttt{ICML\_2023\_0519} \textemdash\ Derives a variational lower bound that turns a latent-variable generative model into the policy proper, so the multimodal trajectory distribution is expressed and optimized as a generative density rather than as a Gaussian policy with auxiliary modes.
  \\ \textcolor{gray}{\footnotesize \textit{Lesson:} Variational lower bound that turns a latent-variable generative model into the policy --- multimodal trajectory distribution as a generative density, not Gaussian-with-auxiliary-modes.}
  \item \textbf{[Reject]} \texttt{NeurIPS\_2023\_0616} \textemdash\ Keeps behavioral cloning as the policy and uses a diffusion model only as an inference-time scorer over (state, action) joint density, so the generative reframing is auxiliary rather than load-bearing.
  \\ \textcolor{gray}{\footnotesize \textit{Lesson:} Keeping behavioral cloning as the policy and using diffusion only as inference-time scorer makes the generative reframing auxiliary, not load-bearing.}
  \item \textbf{[Reject]} \texttt{ICLR\_2023\_0298} \textemdash\ Applies a VQ-VAE codebook to compress continuous actions into discrete tokens but then runs standard DQN over the codebook, treating the generative model as preprocessing rather than as the decision-making mechanism.
  \\ \textcolor{gray}{\footnotesize \textit{Lesson:} VQ-VAE codebook + standard DQN over the codebook treats the generative model as preprocessing, not the decision-making mechanism.}
  \item \textbf{[Reject]} \texttt{NeurIPS\_2024\_1210} \textemdash\ Adds a diffusion model that synthesizes performance trajectories to assist a hierarchical UED teacher, leaving the underlying sparse-reward RL teacher loop intact instead of replacing the curriculum-search problem with a generative one.
  \\ \textcolor{gray}{\footnotesize \textit{Lesson:} A diffusion model assisting a hierarchical UED teacher leaves the underlying RL teacher loop intact --- the curriculum-search problem isn't replaced by a generative one.}
\end{itemize}
\end{tcolorbox}

\begin{tcolorbox}[colback=bgskill!50, colframe=cgreen!70, title={\textbf{C28}: Principled Equivalence Recovery via Objective Reformulation \textcolor{gray}{\small (O11/H0/R5, conf: \textbf{medium})}}, breakable, before skip=4pt, after skip=4pt]
\textbf{Tactical pattern.} The shared move is to take an *existing* empirically-developed heuristic objective (SAM update, diffusion loss with monotonic weighting, Expected Improvement, JMMD, etc.) and prove via a specific mathematical bridge --- Fenchel biconjugate, ELBO under implicit data augmentation, Tweedie/score reformulation, Representer Theorem, variational posterior approximation, or minimax-rate analysis --- that the heuristic *is exactly* (or a special case of) a known principled object. The signature verbs are "show that L\_heuristic = L\_classical", "recover X as a special case", "the existing update equals the optimum of the dual". Crucially, the classical target is named in advance (ELBO, Bayes, MMD, minimax-optimal estimator, JSD) and the bridge is a single invocation of a textbook tool, not a new framework. The payoff is retroactive justification plus a controlled generalization (e.g., $\alpha$-$\beta$-divergence, VES-Gamma) obtained by perturbing one ingredient of the now-revealed classical object.

\textbf{Differentiation.} \textit{Versus sibling cluster 16 (Mechanism Reframing via Mathematical Equivalence), that cluster reformulates *core operators* (attention, message-passing) to obtain cheaper implementations --- the goal is computational, and the original and reformulated mechanisms can have different behaviors that happen to coincide on a regime of interest. Cluster 28 instead reformulates *training/inference objectives* and proves *exact* equivalence to a named classical object so the heuristic inherits its guarantees --- no efficiency claim is needed. Versus cluster 45 (Generative Framework Generalization), cluster 45 moves *forward* --- relaxing a constraint in an already-principled framework to broaden it. Cluster 28 moves *backward* --- taking something with no theory and showing it was already an instance of theory all along. Versus cluster 8 (Formal Abstraction for Expressivity), cluster 8 derives capability bounds for *architectures* via algebra/circuits, whereas cluster 28 derives an equivalence for a scalar *objective*.}

\textbf{When to pick.} Pick this tactic when the target is a single empirically successful loss/acquisition/regularizer/update rule that has no probabilistic story, and you can name in advance the classical object it most resembles (ELBO, Bayes posterior, KSD, MMD, minimax-optimal density estimator). Skip in favor of siblings if the target is an architecture ($\to$ cluster 8/14) or a mechanism that needs to be made cheaper rather than justified ($\to$ cluster 16).

\textbf{Tactical failure.} The dominant failure mode in the reject papers is *substitution dressed as equivalence*: replacing one heuristic with a slightly more principled-looking variant without recovering an exact, named classical object. ICML\_2025\_1680 swaps KL for an f\_$\chi$$^n$ regularizer, motivating it via gradient-field analysis but not deriving it as the unique optimum of any classical principle; ICLR\_2024\_0813 replaces a supremum-based distributional Bellman error with an expectation-based energy distance to dodge a vacuous constant --- an analytic tightening rather than a recovery of an existing object as a known optimum. ICML\_2025\_1779 bundles three fixes under one length-regularized DPO-COV objective; the "equivalence" surface (DPO-COV $\leftrightarrow$ RLHF-COV) is internal to the new construction rather than a bridge to classical statistics. ICLR\_2024\_0941 transports score-based convergence proofs to consistency models --- an *adaptation* of analysis machinery, not an equivalence to a textbook framework. NeurIPS\_2025\_1923 does derive an equivalence (to single-sample fitting) but the equivalence demystifies negatively rather than legitimizing the heuristic --- reviewers want the punchline "the heuristic was always optimal", not "the heuristic was always degenerate". The common thread: lacking a clean "L\_heuristic = L\_classical\_object" identity that a reviewer can verify in one line.

\textbf{Examples.}
\begin{itemize}[leftmargin=12pt,topsep=2pt,itemsep=1pt]
  \item \textbf{[Oral]} \texttt{ICLR\_2023\_0331} \textemdash\ Canonical instance of the tactic --- applies the Fenchel biconjugate to the expected log-loss in the Bayes objective and shows the resulting minimax relaxation is *exactly* the SAM update, converting an empirically-derived sharpness heuristic into a one-line consequence of convex duality.
  \\ \textcolor{gray}{\footnotesize \textit{Lesson:} Apply Fenchel biconjugate to expected log-loss in Bayes objective --- minimax relaxation IS exactly SAM, an empirical heuristic recovered as a one-line consequence of convex duality.}
  \item \textbf{[Oral]} \texttt{NeurIPS\_2023\_0753} \textemdash\ Identifies an implicit data augmentation that makes the family of monotonically-weighted diffusion training objectives equal the ELBO of an augmented model, recovering a perceptual heuristic as a principled likelihood objective.
  \\ \textcolor{gray}{\footnotesize \textit{Lesson:} Identify implicit data augmentation making monotonically-weighted diffusion objectives equal the ELBO of an augmented model --- recover a perceptual heuristic as a principled likelihood.}
  \item \textbf{[Oral]} \texttt{ICML\_2025\_1662} \textemdash\ Reinterprets Expected Improvement as variational inference over the optimum distribution with a specific approximating family --- EI is recovered exactly as a special case, and swapping the variational family yields the principled VES-Gamma extension.
  \\ \textcolor{gray}{\footnotesize \textit{Lesson:} Reinterpret Expected Improvement as variational inference over the optimum distribution --- EI recovered exactly as a special case, swapping the variational family yields the principled extension.}
  \item \textbf{[Oral]} \texttt{ICML\_2024\_1118} \textemdash\ Uses Tweedie/score relationships to rewrite the entropy-based definition of O-information into a score-only expression, showing the same target quantity equals an estimable score functional and bypassing density estimation without redefining the object.
  \\ \textcolor{gray}{\footnotesize \textit{Lesson:} Tweedie/score relationships rewrite entropy-based O-information into a score-only expression --- same target equals an estimable score functional, bypassing density estimation.}
  \item \textbf{[Reject]} \texttt{ICML\_2025\_1680} \textemdash\ Substitutes KL with an f\_$\chi$$^n$ regularizer motivated by gradient-field analysis but never recovers any classical object as a special case --- the work reads as a better-tuned heuristic rather than a principled equivalence.
  \\ \textcolor{gray}{\footnotesize \textit{Lesson:} f\_$\chi$$^n$ regularizer motivated by gradient-field analysis without recovering any classical object as a special case is a better-tuned heuristic, not a principled equivalence.}
  \item \textbf{[Reject]} \texttt{ICLR\_2024\_0813} \textemdash\ Tightens a distributional Bellman bound by replacing a supremum metric with an expectation-based energy distance --- an analytic improvement, but the new objective is not shown to *be* a known classical estimator, only a less vacuous one.
  \\ \textcolor{gray}{\footnotesize \textit{Lesson:} Replacing supremum with energy distance is analytic improvement, not equivalence --- the new objective must BE a known classical estimator, not just a less vacuous one.}
  \item \textbf{[Reject]} \texttt{NeurIPS\_2025\_1923} \textemdash\ Derives an equivalence (high-dimensional diffusion objective collapses to single-sample fitting) that demystifies the heuristic as degenerate rather than recovering it as an instance of a principled classical framework --- same tactical move, opposite rhetorical payoff.
  \\ \textcolor{gray}{\footnotesize \textit{Lesson:} Equivalence showing the target collapses to single-sample fitting demystifies as degenerate rather than recovering as a principled instance --- same tactical move, opposite payoff.}
\end{itemize}
\end{tcolorbox}

\begin{tcolorbox}[colback=bgskill!50, colframe=cgreen!70, title={\textbf{C33}: Reframing Through Structural Analogy and Decomposition \textcolor{gray}{\small (O10/H2/R1, conf: \textbf{low})}}, breakable, before skip=4pt, after skip=4pt]
\textbf{Tactical pattern.} The shared move is to take a debate or measurement problem that is stuck on a binary/unmeasurable framing and introduce a NEW analytical axis along which the original problem decomposes into tractable sub-problems --- or to map the problem onto a structurally analogous domain with established solutions. Concrete instances: split "open vs closed AI" along capability$\times$timescale axes (ICML\_2024\_1090) or along counterfactual-availability (ICML\_2024\_1091); replace "AI personhood" with causal accountability chains traced to specific human actors (ICML\_2024\_1085); convert an unmeasurable normative quantity ("is X more harmful?") into a falsifiable surrogate ("is X harder to counteract?", ICLR\_2024\_0785); reframe individual-instance detection as population-level mixture prevalence (ICML\_2024\_1069); or import cybersecurity's Coordinated Vulnerability Disclosure as an analog for AI safe harbor (ICML\_2024\_1082). Crucially, the "formalism" being substituted is conceptual/policy-analytic (a new axis, a domain analogy, a complexity-class constraint, a population-vs-individual reframe), NOT a mathematical equivalence --- and the target is almost always a societal, governance, or measurement question rather than a technical algorithm.

\textbf{Differentiation.} \textit{Where siblings reformulate technical objects through formal mathematical equivalences --- Cluster 6 maps optimization onto Nash games with no-regret solvers, Cluster 16 swaps attention for structured-matrix duality, Cluster 28 reveals an empirical objective IS a known classical objective via convex duality --- Cluster 33 operates on policy/safety/measurement debates and substitutes a CONCEPTUAL reframing rather than a mathematical isomorphism. The papers do not "solve" the original question; they dissolve a stuck framing by introducing a missing axis (counterfactual, timescale, causal chain) or by importing an analogy (cybersecurity CVD $\to$ AI safe harbor in ICML\_2024\_1082). Compared to Cluster 8's tight expressivity bounds via algebraic frameworks, Cluster 33 rarely yields tight quantitative results --- its outputs are taxonomies, accountability chains, or capability-conditional recommendations. The closest sibling is Cluster 28 (Principled Equivalence Recovery), but Cluster 28 requires showing that two formal objects are mathematically equivalent, while Cluster 33's analogies are structural-conceptual (e.g., research-access barriers in AI parallel those in pre-CVD cybersecurity) and the deliverable is a reframed policy question, not a recovered theorem.}

\textbf{When to pick.} Pick this tactic when the target is a policy, safety, governance, or measurement debate that is deadlocked on a binary framing or an unmeasurable normative quantity --- and you can name a missing analytical axis (timescale, counterfactual availability, capability gap, causal chain, population vs individual) or a structurally analogous domain whose solution can be transferred. If the target is a technical algorithm or architecture rather than a discourse-level question, prefer a sibling cluster that uses a true mathematical reformulation.

\textbf{Tactical failure.} With only one reject in this cluster, the failure characterization is necessarily speculative. The single reject (ICLR\_2023\_0242) attempts the spirit of the tactic --- proactively framing a not-yet-realized misuse problem (generated graph detection) as a classification task before harm materializes --- but the move stops at problem-formalization and the paper falls back on standard discriminative classifiers benchmarked across generalization scenarios. It does not introduce a new conceptual axis (no counterfactual, no causal chain, no domain analogy that brings established solutions) and offers no reframing that explains why the community's current framing was wrong. Reviewers reject because the contribution reduces to "apply known methods to an anticipated problem," and absent demonstrated impact the framing alone is not load-bearing. The implied failure mode for this tactic: proposing a reframing that is merely a relabeling of the existing problem rather than one that exposes a previously-conflated axis or transports a non-trivial solution from an analog domain. Confidence is low because n\_reject=1 prevents pattern-level inference.

\textbf{Examples.}
\begin{itemize}[leftmargin=12pt,topsep=2pt,itemsep=1pt]
  \item \textbf{[Oral]} \texttt{ICML\_2024\_1091} \textemdash\ Exemplifies the cluster's signature move of introducing a missing conceptual axis (counterfactual availability) that dissolves a binary open-vs-closed framing into a multidimensional impact assessment --- the reframing, not any mathematical result, is the contribution.
  \\ \textcolor{gray}{\footnotesize \textit{Lesson:} Counterfactual availability dissolves a binary open-vs-closed framing into multi-dimensional impact assessment --- the reframing IS the contribution.}
  \item \textbf{[Oral]} \texttt{ICML\_2024\_1082} \textemdash\ Pure structural-analogy instance: maps AI red-teaming's legal barriers onto pre-CVD cybersecurity's research-access problem and transports the safe-harbor solution across domains, with adaptation for AI-specific properties.
  \\ \textcolor{gray}{\footnotesize \textit{Lesson:} Map AI red-teaming's legal barriers onto pre-CVD cybersecurity's research-access problem and transport the safe-harbor solution --- pure structural-analogy.}
  \item \textbf{[Oral]} \texttt{ICML\_2024\_1085} \textemdash\ Replaces an unresolvable foundational question (AI personhood) with a tractable functional substitute (causal accountability chains traced to human actors) --- canonical example of decomposing a stuck normative debate by switching the analytical primitive.
  \\ \textcolor{gray}{\footnotesize \textit{Lesson:} Replace AI personhood with causal accountability chains traced to human actors --- decompose a stuck normative debate by switching the analytical primitive.}
  \item \textbf{[Oral]} \texttt{ICML\_2024\_1069} \textemdash\ Reframes unreliable individual-instance detection as population-level mixture-proportion estimation, swapping the analytical level rather than improving the classifier --- a cluster-defining 'change the axis, not the algorithm' move.
  \\ \textcolor{gray}{\footnotesize \textit{Lesson:} Recast unreliable individual-instance detection as population-level mixture-proportion estimation --- change the analytical level, not the algorithm.}
  \item \textbf{[Reject]} \texttt{ICLR\_2023\_0242} \textemdash\ Attempts proactive problem reframing (generated-graph detection before misuse peaks) but stops at formalization and benchmarks standard classifiers, never introducing a new conceptual axis or domain analogy that would make the reframing load-bearing.
  \\ \textcolor{gray}{\footnotesize \textit{Lesson:} Proactive problem reframing must introduce a NEW conceptual axis or domain analogy --- formalization + standard classifiers without that is description, not reframing.}
\end{itemize}
\end{tcolorbox}

\begin{tcolorbox}[colback=bgskill!50, colframe=cgreen!70, title={\textbf{C43}: Combinatorial Search via Sequential Decision Decomposition \textcolor{gray}{\small (O18/H3/R14, conf: \textbf{high})}}, breakable, before skip=4pt, after skip=4pt]
\textbf{Tactical pattern.} Take a discrete combinatorial search problem that has historically been attacked by enumeration, hand-coded heuristics, or manual design --- neural architectures, symbolic expressions, agent workflows, prompt programs, logic rules, instruction curricula --- and explicitly construct a sequential decision process out of it (state = partial construction, action = atomic extension, reward = incremental fit/feasibility/feedback), then attach a guided search engine that exploits this decomposition: MCTS over expression grammars (ICLR\_2023\_0356, ICLR\_2025\_1348), evolutionary search whose mutation operators are LLM code edits or grammar rewrites (ICML\_2024\_1031, ICML\_2024\_1107, ICLR\_2025\_1493), RL/policy-gradient over construction traces (ICLR\_2022\_0144, ICLR\_2025\_1351), or differentiable tree construction via stochastic branching (ICLR\_2025\_1487). The shared move is converting "search a static space of structures" into "build a structure step-by-step under guidance," which buys principled exploration-exploitation, intermediate scoring, and the ability to inject learned priors at each decision rather than only at the leaf.

\textbf{Differentiation.} \textit{Versus sibling cluster 22 (Generative Reframing of Sequential Decision Problems), the directionality is inverted: cluster 22 dissolves an existing optimization/planning problem by replacing it with conditional generative sampling, whereas cluster 43 manufactures a sequential decision process where none existed (the original problem is static enumeration over architectures, expressions, or rule sets) so that RL, MCTS, and evolutionary machinery become applicable. Versus cluster 28 (Principled Equivalence Recovery via Objective Reformulation), this cluster leaves the objective intact and changes the algorithmic substrate; cluster 28 leaves the algorithm intact and rewrites the objective to expose a classical equivalence. Versus cluster 6 (Game-Theoretic Equilibrium Reformulation), cluster 43 stays single-agent and constructive --- there is no second player and no equilibrium concept; the search agent navigates a tree of build-actions, not a strategic interaction.}

\textbf{When to pick.} Pick this instantiation when the target object can be naturally built by a sequence of compositional choices (token-by-token, layer-by-layer, edge-by-edge, clause-by-clause) AND partial constructions admit a cheap incremental score or feasibility check --- these two together make tree/evolutionary/RL search strictly cheaper than enumeration and let learned priors steer at every step rather than only at the leaf.

\textbf{Tactical failure.} Reject papers in this cluster reveal three recurring traps when the recipe is followed mechanically. (1) Stock-MDP-with-stock-search: the authors simply name a state, action, and reward and plug in standard RL/NAS/MCTS without addressing the structural pathology of their domain --- NeurIPS\_2023\_0614 (DensEMANN) couples arch growth to training signals but reviewers see a finite-state heuristic, ICLR\_2024\_0926 just runs evolutionary NAS over backward-connection placements, ICLR\_2023\_0322 reformulates rule discovery as RL without taming the combinatorial action space, and ICLR\_2025\_1411 reformulates quantum scheduling as an MDP without coping with the stochastic dynamics that block sample-efficient convergence. (2) Metaphor in place of mechanism: NeurIPS\_2025\_1930 dresses evolutionary NAS in a "social hierarchy" framing without theoretical justification for why rank-stratified operators should converge faster. (3) Classical-technique-rebrand: ICML\_2025\_1758 (msf-CNN) presents exhaustive graph traversal of a configuration DAG as algorithmic novelty when reviewers correctly classify it as standard systems-engineering search; NeurIPS\_2023\_0668 (LATS) and NeurIPS\_2023\_0671 (ToT-guided) suffer the parallel problem in the LLM space, where attaching tree search to a generator is read as engineering composition rather than a tactical insight. The unifying signal: the reformulation alone is treated as the contribution, with no problem-specific structural lever (action pruning, learned mutation operator, differentiable construction, multi-granularity candidates) that distinguishes the move from off-the-shelf RL/NAS/MCTS.

\textbf{Examples.}
\begin{itemize}[leftmargin=12pt,topsep=2pt,itemsep=1pt]
  \item \textbf{[Oral]} \texttt{ICLR\_2023\_0356} \textemdash\ Canonical instantiation --- recasts symbolic equation discovery as MCTS over a grammar tree with parsimony as a structural prior, exemplifying state=partial-expression, action=production-rule, reward=data-fit.
  \\ \textcolor{gray}{\footnotesize \textit{Lesson:} Recast symbolic equation discovery as MCTS over a grammar tree with parsimony as structural prior --- state=partial-expression, action=production-rule, reward=data-fit.}
  \item \textbf{[Oral]} \texttt{ICML\_2024\_1031} \textemdash\ Pushes the tactic by replacing template-based mutation operators with LLM-generated code edits, dramatically expanding the reachable search space within the evolutionary-search variant of this cluster.
  \\ \textcolor{gray}{\footnotesize \textit{Lesson:} Replace template mutation with LLM-generated code edits --- dramatically expand the reachable search space within evolutionary search.}
  \item \textbf{[Oral]} \texttt{ICLR\_2025\_1348} \textemdash\ Represents agent workflows as executable code so that MCTS modifications are valid programs and partial workflows can be scored by execution feedback --- the executable-code substrate is what makes tree search applicable.
  \\ \textcolor{gray}{\footnotesize \textit{Lesson:} Represent agent workflows as executable code so MCTS modifications are valid programs --- execution feedback scores partial workflows.}
  \item \textbf{[Oral]} \texttt{ICLR\_2025\_1487} \textemdash\ Makes the search itself end-to-end differentiable by parameterizing best-first tree construction as a stochastic process, showing the cluster's tactic is compatible with gradient-based learning of the search policy from demonstrations.
  \\ \textcolor{gray}{\footnotesize \textit{Lesson:} Make the search itself end-to-end differentiable by parameterizing best-first tree construction --- gradient-learnable search policy from demonstrations.}
  \item \textbf{[Reject]} \texttt{NeurIPS\_2025\_1930} \textemdash\ Imposes a 'social hierarchy' metaphor on evolutionary NAS without principled justification --- the classic failure of dressing an off-the-shelf evolutionary search in a new analogy and treating that as the contribution.
  \\ \textcolor{gray}{\footnotesize \textit{Lesson:} Imposing a 'social hierarchy' metaphor on evolutionary NAS without principled justification dresses an off-the-shelf search in a new analogy.}
  \item \textbf{[Reject]} \texttt{ICML\_2025\_1758} \textemdash\ Reviewers correctly read exhaustive DAG traversal of a configuration space as a systems-engineering technique applied to a new domain, not a tactical novelty in combinatorial-search-via-decomposition.
  \\ \textcolor{gray}{\footnotesize \textit{Lesson:} Exhaustive DAG traversal of a configuration space is systems-engineering applied to a new domain, not tactical novelty in combinatorial search.}
  \item \textbf{[Reject]} \texttt{ICLR\_2025\_1411} \textemdash\ Mechanical MDP reformulation of quantum resource scheduling --- names state/action/reward and applies stock RL without addressing the stochastic dynamics that prevent sample-efficient convergence in the reformulated problem.
  \\ \textcolor{gray}{\footnotesize \textit{Lesson:} Mechanical MDP reformulation of quantum scheduling --- names state/action/reward and applies stock RL without addressing the stochastic dynamics that block sample-efficient convergence.}
\end{itemize}
\end{tcolorbox}

\begin{tcolorbox}[colback=bgskill!50, colframe=cgreen!70, title={\textbf{C45}: Generative Framework Generalization via Mathematical Abstraction \textcolor{gray}{\small (O17/H0/R3, conf: \textbf{low})}}, breakable, before skip=4pt, after skip=4pt]
\textbf{Tactical pattern.} Take an established generative framework (diffusion SDE, flow matching, GFlowNet, score-based model) and identify a more general mathematical object or relaxed constraint that contains the original as a special case, then re-derive the training/sampling machinery on the abstracted object. Concrete moves include: replacing vector fields with infinitesimal generators of arbitrary Markov processes (ICLR\_2025\_1437), lifting the Bellman equation from scalars to probability paths solved via flow matching (ICML\_2025\_1804), changing the modeled quantity from instantaneous to average velocity to derive a closed-form self-consistency identity (ICLR\_2025\_1502), dropping mass conservation in flow matching to get a tunable efficiency-coverage tradeoff (ICLR\_2025\_1378), generalizing flow matching to discrete paths via a kinetic-optimal velocity criterion (ICLR\_2025\_1434), or fusing mean-field interaction SDEs with score-based diffusion SDEs because both are distribution-valued dynamics (ICML\_2024\_1065). The result must preserve (or extend) the original framework's theoretical guarantees while opening it to problem classes the original could not handle.

\textbf{Differentiation.} \textit{Unlike cluster 22 (Generative Reframing of Sequential Decision Problems), which moves a non-generative problem like RL planning INTO the generative paradigm, cluster 45 starts inside generative modeling and pushes its mathematical scaffolding outward --- the target is the framework itself, not a downstream task. Unlike cluster 28 (Principled Equivalence Recovery via Objective Reformulation), which proves an existing empirical objective is secretly equivalent to a known classical principle (e.g., NeurIPS\_2023\_0753 showing diffusion objectives are ELBOs), cluster 45 papers do not unify with a classical equivalent --- they construct a strictly larger object (generators $\supset$ vector fields, average $\supset$ instantaneous velocity, distributional $\supset$ scalar Bellman) and inherit derivations into the new regime. Unlike cluster 16 (Mechanism Reframing via Mathematical Equivalence), which swaps out a single bottleneck operation (attention, message passing) inside a fixed architecture, cluster 45 reformulates the entire generative process at the level of its governing equation.}

\textbf{When to pick.} Pick this tactic when your starting point is already a working generative framework whose formulation rests on a specific mathematical object (a vector field, an SDE, a scalar value, a balanced transport plan), AND a strictly more general object exists in an adjacent literature (Markov process theory, optimal transport, mean-field physics, dynamic programming on distributions) that subsumes the original as a special case. Avoid this tactic if you are merely tuning hyperparameters of an existing coupling or schedule --- that lands in the failure mode below.

\textbf{Tactical failure.} The dominant failure mode is offering a parametric tweak to source-target coupling or potential design while branding it as a "generalization." All three rejects (ICLR\_2024\_0954, ICML\_2025\_1689, NeurIPS\_2025\_1907) modify how source-target pairs are constructed in flow/interpolant training --- data-dependent couplings, Lagrangian-derived velocity fields, and Sinkhorn-optimal pairings respectively --- but none replaces the framework's underlying mathematical object with a strictly larger one, and none derives new closed-form identities or guarantees that the original could not produce. Reviewers read these as variance-reduction or coupling-engineering papers rather than abstractions: the "more general" framing collapses on inspection because the generalized object is still a flow/interpolant with adjusted endpoints, the empirical gain is task-localized (imputation, conditional generation), and no theorem demonstrates a capability outside the original framework's reach. With only 3 rejects in the cluster, this characterization is suggestive rather than definitive --- confidence is correspondingly limited.

\textbf{Examples.}
\begin{itemize}[leftmargin=12pt,topsep=2pt,itemsep=1pt]
  \item \textbf{[Oral]} \texttt{ICLR\_2025\_1437} \textemdash\ Canonical example of the tactic: replaces vector fields with infinitesimal generators of arbitrary Markov processes, the strictly more general mathematical object that contains continuous flows as one special case while inheriting the framework's theoretical scaffolding.
  \\ \textcolor{gray}{\footnotesize \textit{Lesson:} Replace vector fields with infinitesimal generators of arbitrary Markov processes --- a strictly more general object that contains continuous flows as one special case.}
  \item \textbf{[Oral]} \texttt{ICLR\_2025\_1502} \textemdash\ Changes the modeled quantity from instantaneous to average velocity and derives the MeanFlow Identity --- a new closed-form differential relation that did not exist in the original formulation, exemplifying how abstracting the learned object unlocks novel training constraints.
  \\ \textcolor{gray}{\footnotesize \textit{Lesson:} Change modeled quantity from instantaneous to average velocity; derive the MeanFlow Identity --- a new closed-form differential relation that did not exist in the original formulation.}
  \item \textbf{[Oral]} \texttt{ICML\_2025\_1804} \textemdash\ Lifts the Bellman equation from scalar values to probability paths and solves the resulting consistency equation with flow matching, demonstrating the cluster's pattern of generalizing a foundational equation to operate on distribution-valued objects.
  \\ \textcolor{gray}{\footnotesize \textit{Lesson:} Lift the Bellman equation from scalar values to probability paths and solve with flow matching --- generalize a foundational equation to operate on distribution-valued objects.}
  \item \textbf{[Oral]} \texttt{ICLR\_2025\_1378} \textemdash\ Relaxes the mass-conservation constraint that flow matching had treated as essential, formalizing the relaxation degree as a tunable knob with provable efficiency-approximation tradeoffs --- a constraint-removal abstraction rather than a coupling tweak.
  \\ \textcolor{gray}{\footnotesize \textit{Lesson:} Relax the mass-conservation constraint flow matching treated as essential --- formalize relaxation degree as a tunable knob with provable efficiency-approximation tradeoffs.}
  \item \textbf{[Reject]} \texttt{ICLR\_2024\_0954} \textemdash\ Frames data-dependent couplings as a generalization of stochastic interpolants, but the modification is ultimately a reparameterization of the base distribution --- no new mathematical object is introduced and no capability outside the original framework is established.
  \\ \textcolor{gray}{\footnotesize \textit{Lesson:} Data-dependent couplings as 'generalization' --- reparameterizing the base distribution introduces no new mathematical object and no capability outside the original framework.}
  \item \textbf{[Reject]} \texttt{NeurIPS\_2025\_1907} \textemdash\ Replaces random source-target pairings with Sinkhorn-optimal ones to reduce velocity-field variance, which is a coupling-engineering improvement to flow training rather than the abstraction-of-the-governing-equation move that distinguishes the successful cluster members.
  \\ \textcolor{gray}{\footnotesize \textit{Lesson:} Sinkhorn-optimal source-target pairings for variance reduction is coupling engineering, not the abstraction-of-the-governing-equation move.}
  \item \textbf{[Reject]} \texttt{ICML\_2025\_1689} \textemdash\ Imports a Lagrangian-mechanics flavored velocity field into flow-based imputation as a variance-reduction heuristic, but does not yield a strictly larger generative framework or new theoretical identity --- the contribution is task-localized rather than abstraction-level.
  \\ \textcolor{gray}{\footnotesize \textit{Lesson:} Lagrangian-mechanics velocity field for variance reduction is task-localized, not abstraction-level --- yields no strictly larger framework or new theoretical identity.}
\end{itemize}
\end{tcolorbox}

\subsection{Decompose and Recombine}
\textit{Parent methodology id: \texttt{decomposition\_into\_tractable\_stages}; 5 tactical clusters.}

\begin{tcolorbox}[colback=bgskill!50, colframe=cgreen!70, title={\textbf{C07}: Timescale Separation for Staged Dynamics Analysis \textcolor{gray}{\small (O11/H0/R0, conf: \textbf{low})}}, breakable, before skip=4pt, after skip=4pt]
\textbf{Tactical pattern.} The shared move is: take a coupled nonlinear training/inference dynamical system, identify two-or-three natural timescales (or phases) inside it, and then derive each phase's behavior via a closed-form analytical reduction --- typically mean-field PDEs (NeurIPS\_2025\_1830, ICLR\_2025\_1417, ICML\_2024\_1138, NeurIPS\_2025\_1864), bifurcation/fixed-point analysis (NeurIPS\_2025\_1938, NeurIPS\_2025\_1862), or two-timescale separation in linearized/NTK-like limits (NeurIPS\_2025\_1873, ICML\_2024\_1138). The verbs are concrete: "linearize", "take mean-field limit", "compute fixed points and stability", "prove timescale gap grows with n". The objects are concrete: condensation-then-rank-collapse stages, tau\_gen vs tau\_mem windows, fast-attention/slow-FFN dynamics, sparsity bifurcation points. The output is always a sharp analytical characterization of qualitatively distinct regimes (cluster, plateau, collapse, generalize, memorize) plus a quantitative law for when each regime ends.

\textbf{Differentiation.} \textit{Versus sibling 41 (Decomposed Intermediate State Chaining), which decomposes reasoning into discrete sub-steps with verifiable intermediate inputs/outputs, this cluster decomposes a continuous dynamical trajectory into time-ordered phases analyzed by physics-style PDE/bifurcation tools --- not by routing engineered intermediates between modules. Versus sibling 42 (Decompose-then-Recombine Multi-Objective Search), which splits a scalar objective into sub-objectives solved by separate optimizers and recombined into a Pareto front, cluster 7 keeps a single objective but splits its temporal evolution; the "recombination" is analytical (composing the two phases into a full loss curve as in ICML\_2025\_1738) rather than algorithmic. Versus sibling 13 (global-via-local-proxy under privacy/distribution constraints) and sibling 44 (factorizing generative tasks into disentangled latent subproblems), cluster 7's split is along the time/dynamics axis with closed-form regime characterizations as the deliverable, whereas siblings split along agents/modalities and deliver constructive algorithms.}

\textbf{When to pick.} Pick this tactic when your phenomenon is a single training/inference run that empirically shows phase-like behavior (plateau-then-drop, generalize-then-memorize, condense-then-collapse, abrupt sparsification) and you can plausibly take a tractable scaling limit (mean-field, NTK, large-N, linearized loss) in which the dynamics decouple into fast and slow components. Choose a sibling instead if your decomposition axis is agents/data shards (sibling 13), reasoning steps with discrete inputs/outputs (sibling 41), competing objectives (sibling 42), or generative output components (sibling 44).

\textbf{Tactical failure.} This cluster contains 0 Reject papers, so the tactical failure mode is necessarily speculative --- confidence is set to low on this dimension. Based on the parent methodology's general failure pattern and the structure of the Oral submissions, plausible failure modes for this specific tactic include: (a) deriving timescale-separated PDEs in a toy/linearized regime without empirically demonstrating that real architectures exhibit the same phase structure (the Oral papers all anchor their analytical phases to observed transformer/diffusion behavior, e.g. NeurIPS\_2025\_1830, NeurIPS\_2025\_1945); (b) producing a multi-phase decomposition that is qualitative/post-hoc rather than yielding a quantitative law for the phase boundary (the Orals all deliver sharp predictions --- bifurcation thresholds, tau\_mem=Omega(n), power-law plateau lengths); (c) assuming independence of timescales without proving it, where the Orals like ICML\_2025\_1738 explicitly verify multiplicative decomposability. No reject paper\_ids can be cited because none exist in this cluster's data.

\textbf{Examples.}
\begin{itemize}[leftmargin=12pt,topsep=2pt,itemsep=1pt]
  \item \textbf{[Oral]} \texttt{NeurIPS\_2025\_1873} \textemdash\ Cleanest exemplar of the two-stage decomposition tactic: linearizes the architecture, identifies Stage-1 outer-layer condensation and Stage-2 inner-layer rank collapse as sequential rather than simultaneous, with each stage's dynamics characterized via Jacobian structure.
  \\ \textcolor{gray}{\footnotesize \textit{Lesson:} Linearize the architecture, separate Stage-1 outer-layer condensation from Stage-2 inner-layer rank collapse --- sequential rather than simultaneous, each characterized via Jacobian structure.}
  \item \textbf{[Oral]} \texttt{NeurIPS\_2025\_1945} \textemdash\ Quintessential timescale-gap argument: defines tau\_gen and tau\_mem analytically and proves their ratio grows with dataset size n, turning early stopping into a derived consequence of two-timescale separation rather than an empirical heuristic.
  \\ \textcolor{gray}{\footnotesize \textit{Lesson:} Prove $\tau$\_gen / $\tau$\_mem ratio grows with dataset size n --- early stopping becomes a derived consequence of two-timescale separation, not an empirical heuristic.}
  \item \textbf{[Oral]} \texttt{ICML\_2024\_1138} \textemdash\ Most explicit statement of the tactic's machinery: identifies that FFN parameters evolve on a slower effective timescale than attention parameters in the NTK-like limit, exploits the separation to reduce a nonconvex coupled ODE to a tractable mean-field PDE.
  \\ \textcolor{gray}{\footnotesize \textit{Lesson:} FFN parameters evolve on a slower effective timescale than attention; exploit the gap to reduce a coupled nonconvex ODE to a tractable mean-field PDE.}
  \item \textbf{[Oral]} \texttt{NeurIPS\_2025\_1862} \textemdash\ Pure bifurcation-analysis instantiation: locates the learning-rate threshold at which SGB softmax dynamics undergo a fixed-point bifurcation, yielding a sharp regret regime boundary that O(1/T) prescriptions miss.
  \\ \textcolor{gray}{\footnotesize \textit{Lesson:} Locate the LR threshold where SGB softmax dynamics undergo a fixed-point bifurcation --- sharp regret regime boundaries that O(1/T) analyses miss.}
\end{itemize}
\end{tcolorbox}

\begin{tcolorbox}[colback=bgskill!50, colframe=cgreen!70, title={\textbf{C13}: Global Structure Recovery via Local Proxy Decomposition \textcolor{gray}{\small (O3/H0/R11, conf: \textbf{low})}}, breakable, before skip=4pt, after skip=4pt]
\textbf{Tactical pattern.} Take a centralized algorithm whose core step requires a specific global object --- pairwise distance matrix, summed gradient, supernet weights, surrogate-model posterior, or class prototypes --- and replace that object with a synthetic reconstruction built from quantities each client can already compute locally. The verbs are concrete: aggregate local gradients to drive global virtual-data distillation (NeurIPS\_2023\_0635), run matrix completion on local distance blocks to approximate the global similarity matrix (NeurIPS\_2025\_1929), perturb a shared supernet locally to obtain per-client architectures (ICLR\_2023\_0352), accumulate compression error feedback inside FedAvg rounds (ICML\_2023\_0413), or share class prototypes instead of raw subgraph data (ICLR\_2025\_1596). The deliverable is always paired: a federated procedure plus a convergence/recovery theorem showing the proxy-based version matches the centralized solution up to bounded error.

\textbf{Differentiation.} \textit{Versus sibling cluster 42 (Decompose-then-Recombine Multi-Objective Search), which decomposes along the OBJECTIVE axis --- splitting a scalar-weighted target into structured sub-objectives whose recombination yields a richer Pareto-style solution space --- this cluster decomposes along the DATA/AGENT axis, where the partition is forced exogenously by privacy or communication constraints rather than chosen for objective expressiveness. Versus sibling cluster 44 (Disentangled Generative Subproblems), which factors a generative model into independent latent or part-level streams for controllability, cluster 13's recombination must reproduce a SPECIFIC global object (a distance matrix, a gradient, a posterior) that the centralized algorithm explicitly references, making the recombination step a structured-recovery problem rather than a compositional design choice. The signature move unique to cluster 13 is naming the global object the centralized method consumes and proving a local-proxy substitute preserves its solution.}

\textbf{When to pick.} Pick this instantiation when your centralized baseline has an explicit dependency on a global statistic (all-pairs comparisons, full-batch gradient, joint posterior, shared supernet) that the partition constraint forbids --- and that statistic admits a low-rank, additive, or otherwise structured factorization computable from local pieces. If the bottleneck is instead a single scalar trade-off or a coupled generative process, sibling clusters 42 or 44 fit better.

\textbf{Tactical failure.} The dominant failure is "mechanical federation port": the paper takes algorithm X, observes that quantity Q is the only globally-coupled term, swaps in an obvious local aggregate of Q, and reproves convergence under the now-standard FL machinery --- yielding a result that reviewers read as a routine extension rather than a structural insight. NeurIPS\_2024\_1160 stacks blockchain + FL + anomaly detection without showing the proxy decomposition unlocks anything beyond a standard FedAvg pipeline; NeurIPS\_2023\_0634 (FAVAS) and ICLR\_2024\_0960 derive partial-participation/asynchronous variants whose theorems track the heterogeneity in the obvious way; ICLR\_2025\_1429 ports MLE-IRL to the federated setting with a routine bilevel convergence proof. A second failure pattern is overlooking that the "non-sensitive" proxy itself leaks: ICML\_2025\_1699 (DictPFL) is explicit that the locally-retained basis can encode private inputs even when only coefficients are transmitted, and SENSE (NeurIPS\_2025\_1929) reconstructs a global distance matrix from local distances without resolving whether that reconstruction itself violates the privacy budget the federation was meant to enforce. Rejects rarely fail at the math --- they fail because the proxy substitution feels like a deployment patch rather than a new algorithmic primitive.

\textbf{Examples.}
\begin{itemize}[leftmargin=12pt,topsep=2pt,itemsep=1pt]
  \item \textbf{[Oral]} \texttt{ICLR\_2025\_1596} \textemdash\ Shifts what is federated from model parameters to class prototypes and uses them to synthesize global-distribution training data --- a textbook instance of recovering global structure (the cross-client distribution) from a compact locally-computable proxy.
  \\ \textcolor{gray}{\footnotesize \textit{Lesson:} Federate class prototypes (not parameters) and use them to synthesize global-distribution training data --- recover global structure (cross-client distribution) from a compact locally-computable proxy.}
  \item \textbf{[Oral]} \texttt{ICML\_2023\_0413} \textemdash\ Reconstructs the centralized error-feedback bias-correction trajectory by aggregating per-client compressed-gradient residues, proving the decomposed FL procedure recovers the full-precision convergence rate the centralized algorithm enjoys.
  \\ \textcolor{gray}{\footnotesize \textit{Lesson:} Aggregate per-client compressed-gradient residues to reconstruct centralized error-feedback bias-correction trajectory --- prove the decomposed FL procedure recovers full-precision rate.}
  \item \textbf{[Oral]} \texttt{ICLR\_2025\_1541} \textemdash\ Removes the need for a global problem-parameter estimate (the very structure prior FL convergence proofs assumed centralized access to) by adapting from observable per-client gradient norms --- replacing a global-knowledge requirement with a locally-computed proxy.
  \\ \textcolor{gray}{\footnotesize \textit{Lesson:} Replace the global problem-parameter estimate with adaptation from observable per-client gradient norms --- convert a global-knowledge requirement into a locally-computed proxy.}
  \item \textbf{[Reject]} \texttt{NeurIPS\_2025\_1929} \textemdash\ Substitutes the global pairwise distance matrix that neighbor-embedding requires with a matrix-completion reconstruction from local distance blocks --- exactly this cluster's move, but the proxy reconstruction reads as a routine application of low-rank completion rather than a new federation primitive.
  \\ \textcolor{gray}{\footnotesize \textit{Lesson:} Matrix-completion reconstruction of a global pairwise distance matrix from local blocks is the cluster's move, but reads as routine low-rank application without a new federation primitive.}
  \item \textbf{[Reject]} \texttt{ICML\_2025\_1699} \textemdash\ Decomposes the shared model into a locally-retained basis plus transmittable coefficients so only the proxy layer is encrypted and aggregated, but the basis itself becomes a side-channel --- illustrating how the proxy substitution can silently violate the privacy constraint that motivated decomposition.
  \\ \textcolor{gray}{\footnotesize \textit{Lesson:} Local basis + transmittable coefficients can silently violate the privacy constraint that motivated decomposition --- the proxy itself becomes a side channel.}
  \item \textbf{[Reject]} \texttt{ICLR\_2023\_0352} \textemdash\ Treats the global supernet as a shared anchor and lets each client locally perturb operations to find a per-client sub-architecture, but the contribution lands as a mechanical pairing of NAS supernet weight-sharing with FL aggregation rather than a new structural insight about either.
  \\ \textcolor{gray}{\footnotesize \textit{Lesson:} Mechanically pairing NAS supernet weight-sharing with FL aggregation --- neither component's structural insight is advanced.}
\end{itemize}
\end{tcolorbox}

\begin{tcolorbox}[colback=bgskill!50, colframe=cgreen!70, title={\textbf{C41}: Decomposed Intermediate State Chaining \textcolor{gray}{\small (O8/H3/R5, conf: \textbf{high})}}, breakable, before skip=4pt, after skip=4pt]
\textbf{Tactical pattern.} Insert an explicit, inspectable intermediate state between each reasoning step so that step k's output becomes a typed, conditioning input to step k+1, and use that intermediate channel as the locus for verification, refinement, or selective correction. Concrete moves include: alternating "select facts" and "derive conclusion" calls (ICLR\_2023\_0337), predicting post-execution program states between synthesis steps (ICLR\_2024\_0826), generating-then-counterexample-refining hypotheses (ICLR\_2024\_0922), maintaining a global state across a reasoning tree with a dedicated rethink pass (NeurIPS\_2024\_1189), gating a slow symbolic corrector by a learned per-location violation predictor (ICLR\_2023\_0272), and growing a library of intermediate lemmas during proving (ICLR\_2024\_0881). The unit being chained is always a semantic intermediate (state, hypothesis, lemma, plan-fragment), not just a token sequence.

\textbf{Differentiation.} \textit{Where Cluster 44 decomposes a problem into structurally parallel, disentangled subproblems (separate latent streams or part-level conditioning that can be solved independently), Cluster 41 decomposes along the temporal axis into a strictly sequential chain whose correctness depends on the verifiability of each link's intermediate output. Cluster 42 decomposes the objective itself (scalarization $\to$ Pareto fronts, multi-criteria oracles), whereas Cluster 41 keeps the objective fixed and decomposes the reasoning path to it. Cluster 7 uses staged decomposition for post-hoc theoretical analysis of an existing dynamical system; Cluster 41 prescribes the staged decomposition as a constructive system architecture. The defining contrast with all three siblings is that Cluster 41's intermediate states are operationally consumed (read, verified, or corrected) by the next stage, not merely combined at the end.}

\textbf{When to pick.} Pick this tactic when the failure mode you observe is end-to-end correctness loss specifically attributable to one missing or unreliable cognitive sub-operation in a chain (selection vs. application, knowledge vs. search, hypothesis generation vs. testing) --- i.e., when you can name a capability asymmetry between sub-steps. If instead the bottleneck is parallel/structural (multiple modalities, multiple objectives, multiple regions), prefer Cluster 42 or 44.

\textbf{Tactical failure.} Rejects in this cluster cluster around chains that re-express the I/O flow as a multi-step pipeline without diagnosing a specific capability asymmetry the chain resolves. CogCoM (NeurIPS\_2024\_1177) mandates intermediate operation-result pairs but the chain is justified by anti-shortcut hygiene rather than by isolating a sub-operation the model provably fails at; reviewers see chain-of-manipulations as architectural decoration over conventional VL training. RePLan (ICLR\_2024\_0936) inserts a perception-grounded replanning step but the closed-loop verification is treated as engineering for plan execution rather than as a learned localization of where reasoning fails. The o1 underthinking paper (ICML\_2025\_1812) operates on the chain at the decoding level with a switching penalty --- a surface-level patch on the chain rather than a redesign of the intermediate state. ICLR\_2024\_0804 inverts the chain (decouple plan from execution) but lacks the ablation showing the coupling itself was the bottleneck. MSR-ViR (ICLR\_2025\_1514) adopts planner+executor decomposition but trains the planner with sparse downstream reward, so the intermediate states never become verifiable artifacts. The pattern: rejects add chain structure as a method ingredient; orals start from a named asymmetry between sub-steps and show the chain is the minimal intervention that closes it.

\textbf{Examples.}
\begin{itemize}[leftmargin=12pt,topsep=2pt,itemsep=1pt]
  \item \textbf{[Oral]} \texttt{ICLR\_2023\_0337} \textemdash\ Cleanest instantiation of the tactic: explicitly alternates Selection and Inference as two distinct sub-operations and shows the chained intermediate facts are what unlocks multi-step reasoning the same model cannot do monolithically.
  \\ \textcolor{gray}{\footnotesize \textit{Lesson:} Alternate Selection and Inference as two distinct sub-operations --- chained intermediate facts unlock multi-step reasoning the same model cannot do monolithically.}
  \item \textbf{[Oral]} \texttt{ICLR\_2024\_0826} \textemdash\ Defines the chained intermediate as an execution state (semantic) rather than a code prefix (syntactic), demonstrating that the choice of what gets passed between links is the load-bearing design decision in this tactic.
  \\ \textcolor{gray}{\footnotesize \textit{Lesson:} Define the chained intermediate as execution state (semantic), not code prefix (syntactic) --- the choice of WHAT gets passed between links is the load-bearing decision.}
  \item \textbf{[Oral]} \texttt{NeurIPS\_2024\_1189} \textemdash\ Extends the linear chain to a tree with a global revision pass, showing the tactic generalizes when you let later intermediate states retroactively rewrite earlier ones.
  \\ \textcolor{gray}{\footnotesize \textit{Lesson:} Extend the linear chain to a tree with global revision pass --- later intermediate states retroactively rewrite earlier ones.}
  \item \textbf{[Oral]} \texttt{ICLR\_2024\_0881} \textemdash\ Treats the chain's intermediate artifacts (helper lemmas) as a persistent, growing resource the system writes back to itself, turning each verified intermediate state into reusable scaffolding for future chains.
  \\ \textcolor{gray}{\footnotesize \textit{Lesson:} Treat verified intermediate artifacts (helper lemmas) as persistent growing scaffolding for future chains --- chains write back to themselves.}
  \item \textbf{[Reject]} \texttt{NeurIPS\_2024\_1177} \textemdash\ Mandates intermediate operation-result pairs but the chain is motivated by anti-shortcut training hygiene rather than by a named capability the chain unlocks, so the intermediate states read as supervision targets, not as the diagnostic locus.
  \\ \textcolor{gray}{\footnotesize \textit{Lesson:} Chains motivated by anti-shortcut training hygiene rather than a named capability the chain unlocks --- intermediates become supervision targets, not the diagnostic locus.}
  \item \textbf{[Reject]} \texttt{ICLR\_2024\_0936} \textemdash\ Adds perceptual replanning between plan steps but frames the intermediate verification as runtime engineering for execution robustness rather than as resolving a sub-step capability asymmetry, leaving the chain looking like a closed-loop control wrapper.
  \\ \textcolor{gray}{\footnotesize \textit{Lesson:} Perceptual replanning as runtime engineering for execution robustness, not resolving a sub-step capability asymmetry --- looks like closed-loop control wrapper.}
  \item \textbf{[Reject]} \texttt{ICML\_2025\_1812} \textemdash\ Acts on the chain only at the decoding layer via a thought-switching token penalty --- surface-level intervention on chain dynamics rather than redesigning what the intermediate state carries.
  \\ \textcolor{gray}{\footnotesize \textit{Lesson:} Decoding-layer thought-switching token penalty is surface intervention on chain dynamics, not redesigning what the intermediate state carries.}
\end{itemize}
\end{tcolorbox}

\begin{tcolorbox}[colback=bgskill!50, colframe=cgreen!70, title={\textbf{C42}: Decompose-then-Recombine Multi-Objective Search \textcolor{gray}{\small (O1/H0/R10, conf: \textbf{low})}}, breakable, before skip=4pt, after skip=4pt]
\textbf{Tactical pattern.} The shared move is to take an optimization problem currently solved as a single scalarized objective with manually-tuned weights (or as a separate iterative solve per instance) and either (a) replace the scalar with an explicit multi-objective search that recovers a Pareto frontier via evolutionary populations (ICLR\_2024\_0902), gradient-based MOO (ICLR\_2025\_1424), or sequential single-criterion subproblems with monotone-shrinkage selection (ICLR\_2024\_0814), or (b) amortize the per-instance optimization itself by training a meta-predictor that maps inputs to near-optimal parameters/initializations (NeurIPS\_2024\_1311, ICLR\_2023\_0280, ICLR\_2023\_0271). A third sub-move uses lexicographic ordering or MTO-derived statistics to set scalar weights principally rather than by hand (ICLR\_2023\_0359, ICLR\_2025\_1356). The verbs are concrete: partition the objective, run a population/oracle/MTO pass per partition, then recombine into a Pareto set, a lexicographic chain, or a learned warm-start. The decomposition axis is always the optimization target itself --- sub-objectives, sub-instances, or sub-regions of the solution space --- not time, modality, or reasoning steps.

\textbf{Differentiation.} \textit{Sibling 7 (timescale separation) decomposes a *dynamical system* across phases of training to enable analysis; cluster 42 instead decomposes a *static optimization target* across competing objectives to enable search. Sibling 41 (intermediate state chaining) decomposes a *reasoning trajectory* into verifiable sub-steps where each output conditions the next; cluster 42's sub-problems are run in parallel or amortized into a learned mapping, never chained as conditioned reasoning. Sibling 44 (disentangled generative subproblems) factors a *generative model's conditioning* into independent latent streams; cluster 42 leaves the model architecture intact and instead restructures the training objective or per-instance solve. Sibling 13 (federated local proxies) decomposes along *data/privacy boundaries*; cluster 42's decomposition is along the objective itself, with no privacy or distribution constraint forcing the split.}

\textbf{When to pick.} Pick this tactic when your problem currently uses a hand-tuned scalar weight (or a per-instance optimizer loop) over multiple competing criteria, AND you have evidence the trade-off frontier is non-convex, the weight choice is brittle, or the per-instance cost is the bottleneck --- not when you need staged analysis (sibling 7), iterative reasoning (sibling 41), or modular generation (sibling 44).

\textbf{Tactical failure.} The dominant failure pattern across rejects in this cluster is that the decomposition is methodologically natural enough to feel "obviously correct," which lets authors skip the ablation that isolates *which* component of the decomposition actually drives the gain. ICLR\_2025\_1424 reformulates IB as MOO but never separates frontier non-convexity from model capacity or data effects, so reviewers cannot tell if MOO itself is doing the work. ICLR\_2024\_0814 proves convergence of a sequence-of-subproblems scheme but does not ablate the subproblem-selection rule against the underlying single-criterion oracle, leaving open whether the selection ordering or the solver is binding. ICLR\_2025\_1356 (AutoScale) and ICLR\_2024\_0902 (MOESR) similarly bolt a multi-objective layer onto an existing pipeline without isolating it. The amortization variants (NeurIPS\_2024\_1311, ICLR\_2023\_0280, ICLR\_2023\_0271) fail along an adjacent axis: they assume the input-to-optimum mapping is smooth across instances but do not stress-test out-of-distribution targets, so reviewers read the speedup as overfit-to-training-distribution rather than genuine amortization. ICLR\_2023\_0234 and ICLR\_2025\_1605 add elaborate decomposition machinery (ranking-preserving surrogates, semantic facet partitioning) whose value collapses if the underlying assumption --- heterogeneous sub-structure, ordinal-only fidelity --- is not separately validated.

\textbf{Examples.}
\begin{itemize}[leftmargin=12pt,topsep=2pt,itemsep=1pt]
  \item \textbf{[Oral]} \texttt{ICLR\_2023\_0359} \textemdash\ Cleanly imports the classical lexicographic ordering structure to replace scalar weighting with a sequentially-constrained decomposition of objectives --- the tactical move (decompose the objective, recombine via priority chain) is explicit and the choice is principled rather than empirical.
  \\ \textcolor{gray}{\footnotesize \textit{Lesson:} Import classical lexicographic ordering to replace scalar weighting with sequentially-constrained decomposition --- explicit and principled, not empirical.}
  \item \textbf{[Oral]} \texttt{ICLR\_2023\_0260} \textemdash\ Strongest non-Oral demonstration in the cluster: builds a hierarchical tree decomposition of the solution space to seed an iterative optimizer, exemplifying the decompose-then-recombine pattern with a structural proof rather than just a pipeline (included to satisfy minItems given only 1 Oral exists in this cluster).
  \\ \textcolor{gray}{\footnotesize \textit{Lesson:} Hierarchical tree decomposition of the solution space seeded an iterative optimizer --- decompose-then-recombine with a structural proof, not a pipeline.}
  \item \textbf{[Reject]} \texttt{ICLR\_2025\_1424} \textemdash\ Replaces the IB scalar weight with a gradient-based multi-objective solver to recover a non-convex Pareto front, but provides no ablation isolating frontier non-convexity from confounding factors --- exactly the methodologically-natural-but-unisolated failure pattern.
  \\ \textcolor{gray}{\footnotesize \textit{Lesson:} IB scalar weight $\to$ gradient-based multi-objective solver needs ablation isolating frontier non-convexity from confounding factors --- methodologically natural is not validation.}
  \item \textbf{[Reject]} \texttt{ICLR\_2024\_0814} \textemdash\ Decomposes Pareto-front recovery into a sequence of single-criterion subproblems with a monotone-shrinkage selector, but never ablates the selection rule independently of the oracle, so reviewers cannot attribute convergence to the proposed decomposition.
  \\ \textcolor{gray}{\footnotesize \textit{Lesson:} Pareto-front decomposed into single-criterion subproblems with monotone-shrinkage selector needs to ablate the selection rule independently of the oracle.}
  \item \textbf{[Reject]} \texttt{ICLR\_2025\_1356} \textemdash\ Two-phase decomposition that uses MTO exploration to set linear-scalarization weights --- a clean recombination move, but the calibration step is not validated against a wide enough range of weight regimes to prove the hybrid actually outperforms either pure approach.
  \\ \textcolor{gray}{\footnotesize \textit{Lesson:} Two-phase MTO + linear-scalarization weights needs validation against a wide enough range of weight regimes to prove the hybrid outperforms either pure approach.}
\end{itemize}
\end{tcolorbox}

\begin{tcolorbox}[colback=bgskill!50, colframe=cgreen!70, title={\textbf{C44}: Structured Decomposition into Disentangled Generative Subproblems \textcolor{gray}{\small (O8/H2/R6, conf: \textbf{high})}}, breakable, before skip=4pt, after skip=4pt]
\textbf{Tactical pattern.} Identify a compositional axis WITHIN the generative target itself (identity vs motion vs expression for talking faces; head/body/hand streams for egocentric motion; content vs phase/transition for long sequences; prosody vs timbre vs content for speech; texture vs illumination for 3D), then instantiate that decomposition as ARCHITECTURAL channels in the generator: separate latent codebooks, dedicated conditioning streams routed to different attention layers, distinct sub-diffusion models per factor, or staged generators where each stage owns one factor. Recombination is built into the model --- not done post-hoc --- typically by feeding the disentangled factors back into a shared denoiser/decoder under a disentanglement-promoting training objective (compositional distillation, factorized codec losses, phase-space alignment). The signature verbs are "factorize the codec into N streams", "inject part-level conditioning separately", "train sub-diffusion-X conditioned on factor-X's latent", and "anchor stage-2 generation on stage-1's boundary state".

\textbf{Differentiation.} \textit{Versus sibling cluster 41 (Decomposed Intermediate State Chaining), the decomposition here is across SIMULTANEOUS factors of the OUTPUT signal (identity $\perp$ motion $\perp$ expression coexisting in every frame), whereas cluster 41 decomposes a REASONING TRAJECTORY into sequential verifiable sub-steps where step k's output feeds step k+1's input --- temporal chaining of computations, not factorization of a generative signal. Versus cluster 42 (Decompose-then-Recombine Multi-Objective Search), this cluster decomposes the GENERATIVE OUTPUT structure to gain controllability and quality, while cluster 42 decomposes an OPTIMIZATION OBJECTIVE (Pareto front, scalarization weights) to explore solution diversity --- there is no generative model to "recompose" in cluster 42, just aggregated weights or initializations. Versus cluster 7 (Timescale Separation), that cluster decomposes for THEORETICAL ANALYSIS of dynamics, not for engineering controllable synthesis.}

\textbf{When to pick.} Pick this tactic when your generative target admits a known semantic factorization (identity/motion, content/style, geometry/appearance, foreground/background) AND prior end-to-end generators visibly conflate these factors --- symptoms include identity drift under prompt changes, motion artifacts when conditioning is weak, or inability to control one factor without disturbing another. If instead your problem is sequential reasoning or scalar optimization, pick cluster 41 or 42.

\textbf{Tactical failure.} The recurring failure for this tactic in rejected papers is that the proposed factorization is either INCOMPLETE (one factor still leaks into another) or NOT VALIDATED as the right axis of decomposition for the task. NeurIPS\_2023\_0666 (ITEM3D) disentangles texture from illumination via an SDS directional prior but reviewers were unconvinced the prior actually breaks the ambiguity rather than just shifting it. NeurIPS\_2024\_1153 (AdaFace) introduces a compositional distillation loss to balance identity vs text --- but the loss is a regularizer, not a true architectural factorization, so reviewers saw it as a heuristic balance trick rather than disentanglement. ICLR\_2024\_0977 (Conceptor autoencoder) borrows conceptor regularization as the disentanglement mechanism but applies it to a toy two-pattern interpolation setting, which reviewers read as a domain-specific gadget without evidence the latent is genuinely convex/disentangled at scale. ICLR\_2023\_0395 (VQR) borrows a localization mask from spatial detection and grafts it onto cross-attention as the "decomposition" mechanism --- reviewers saw this as a one-off architectural patch rather than a principled factorization of the generative problem. NeurIPS\_2024\_1174 (CI-GAN) injects a character-structure embedding into a conditional GAN but never demonstrates the structural conditioning produces independent control over the factor it claims to disentangle. Across all rejects, the missing element is direct evidence --- through factor-swap experiments, controllability metrics on each axis, or recombination ablations --- that the proposed channels are actually disentangled and that joint quality is preserved.

\textbf{Examples.}
\begin{itemize}[leftmargin=12pt,topsep=2pt,itemsep=1pt]
  \item \textbf{[Oral]} \texttt{NeurIPS\_2025\_1915} \textemdash\ Explicitly instantiates part-level conditioning streams (head, body, hand) as separate injection modules into the generator --- the canonical 'factor $\to$ architectural channel' move of this cluster.
  \\ \textcolor{gray}{\footnotesize \textit{Lesson:} Part-level conditioning streams (head, body, hand) as separate injection modules --- the canonical 'factor $\to$ architectural channel' move.}
  \item \textbf{[Oral]} \texttt{ICML\_2024\_1072} \textemdash\ Factorizes the speech codec itself into disentangled streams (prosody, content, acoustics) and trains a dedicated diffusion model per factor --- the most architecturally complete instantiation of factor-then-recombine generation.
  \\ \textcolor{gray}{\footnotesize \textit{Lesson:} Factorize the speech codec itself into prosody / content / acoustics --- train a dedicated diffusion model per factor, the most architecturally complete instantiation.}
  \item \textbf{[Oral]} \texttt{NeurIPS\_2024\_1328} \textemdash\ Constructs a disentangled latent space (identity vs dynamic motion) and runs the generative process inside it, demonstrating that the right place for decomposition is the latent representation, not the output space.
  \\ \textcolor{gray}{\footnotesize \textit{Lesson:} Disentangled latent space (identity vs dynamic motion) --- the right place for decomposition is the latent representation, not the output space.}
  \item \textbf{[Oral]} \texttt{NeurIPS\_2025\_1854} \textemdash\ Splits long-motion generation into a content sub-diffusion (SPDM) and a transition sub-diffusion (TPDM) operating in a shared phase space --- staged-generator instance of the tactic with explicit boundary-state anchoring for recombination.
  \\ \textcolor{gray}{\footnotesize \textit{Lesson:} Split long-motion generation into content sub-diffusion (SPDM) + transition sub-diffusion (TPDM) in shared phase space --- explicit boundary-state anchoring for recombination.}
  \item \textbf{[Reject]} \texttt{NeurIPS\_2024\_1153} \textemdash\ The 'decomposition' is a balancing loss between identity and text rather than an architectural factorization, leaving reviewers unconvinced the two factors are genuinely separated as opposed to traded off.
  \\ \textcolor{gray}{\footnotesize \textit{Lesson:} A balancing loss between identity and text is not architectural factorization --- reviewers stay unconvinced the factors are genuinely separated.}
  \item \textbf{[Reject]} \texttt{NeurIPS\_2023\_0666} \textemdash\ Asserts a texture/illumination disentanglement guided by a generative prior but does not demonstrate that the directional SDS signal actually resolves the ambiguity rather than absorbing it into the prior.
  \\ \textcolor{gray}{\footnotesize \textit{Lesson:} Asserting texture/illumination disentanglement guided by a generative prior without showing the directional SDS signal resolves rather than absorbs the ambiguity.}
  \item \textbf{[Reject]} \texttt{ICLR\_2023\_0395} \textemdash\ Borrows a localization mask from object detection and grafts it into cross-attention as the decomposition mechanism --- a one-off architectural patch that reviewers read as a heuristic rather than a principled factorization of generation.
  \\ \textcolor{gray}{\footnotesize \textit{Lesson:} Borrowing a localization mask from object detection and grafting it into cross-attention is a one-off architectural patch, not a principled factorization.}
\end{itemize}
\end{tcolorbox}

\subsection{Diagnose Then Surgically Fix}
\textit{Parent methodology id: \texttt{root\_cause\_surgical\_remediation}; 2 tactical clusters.}

\begin{tcolorbox}[colback=bgskill!50, colframe=cgreen!70, title={\textbf{C02}: Root Cause Diagnosis with Surgical Remediation \textcolor{gray}{\small (O11/H3/R3, conf: \textbf{low})}}, breakable, before skip=4pt, after skip=4pt]
\textbf{Tactical pattern.} Papers in this cluster localize the failure to a SCALAR-LEVEL operation inside the training pipeline --- a single loss term, a gradient estimator's expectation, a numerical quantity at one timestep, a per-component update rule --- and then patch exactly that operation with one of three concrete edits: (a) swap the loss/penalty term while keeping the rest of the objective intact (Pseudo-Huber for MSE in ICLR\_2024\_0855, max-logit for ground-truth-logit in NeurIPS\_2025\_1905, energy-weighted phase loss for L2 in ICLR\_2021\_0046); (b) insert a debiasing/rescaling factor that restores the unbiasedness of an estimator that the original pipeline silently broke (per-batch reweighting in ICLR\_2024\_0860, per-objective gradient correction before aggregation in ICLR\_2023\_0282, non-parametric pseudo-count supplement in ICML\_2025\_1693); or (c) carve out the pathological numerical region and apply a region-specific intervention (timestep truncation around the Lipschitz singularity in ICLR\_2024\_0884, output scaling on identified unstable derivatives in ICLR\_2025\_1583, gradient-norm filtering plus SAM in ICLR\_2023\_0376). The signature is that the rest of the network, dataset, and training loop are byte-for-byte unchanged --- the diff is a few lines around one mathematical object.

\textbf{Differentiation.} \textit{The contrast with sibling cluster 31 (Bottleneck Audit with Targeted Component Substitution) is the granularity of the edit. Cluster 31 papers swap WHOLE MODULES --- a CLIP visual encoder for an ImageNet backbone (ICLR\_2022\_0099), a stronger pre-trained encoder, an entire counterfactual-prediction submodule (ICLR\_2024\_0944) --- and the diagnostic framing is "this slot in the pipeline is undersized/mismatched, drop in a better-formed slot." Cluster 2 papers, in contrast, do not introduce or swap any reusable architectural component; they edit a mathematical expression (a loss term, an estimator, a numerical bound) inside a slot that remains structurally unchanged. Concretely: ICLR\_2024\_0860 keeps the entire training pipeline and merely multiplies retained-example losses by a scalar; ICLR\_2024\_0884 keeps the diffusion architecture and merely truncates a timestep range. A reader cannot reuse a cluster-2 fix as a drop-in module --- they must reproduce the diagnostic chain to know where to apply the scalar-level patch.}

\textbf{When to pick.} Pick this cluster when the failure mode survives every architectural and data-level fix you can think of, AND you can write down the offending mathematical object on one line (a specific loss term, an estimator's expectation, a quantity that diverges at a known regime). If your candidate fix is "replace module M with module M'," go to cluster 31 instead.

\textbf{Tactical failure.} The failure mode for this tactic is delivering only HALF of the diagnose-then-fix arc. ICLR\_2024\_0864 (PINNs near finite-time blow-up) derives a generalization bound that ties solver error to solution stiffness near singularities --- a textbook diagnosis --- but stops there and provides no surgical fix that demonstrably restores accuracy, leaving reviewers with analysis-without-remediation. ICLR\_2023\_0225 (ELBO-ing Stein Mixtures) executes the diagnostic correctly (the prior method implicitly fixes a degenerate boundary value of a parameterized family) but its "remedy" is to expose the parameter as a sweep dial and pick the empirically-best index --- this reads as hyperparameter tuning rather than the principled minimal edit that orals deliver, because it does not commit to a single corrected operation. ICML\_2025\_1794 (Semi-gradient DICE) correctly identifies that semi-gradient updates compute the wrong fixed point, but the proposed redesign is not validated as the minimal change that resolves the issue, blurring the line between "surgical fix" and "method redesign." Common thread across these rejects: the diagnosis is sharp, but the intervention is either absent (0864), parameterized rather than committed (0225), or non-minimal (1794) --- failing the verify-F-resolved-while-unrelated-behavior-preserved leg of the parent operational signature.

\textbf{Examples.}
\begin{itemize}[leftmargin=12pt,topsep=2pt,itemsep=1pt]
  \item \textbf{[Oral]} \texttt{ICLR\_2024\_0884} \textemdash\ Archetypal scalar-level fix: prove the score function has an unbounded Lipschitz constant exactly at t=0, then surgically truncate that timestep region rather than redesigning the diffusion model --- the patch is a region carve-out, not a module swap.
  \\ \textcolor{gray}{\footnotesize \textit{Lesson:} Prove a Lipschitz singularity at t=0 as a theorem, then truncate THAT region only --- a region carve-out beats a module swap.}
  \item \textbf{[Oral]} \texttt{NeurIPS\_2025\_1905} \textemdash\ Decomposes the label-smoothing loss analytically, isolates the ground-truth-logit penalty term as the cause of representation collapse, and replaces only that term with max-logit penalization --- one loss-term swap, everything else byte-identical.
  \\ \textcolor{gray}{\footnotesize \textit{Lesson:} Decompose the loss analytically, isolate the offending term, replace ONLY that term --- preservation by construction is the surgical-fix gold standard.}
  \item \textbf{[Oral]} \texttt{ICLR\_2024\_0860} \textemdash\ Diagnoses dynamic-pruning degradation as gradient-estimator BIAS rather than information loss, then patches the estimator with a per-batch rescaling scalar that proves unbiasedness --- the canonical 'add a scalar correction at one line' move.
  \\ \textcolor{gray}{\footnotesize \textit{Lesson:} Reframe failure as gradient-estimator BIAS (correctable in expectation), not information loss (architectural redesign needed) --- the diagnostic frame determines the fix's locality.}
  \item \textbf{[Oral]} \texttt{ICLR\_2021\_0046} \textemdash\ Decomposes the L2 audio loss into energy and phase components, identifies that uniform L2 systematically under-penalizes phase errors at high-energy regions, and reweights only the phase term --- a within-loss rebalancing rather than an architectural change.
  \\ \textcolor{gray}{\footnotesize \textit{Lesson:} Decomposing a single L2 loss into energy and phase components and rebalancing only one component is within-loss surgery --- no architecture change, no extra parameters.}
  \item \textbf{[Reject]} \texttt{ICLR\_2024\_0864} \textemdash\ Delivers the diagnosis half (generalization bounds tied to solution stiffness near blow-up) but no surgical remediation that resolves the failure --- analysis without a committed minimal edit.
  \\ \textcolor{gray}{\footnotesize \textit{Lesson:} Diagnosis without committed minimal edit fails --- a generalization bound that explains the pathology must come paired with a one-line fix the bound justifies.}
  \item \textbf{[Reject]} \texttt{ICLR\_2023\_0225} \textemdash\ Diagnoses that the established method picks a degenerate boundary value of a parameterized family, but the 'fix' is exposing the parameter as a sweep dial rather than a principled single-point correction --- fails the minimal-edit criterion.
  \\ \textcolor{gray}{\footnotesize \textit{Lesson:} Exposing a parameter as a sweep dial is not a fix --- the methodology demands a principled single-point correction, not 'choose your own value'.}
  \item \textbf{[Reject]} \texttt{ICML\_2025\_1794} \textemdash\ Identifies that semi-gradient DICE silently computes the wrong fixed point, but the proposed full-gradient redesign is not framed or validated as the minimal change isolating that pathology, leaving the surgical-fix bar unmet.
  \\ \textcolor{gray}{\footnotesize \textit{Lesson:} Identifying the wrong fixed point (semi-gradient $\ne$ full-gradient DICE) is correct diagnosis, but the proposed redesign must be framed AS the minimal change isolating that pathology.}
\end{itemize}
\end{tcolorbox}

\begin{tcolorbox}[colback=bgskill!50, colframe=cgreen!70, title={\textbf{C31}: Bottleneck Audit with Targeted Component Substitution \textcolor{gray}{\small (O17/H44/R24, conf: \textbf{high})}}, breakable, before skip=4pt, after skip=4pt]
\textbf{Tactical pattern.} Take a working multi-component pipeline, audit each module against the structural properties of the target setting (or against what an adjacent community has since improved), localize a single component whose implicit assumptions mismatch the actual problem, and SWAP that component out --- replacing a tokenizer/codebook/aggregation rule/auxiliary mechanism/encoder/training stage with a structure-matched alternative while leaving the rest of the pipeline untouched. The verbs are concrete and architectural: 'substitute the featurizer with CLIP' (ICLR\_2022\_0099), 'replace the global softmax aggregation with cheaper structure-aware ops' (ICLR\_2025\_1573), 'transplant the diffusion framework wholesale to escape multiple incumbent constraints simultaneously' (ICLR\_2021\_0016), 'remove the noisy duration-predictor proxy and inject ground-truth pitch/energy directly' (ICLR\_2021\_0021), 'reformulate the inherited noise-schedule sampling to match the new flow geometry' (ICML\_2024\_1121). The deliverable is a drop-in replacement at a named module boundary, justified by an empirical gap in that module's expressive capacity or assumption fit.

\textbf{Differentiation.} \textit{Sibling cluster 2 (Root Cause Diagnosis with Surgical Remediation) localizes failures to MATHEMATICAL/NUMERICAL pathologies inside losses or optimization dynamics --- Lipschitz singularities in the diffusion noise schedule, EMA-driven gradient instability in consistency training, attribution of variance to specific numerical sources --- and the surgical fix is a small loss-term, schedule, or numerical correction. Cluster 31 instead localizes failure at a MODULE BOUNDARY in a pipeline (a tokenizer, an encoder, an aggregation operator, a discrete codebook, an auxiliary stage, an inherited noise sampler) and the fix is a wholesale component swap --- often importing a module from an adjacent community (DiffWave transplanting diffusion, ICLR\_2024\_0870 redesigning the tokenizer interface for visual generation, ICML\_2023\_0420 swapping in an upgraded diffusion generator inside an adversarial-training pipeline). Put differently: sibling 2 edits formulas, cluster 31 unplugs and replaces blocks.}

\textbf{When to pick.} Pick this tactic when the failure manifests as 'one named module is the binding constraint' inside an otherwise-working pipeline --- typically when an adjacent community has produced a strictly better version of that module, or when the module's implicit assumption (independence, fixed schedule, hard assignment, single global policy, homogeneous input) is provably violated by your target distribution. Choose sibling 2 instead if the failure is a numerical or loss-landscape pathology with no clean module boundary to swap against.

\textbf{Tactical failure.} Reject papers in this cluster perform the substitution and motivate it with a plausible structural-mismatch story plus theoretical grounding, but fail to demonstrate that the substituted-out component was actually the DOMINANT bottleneck rather than one of several secondary factors. Reviewers see an extension that 'looks like a natural improvement' and demand the controlled ablation that would isolate the swapped module's contribution from data, capacity, and training-recipe confounds. ICLR\_2025\_1573 (replacing attention with structure-aware MPNN ops) and ICLR\_2025\_1415 (substituting pairwise structure-aware encoders) both explicitly note the gap is invisible on standard benchmarks because incumbents are 'already competitive' --- the proposed swap is necessary only if directly stress-measured. ICLR\_2025\_1564 (soft input-conditioned matching replacing VQ codebook lookup) and ICLR\_2025\_1369 (CardiCat structure-aware VAE component) suffer the same: substitution is well-motivated but the targeted mismatch is not shown to be load-bearing versus capacity. ICLR\_2024\_0931 (unified high-order GNN abstraction) is rejected for failing to show the unified interface is better than case-by-case wrappers --- the substitution exists but lacks a head-to-head against the alternative. The pattern: the substitute lands a clean theoretical story, but reviewers cannot tell whether removing the original component would have hurt at all.

\textbf{Examples.}
\begin{itemize}[leftmargin=12pt,topsep=2pt,itemsep=1pt]
  \item \textbf{[Oral]} \texttt{ICLR\_2021\_0016} \textemdash\ DiffWave is the cleanest 'transplant a whole framework as a drop-in module' instance --- auditing constraint-laden incumbents and substituting a diffusion process wholesale to simultaneously dissolve multiple architectural failure modes.
  \\ \textcolor{gray}{\footnotesize \textit{Lesson:} Substitute a diffusion process wholesale to dissolve multiple architectural failure modes simultaneously --- transplant a whole framework as a drop-in module.}
  \item \textbf{[Oral]} \texttt{ICML\_2023\_0465} \textemdash\ Hiera embodies 'swap the learning objective to render compensatory components redundant', then strips the now-unnecessary modules --- a substitution at one interface that cascades into pipeline simplification.
  \\ \textcolor{gray}{\footnotesize \textit{Lesson:} Swap the learning objective to render compensatory components redundant, then strip them --- substitution at one interface cascades into pipeline simplification.}
  \item \textbf{[Oral]} \texttt{ICML\_2024\_1121} \textemdash\ SD3 audits which components were inherited uncritically from the displaced paradigm and replaces only the geometry-mismatched ones (noise sampler, timestep weighting) while keeping the rest of the rectified-flow pipeline intact.
  \\ \textcolor{gray}{\footnotesize \textit{Lesson:} Audit which components were inherited uncritically and replace only the geometry-mismatched ones --- keep the rest of the pipeline intact.}
  \item \textbf{[Oral]} \texttt{ICLR\_2021\_0021} \textemdash\ FastSpeech 2 removes a complex auxiliary teacher module entirely by substituting it with directly-extracted ground-truth conditioning signals --- a textbook 'the bottleneck component was a noisy proxy, swap it out' move.
  \\ \textcolor{gray}{\footnotesize \textit{Lesson:} Remove a complex auxiliary teacher entirely by substituting directly-extracted ground-truth conditioning --- the bottleneck was a noisy proxy, swap it out.}
  \item \textbf{[Reject]} \texttt{ICLR\_2025\_1573} \textemdash\ Substitutes attention with simpler structure-aware operators in large-graph MPNNs but cannot demonstrate the swapped attention block was the binding constraint rather than a secondary factor on standard benchmarks.
  \\ \textcolor{gray}{\footnotesize \textit{Lesson:} Substituting attention with structure-aware operators in MPNNs needs to demonstrate the swapped block was the binding constraint, not a secondary factor.}
  \item \textbf{[Reject]} \texttt{ICLR\_2025\_1564} \textemdash\ Replaces a hard VQ codebook-lookup module with continuous input-conditioned soft matching --- a clear component swap motivated by codebook-mismatch under degradation, but without isolating ablations to prove the original module was the dominant bottleneck.
  \\ \textcolor{gray}{\footnotesize \textit{Lesson:} Hard VQ codebook $\to$ continuous input-conditioned soft matching needs ablations isolating that the original module was the dominant bottleneck.}
  \item \textbf{[Reject]} \texttt{ICLR\_2024\_0931} \textemdash\ Proposes a unified abstraction layer to substitute for ad-hoc per-method high-order GNN code, but lacks the head-to-head comparison showing the unified interface beats simple wrapping --- the substitution is presented without evidence the prior component was actually limiting.
  \\ \textcolor{gray}{\footnotesize \textit{Lesson:} Unified abstraction layer needs head-to-head evidence the prior component was actually limiting --- substitution claims demand binding-constraint proof.}
\end{itemize}
\end{tcolorbox}

\subsection{Soft Probabilistic Relaxation}
\textit{Parent methodology id: \texttt{soft\_probabilistic\_relaxation}; 1 tactical cluster.}

\begin{tcolorbox}[colback=bgskill!50, colframe=cgreen!70, title={\textbf{C04}: Soft Uncertainty Preservation via Probabilistic Relaxation \textcolor{gray}{\small (O3/H0/R12, conf: \textbf{low})}}, breakable, before skip=4pt, after skip=4pt]
\textbf{Tactical pattern.} Locate a specific point in an existing pipeline where the model commits prematurely to a discrete value (argmax pseudolabel, hard constraint satisfaction, fixed Gaussian prior, deterministic latent component, complete-coordinate score-matching, fixed-modality posterior) and replace that single deterministic operation with a distributional object --- a partial-label set over top-k candidates (ICML\_2024\_1008), a log-quadratic relaxation of the Potts energy with soft targets (NeurIPS\_2024\_1280), a latent-variable mixture with EM over annotators (ICLR\_2025\_1540), an importance-weighted score-matching loss over only-observed coordinates (ICML\_2025\_1792), or a learned/score-based prior swapped for the Gaussian (ICLR\_2025\_1451, ICLR\_2025\_1518). Then optimize the relaxed objective via EM, alternating cooperative loops, or differentiable regularization rather than recovering a hard solution. The shared verbs are "replace point with distribution," "swap argmax with candidate set," "lift deterministic component to probabilistic component," and "alternate inference and parameter updates."

\textbf{Differentiation.} \textit{Singleton under parent methodology --- no sibling clusters exist.}

\textbf{When to pick.} Pick this tactic when the bottleneck is a single load-bearing hard commitment (a thresholded label, a delta-distribution prior, an enforced constraint) whose committed value is demonstrably wrong on a meaningful fraction of inputs, and the surrounding training loop already supports iterative inference (so EM/alternation is a drop-in). Avoid it when the hard step is auxiliary or already well-calibrated --- softening it then reads as cosmetic.

\textbf{Tactical failure.} The dominant failure across this cluster's rejects is that the softening reads as a "natural probabilistic upgrade" of an existing model rather than a tactical move that unlocks new capability, and the paper does not isolate the hard commitment as the actual bottleneck. ICLR\_2025\_1451 swaps the VAE's Gaussian prior for a score-based prior as a drop-in change, which reviewers perceive as combining two known components without showing the Gaussian was the binding constraint. ICLR\_2024\_0853 unifies noisy/partial/multi-candidate label settings as MLE-over-latent-labels with EM, but reviewers see integration of prior algorithms rather than a new tactical advance. ICLR\_2023\_0398's soft distributional coherency regularizer reads as a routine relaxation of point-level constraints. ICLR\_2024\_0911 (probabilistic NMF via EM), NeurIPS\_2023\_0674 (Wasserstein autoencoder adapted to DAGs), ICLR\_2023\_0240 (power-set posterior over modalities), ICLR\_2023\_0340 (cooperative generator--discriminator loop), and ICLR\_2025\_1531 (alternating supervised/unsupervised BNN phases) all share this pattern: the relaxation is structurally clean but the experiments don't demonstrate that hard commitment was load-bearing, so the contribution looks like "make X probabilistic" rather than "solve a problem the discrete version could not."

\textbf{Examples.}
\begin{itemize}[leftmargin=12pt,topsep=2pt,itemsep=1pt]
  \item \textbf{[Oral]} \texttt{ICML\_2024\_1008} \textemdash\ Imports the partial-label training framework wholesale to represent CLIP's top-k pseudolabel uncertainty as a candidate set, rather than merely softening argmax --- the move repurposes an external tool's theory rather than locally relaxing a hard step.
  \\ \textcolor{gray}{\footnotesize \textit{Lesson:} Recognizing CLIP top-k confidence $\cong$ partial-label set lets you transplant the partial-label framework wholesale --- repurpose external theory rather than design a new local relaxation.}
  \item \textbf{[Oral]} \texttt{ICLR\_2025\_1540} \textemdash\ Casts annotator assignment as a latent variable in a probabilistic mixture and uses EM both to learn from partial annotations and to control workload via mixture priors, exemplifying the relaxation-plus-iterative-inference recipe with a second payoff.
  \\ \textcolor{gray}{\footnotesize \textit{Lesson:} EM with mixture priors gets workload distribution control as a free byproduct --- the cleanest soft relaxations yield structural bonuses, not just stability.}
  \item \textbf{[Oral]} \texttt{ICML\_2025\_1792} \textemdash\ Directly rewrites the score-matching objective to use only observed coordinates with importance weighting, demonstrating that softening the complete-data assumption is a mathematical rederivation rather than a preprocessing patch.
  \\ \textcolor{gray}{\footnotesize \textit{Lesson:} Rewrite the objective to use only observed coordinates with importance weighting --- softening at the math level, not as preprocessing, makes the relaxation principled.}
  \item \textbf{[Reject]} \texttt{ICLR\_2025\_1451} \textemdash\ Replacing the VAE's Gaussian prior with a score-based prior is a clean instance of the tactic but reads as a drop-in component swap without isolating the Gaussian as the actual bottleneck.
  \\ \textcolor{gray}{\footnotesize \textit{Lesson:} VAE Gaussian prior $\to$ score-based prior is a drop-in component swap unless you isolate Gaussian as the actual bottleneck.}
  \item \textbf{[Reject]} \texttt{ICLR\_2024\_0853} \textemdash\ Unifies noisy/partial/multi-candidate label settings under MLE-over-latent-labels with EM, but the soft-relaxation move is perceived as integrative bookkeeping rather than a new tactical capability.
  \\ \textcolor{gray}{\footnotesize \textit{Lesson:} Unifying noisy/partial/multi-candidate labels under one MLE template reads as integrative bookkeeping unless the unification produces a new tactical capability.}
  \item \textbf{[Reject]} \texttt{ICLR\_2023\_0398} \textemdash\ Soft distributional coherency regularization replaces hard point-level constraints in hierarchical forecasting, but the relaxation reads as the obvious next step rather than a load-bearing tactical move.
  \\ \textcolor{gray}{\footnotesize \textit{Lesson:} Soft distributional coherency replacing hard point constraints is the obvious next step --- the relaxation must be load-bearing, not aesthetic.}
\end{itemize}
\end{tcolorbox}

\subsection{Exploit Redundancy via Compression}
\textit{Parent methodology id: \texttt{redundancy\_compression\_and\_distillation}; 2 tactical clusters.}

\begin{tcolorbox}[colback=bgskill!50, colframe=cgreen!70, title={\textbf{C17}: Redundancy Elimination via Compact Reparameterization \textcolor{gray}{\small (O24/H11/R23, conf: \textbf{high})}}, breakable, before skip=4pt, after skip=4pt]
\textbf{Tactical pattern.} Take a specific structural artifact of an already-trained system --- a weight matrix, an intermediate cache, an activation pattern, a per-layer numerical format, a stored optimizer state --- and replace it with a tailored compact surrogate (low-bit quantization, vector-cluster codebook, low-rank/sparse factor, structured matrix class, frequency-domain reparameterization, distilled residual, null-space projection) that drops in at the same I/O interface. The unifying verbs are "quantize/decompose/reparameterize/prune the static representation," not "shorten the trajectory." Concretely, papers either (a) profile a heterogeneity (BiLLM's long-tail second-order sensitivity, SqueezeLLM's dense-and-sparse split, F8Net's per-layer format) and assign a stratum-specific encoder, or (b) prove an algebraic equivalence (Monarch matrices for both mixers, RMSNorm absorbing centering into weights, rotation-trick gradients for VQ, Fourier dual of integral operators) that lets a cheaper primitive substitute exactly. The compression target is always a frozen structural object, never a multi-step inference loop.

\textbf{Differentiation.} \textit{Versus sibling cluster 18 (Iterative Process Compression via Step Distillation), this cluster compresses the *static system* --- weights, caches, activations, formats --- while sibling 18 compresses the *temporal procedure* (collapsing N denoising/MCMC/optimization steps into K via student-teacher trajectory matching). A QLoRA-style move freezes the base model and quantizes its weight tensor in place; a Progressive Distillation-style move keeps the network identical but trains a student to skip iterations. Sibling 18 papers retain full per-call cost and reduce call count; cluster 17 papers retain call count and reduce per-call cost. The diagnostic: if the redundancy lives inside one forward pass (low-rank weight, ternary cache, redundant neurons, sortable activations), the move belongs here; if it lives across many passes of an autoregressive or iterative process, it belongs to 18.}

\textbf{When to pick.} Pick this tactic when the binding cost is per-call resource (memory footprint of weights, KV-cache RAM, FLOPs of a normalization, storage of checkpoints, parameter count of an adapter) rather than the number of forward passes. Strongest signal: you can point to a concrete tensor or sub-module whose statistics show heterogeneity (long-tail sensitivity, low effective rank, redundant activation overlap, format-mismatched layers) that a uniform default ignores.

\textbf{Tactical failure.} Rejected papers in this cluster fail in three recurring ways. First, they propose a new compact primitive but neglect to show it is the *binding* compression knob --- Transformer Compression via Subspace Projection (NeurIPS\_2023\_0750) and Large-Scale Spectral GNN via Laplacian Sparsification (ICLR\_2024\_0876) port a clean theoretical object into the pipeline without ablating whether the audited bottleneck dominates. Second, they extend an existing scheme to a more aggressive regime without demonstrating the regime is reachable safely --- KVTQ (ICLR\_2024\_0868) pushes KV cache to ternary, ClusComp (ICLR\_2025\_1374) and decoupleQ (NeurIPS\_2024\_1190) push to 1-2 bit, and Nearly Lossless Adaptive Bit Switching (NeurIPS\_2024\_1264) and MixQuant (ICLR\_2023\_0283) layer wrappers on top of base quantizers, but reviewers cannot tell whether gains come from the new mechanism or from leftover headroom in baselines. Third, they bundle two independently-validated compression axes (XYZ Data Efficiency NeurIPS\_2023\_0761, XoRA ICLR\_2025\_1652 combining LoRA with expander sparsity, NeCGS NeurIPS\_2024\_1265 combining INRs with deformable TSDFs) and inherit each component's accepted prior without isolating the joint benefit. The shared symptom: the compact reparameterization is plausible but the paper has not framed exactly which baseline assumption it overturns, so the contribution reads as engineering combinatorics rather than a pinpointed structural insight.

\textbf{Examples.}
\begin{itemize}[leftmargin=12pt,topsep=2pt,itemsep=1pt]
  \item \textbf{[Oral]} \texttt{NeurIPS\_2023\_0714} \textemdash\ QLoRA quantizes the static base weights to 4-bit NF4 while routing all adaptation through a full-precision LoRA adapter --- a textbook example of replacing one structural artifact (the frozen weight tensor) with a compact reparameterization while leaving inference call count untouched.
  \\ \textcolor{gray}{\footnotesize \textit{Lesson:} QLoRA: 4-bit NF4 base + full-precision LoRA adapter --- replace the static structural artifact (frozen weights) with a compact reparameterization, leave inference call count untouched.}
  \item \textbf{[Oral]} \texttt{ICML\_2023\_0538} \textemdash\ SparseGPT decomposes global pruning into per-layer Hessian-guided reconstruction problems, the canonical 'compress the static weight matrix in place' move that defines this cluster against trajectory-shortening tactics.
  \\ \textcolor{gray}{\footnotesize \textit{Lesson:} SparseGPT: decompose global pruning into per-layer Hessian-guided reconstruction --- compress the static weight matrix in place via local certified surgery.}
  \item \textbf{[Oral]} \texttt{ICML\_2024\_1032} \textemdash\ ExCP discovers that optimizer-momentum tensors are statistically reconstructable from weights and uses that entanglement to compress checkpoints --- exemplifying the cluster's move of finding structural redundancy in a frozen artifact and substituting a compact surrogate.
  \\ \textcolor{gray}{\footnotesize \textit{Lesson:} ExCP: discover optimizer-momentum tensors are statistically reconstructable from weights --- find structural redundancy in a frozen artifact and substitute a compact surrogate.}
  \item \textbf{[Oral]} \texttt{ICLR\_2025\_1510} \textemdash\ MoDeGPT aligns low-rank decomposition granularity with each module's functional sub-components, instantiating the cluster's stratum-specific reparameterization pattern at the architectural level.
  \\ \textcolor{gray}{\footnotesize \textit{Lesson:} Align low-rank decomposition granularity with each module's functional sub-components --- stratum-specific reparameterization at the architectural level.}
  \item \textbf{[Reject]} \texttt{ICLR\_2024\_0868} \textemdash\ KVTQ proposes ternary KV-cache quantization but lacks the binding-constraint argument that distinguishes Oral compression papers --- reviewers can't tell whether ternary specifically (vs. INT2/INT4) is the load-bearing choice.
  \\ \textcolor{gray}{\footnotesize \textit{Lesson:} Ternary KV-cache quantization without binding-constraint argument --- reviewers can't tell whether ternary specifically is the load-bearing choice.}
  \item \textbf{[Reject]} \texttt{NeurIPS\_2023\_0750} \textemdash\ Subspace Projection ports a theoretically clean global low-rank object onto transformer weights but does not ablate whether layer-coupled approximation error was actually the dominant failure mode of layer-local methods, leaving the contribution underspecified.
  \\ \textcolor{gray}{\footnotesize \textit{Lesson:} Subspace projection onto transformer weights without ablating that layer-coupled approximation error was the dominant failure leaves the contribution underspecified.}
  \item \textbf{[Reject]} \texttt{ICLR\_2025\_1374} \textemdash\ ClusComp proposes vector clustering as a drop-in for scalar 1-2 bit quantization, but the case for vector-level structure being the binding constraint at extreme bit-widths is asserted rather than isolated against scalar baselines with comparable headroom.
  \\ \textcolor{gray}{\footnotesize \textit{Lesson:} Vector clustering as drop-in for scalar 1-2 bit quantization needs to isolate that vector-level structure is the binding constraint at extreme bit-widths.}
\end{itemize}
\end{tcolorbox}

\begin{tcolorbox}[colback=bgskill!50, colframe=cgreen!70, title={\textbf{C18}: Iterative Process Compression via Step Distillation \textcolor{gray}{\small (O7/H2/R6, conf: \textbf{medium})}}, breakable, before skip=4pt, after skip=4pt]
\textbf{Tactical pattern.} The shared move is to take an iterative generative or optimization procedure (typically a diffusion model with N sampling steps, or score-based optimization of a parameter field) and train a compressed student that produces the procedure's terminal output in 1 or log(N) forward passes --- by either recursively halving the step count via self-distillation (each round's student becomes the next round's teacher), or by replacing per-step regression targets with distribution-level objectives (moment matching the student's output distribution to the teacher's, matching the score of a data-model mixture, or distribution matching stabilized with a GAN loss). The verbs are recurse, halve, replace-regression, match-distribution, score-distill; the objects are sampling trajectories, score functions, mixture distributions, and two-step composite maps. Critically, the student is supervised on what the unrolled procedure converges to, not on the procedure's internal computations.

\textbf{Differentiation.} \textit{Versus sibling cluster 17 (Compact Reparameterization), the redundancy axis is different in kind: 17 attacks parametric and structural redundancy inside a single forward pass --- FITS strips low-energy frequency components, QLoRA quantizes weights and reroutes adaptation through a low-rank adapter --- leaving each pass cheaper but the procedure itself unchanged. Cluster 18 attacks temporal redundancy in the unrolled execution of an iterative procedure, leaving each forward pass as wide as the teacher's but collapsing 1000 sampling steps into 1--4. This forces a different supervision problem: sibling 17 students can be matched output-for-output on a fixed input distribution, but cluster 18 students see inputs that depend on their own previous outputs, which is why per-sample regression breaks down and distribution-level objectives (moment matching in ICML\_2025\_1731, mixture score in ICML\_2025\_1793, GAN-stabilized DMD in ICML\_2023\_0471) become the load-bearing innovation rather than a stylistic choice.}

\textbf{When to pick.} Pick this tactic when the deployment cost driver is the sequential depth of an iterative procedure (diffusion sampling steps, score-based optimization iterations, MCMC chain length) rather than per-pass FLOPs or parameter count --- i.e., when "1000 steps to 1 step" is what unlocks deployment, not "make one step cheaper." A second signal: the procedure's outputs are defined by a converged distribution rather than a single deterministic target, so distribution-level matching is available as a substitute for per-step supervision.

\textbf{Tactical failure.} Reject papers in this cluster fail in two recognizable ways. (1) Repurposing the teacher-student machinery for a setting where the temporal-redundancy argument does not hold, so the contribution reduces to a generic distillation application without a step-compression payoff: NeurIPS\_2023\_0623 (EmbedDistill) realigns teacher-student embedding geometry for IR, but there is no iterative procedure to compress; ICLR\_2024\_0887 calibrates a binary network using a balanced full-precision teacher, again single-pass; ICLR\_2025\_1406 reuses teacher-guided distillation for TTS unlearning where the teacher emits speaker-randomized targets rather than compressed-step targets. (2) Proposing a structurally novel decomposition without proving the implied sub-targets are easier than the original problem: ICLR\_2023\_0314 (Progressive KD via boosting) trains each student on the previous ensemble's residual but never establishes that those residuals are more learnable by small students than the original task --- the central premise that motivates step-by-step compression goes unverified. ICLR\_2024\_0851 (HyperFields) replaces per-instance NeRF optimization with a meta-network distilled from individually optimized specialists, but treats fidelity preservation as an empirical claim rather than tightly isolating it; ICML\_2025\_1766 (OD$^3$) extends dataset-distillation step-collapse to detection but ends up engineering a placement-and-filter pipeline whose iterative loop is heuristic rather than a principled compression of an iterative procedure.

\textbf{Examples.}
\begin{itemize}[leftmargin=12pt,topsep=2pt,itemsep=1pt]
  \item \textbf{[Oral]} \texttt{ICML\_2025\_1731} \textemdash\ Canonical instance of the tactic's distribution-level pivot: replaces step-by-step bootstrapping (whose intermediate targets are themselves noisy student outputs) with a moment-matching objective on the output distribution, enabling single-stage step compression without an external teacher.
  \\ \textcolor{gray}{\footnotesize \textit{Lesson:} Replace step-by-step bootstrapping with a moment-matching objective on the output distribution --- distribution-level pivot enables single-stage step compression without an external teacher.}
  \item \textbf{[Oral]} \texttt{ICML\_2025\_1793} \textemdash\ Direct one-step training by matching the score of a data-model mixture distribution --- the purest form of bypassing the multi-step procedure entirely rather than distilling it, which is the cluster's aggressive endpoint.
  \\ \textcolor{gray}{\footnotesize \textit{Lesson:} Direct one-step training by matching the score of a data-model mixture distribution --- bypass the multi-step procedure entirely rather than distilling it.}
  \item \textbf{[Oral]} \texttt{ICML\_2023\_0471} \textemdash\ Refines the distribution-matching distillation objective by removing the regression loss and adding a GAN loss, illustrating that the cluster's progress comes from sharpening the step-collapse training objective itself rather than scaling teachers or data.
  \\ \textcolor{gray}{\footnotesize \textit{Lesson:} Refine distribution-matching distillation by removing the regression loss and adding a GAN loss --- progress comes from sharpening the step-collapse objective itself.}
  \item \textbf{[Oral]} \texttt{ICLR\_2025\_1544} \textemdash\ Provides the analytical justification specific to this tactic --- proves that progressive step-halving distillation induces an implicit curriculum giving provably better sample complexity than direct final-teacher distillation, explaining why recursive halving (not one-shot compression) is the dominant design.
  \\ \textcolor{gray}{\footnotesize \textit{Lesson:} Prove progressive step-halving distillation induces an implicit curriculum with provably better sample complexity --- analytic justification specific to the tactic.}
  \item \textbf{[Reject]} \texttt{ICLR\_2023\_0314} \textemdash\ Inverts the cluster's compression direction (decomposes one teacher into a boosted ensemble of small students fit on residuals) but never demonstrates that the residual sub-targets are easier than the original task --- the load-bearing assumption of step-by-step compression goes unverified.
  \\ \textcolor{gray}{\footnotesize \textit{Lesson:} Inverting the compression direction (one teacher $\to$ boosted ensemble) requires demonstrating the residual sub-targets are easier than the original task.}
  \item \textbf{[Reject]} \texttt{ICLR\_2024\_0851} \textemdash\ Applies multi-instance distillation to amortize per-instance NeRF optimization into a meta-network, but the claim that predicted parameters preserve individually-optimized fidelity is asserted via ablations rather than tightly isolated, exposing the cluster's recurring failure to prove the compressed map matches the unrolled one.
  \\ \textcolor{gray}{\footnotesize \textit{Lesson:} Multi-instance distillation amortizing per-instance NeRF needs tight isolation that predicted parameters preserve individually-optimized fidelity.}
  \item \textbf{[Reject]} \texttt{NeurIPS\_2023\_0623} \textemdash\ Uses the cluster's teacher-student vocabulary to realign embedding geometry for IR distillation, but with no iterative procedure to compress the work reads as generic output-level distillation dressed in the cluster's machinery.
  \\ \textcolor{gray}{\footnotesize \textit{Lesson:} Teacher-student vocabulary with no iterative procedure to compress reads as generic output-level distillation dressed in cluster machinery.}
\end{itemize}
\end{tcolorbox}

\subsection{Engineer Auxiliary Supervision Signal}
\textit{Parent methodology id: \texttt{auxiliary\_signal\_engineering}; 5 tactical clusters.}

\begin{tcolorbox}[colback=bgskill!50, colframe=cgreen!70, title={\textbf{C11}: Auxiliary Contrastive Signal for Representation Structure \textcolor{gray}{\small (O2/H3/R6, conf: \textbf{low})}}, breakable, before skip=4pt, after skip=4pt]
\textbf{Tactical pattern.} Add an InfoNCE-style pairwise pull-together / push-apart loss as an auxiliary head on top of an existing primary objective, where the central creative move is defining what counts as a positive vs. negative pair in a domain that lacks the canonical "same-class-vs-different-class" partition. Concrete instantiations across this cluster: pairs from continuous-target distance (ICLR\_2023\_0354, ICLR\_2024\_0901), teacher-output vs. student-output (ICML\_2025\_1702), preferred-vs-pruned branches in a search tree (NeurIPS\_2024\_1172), post-adaptation parameter states keyed by task identity (NeurIPS\_2024\_1246), model-mined globally-hardest negatives (ICLR\_2021\_0004), noised positive vs. noised negative anchor boxes (ICLR\_2023\_0205), or safe-vs-unsafe outcome trajectories (ICLR\_2023\_0330). The contrastive loss is bolted on, not a replacement for the primary task loss; it specifies a representation-space geometry (ordinality, isotropy, task-discriminability, hard-negative separation) that the end-to-end loss alone leaves unconstrained.

\textbf{Differentiation.} \textit{The defining move here is direct pairwise similarity supervision in representation space, with no intermediate discrete labels and no model-self-critique loop. Sibling cluster 21 (Pseudo-Label Guided Representation-Objective Alignment) routes through generated discrete cluster assignments and then aligns embeddings to those assignments --- cluster 11 skips that intermediate, supervising relations between pairs directly via continuous distance/ranking signals. Sibling cluster 32 (Self-Supervised Feedback Loop Internalization) has the model produce critiques or tokens that re-enter as supervision for its own outputs; cluster 11's pair structure comes from external givens (label distance, teacher logits, search-tree pruning order) rather than from the model judging itself. Sibling cluster 29 (Surrogate Supervision Signal Engineering) repurposes a different model or modality as a target signal; cluster 11 instead keeps the primary signal but adds a geometric constraint on how the encoder lays out points pairwise.}

\textbf{When to pick.} Pick this instantiation when the primary loss optimizes outputs correctly but leaves the encoder's intermediate geometry uncontrolled (representation collapse, anisotropy, ordinality mismatch, query-anchor confusion), AND you can name a domain-native pair-construction rule --- output distance, ranking inside a search/pruning step, teacher-vs-student, hard-negative mining, or paired safe/unsafe trajectories. If you cannot articulate the positive/negative rule in one sentence without hand-waving, prefer sibling 21 (go through pseudo-labels) or sibling 29 (repurpose a surrogate supervisor).

\textbf{Tactical failure.} Rejects in this cluster fail in two recurring ways. First, the positive/negative construction is heuristic in a regime where standard contrastive geometry no longer holds: ICLR\_2023\_0354 and ICLR\_2024\_0901 both define positivity via continuous-label proximity for regression, but neither convincingly defended their similarity threshold or showed it dominates simpler smoothness regularizers; ICLR\_2023\_0188 attempts intra-sequence contrast under causal masking, which restricts which hidden states can be legitimately compared and undermines the isotropy claim; ICLR\_2023\_0330 reuses a contrastive classifier as a calibrated risk estimator even though contrastive loss does not minimize calibration error. Second, papers add the auxiliary loss as a low-risk graft and fail to isolate its causal contribution: ICLR\_2024\_0824 jointly tunes alignment loss + schedule + pre-training reuse without ablations attributing gains specifically to the contrastive term, leaving reviewers unable to rule out simpler explanations. The shared signature is "we added a contrastive loss and numbers improved" without a principled defense of the pair definition or a clean ablation that the contrastive geometry --- not co-introduced regularization, augmentation, or schedule --- drove the result.

\textbf{Examples.}
\begin{itemize}[leftmargin=12pt,topsep=2pt,itemsep=1pt]
  \item \textbf{[Oral]} \texttt{ICML\_2025\_1702} \textemdash\ Defines the positive/negative pair as teacher-response vs. student-response within distillation, making the explicit pull-toward-teacher / push-from-self pair construction the core contribution rather than a side regularizer.
  \\ \textcolor{gray}{\footnotesize \textit{Lesson:} Define positive/negative as teacher-vs-student distillation responses --- explicit pull-toward-teacher / push-from-self pair construction becomes the contribution, not a side regularizer.}
  \item \textbf{[Oral]} \texttt{NeurIPS\_2024\_1246} \textemdash\ Lifts the pair definition out of data space entirely and into model-update space (post-adaptation parameters as the contrasted entities), exemplifying how the cluster's value comes from inventing the right pair semantics for the domain.
  \\ \textcolor{gray}{\footnotesize \textit{Lesson:} Lift pair definition out of data space into model-update space (post-adaptation parameters) --- invent the right pair semantics for the domain.}
  \item \textbf{[Reject]} \texttt{ICLR\_2024\_0901} \textemdash\ Uses mixup interpolation to manufacture soft positive pairs for regression but does not show the manufactured pair structure outperforms simple label-smoothness regularization, leaving the contrastive geometry's marginal value unproven.
  \\ \textcolor{gray}{\footnotesize \textit{Lesson:} Mixup-interpolated soft positive pairs for regression must outperform simple label-smoothness regularization to prove the contrastive geometry has marginal value.}
  \item \textbf{[Reject]} \texttt{ICLR\_2023\_0188} \textemdash\ Grafts a contrastive isotropy loss onto causal LMs without resolving how directional masking restricts admissible comparisons, so the claimed representation-space property is not cleanly enforced by the auxiliary objective.
  \\ \textcolor{gray}{\footnotesize \textit{Lesson:} Grafting contrastive isotropy onto causal LMs without resolving how directional masking restricts admissible comparisons leaves the representation property unenforced.}
  \item \textbf{[Reject]} \texttt{ICLR\_2023\_0330} \textemdash\ Repurposes a contrastive classifier as a calibrated risk estimator for safe RL even though contrastive loss does not target calibration error, exposing the failure mode of borrowing the loss form without checking its statistical assumptions.
  \\ \textcolor{gray}{\footnotesize \textit{Lesson:} Repurposing contrastive classifiers as calibrated risk estimators ignores that contrastive loss does not target calibration --- borrowing the loss form requires checking statistical assumptions.}
\end{itemize}
\end{tcolorbox}

\begin{tcolorbox}[colback=bgskill!50, colframe=cgreen!70, title={\textbf{C20}: Supervision Redesign via Training Signal Audit \textcolor{gray}{\small (O1/H2/R7, conf: \textbf{low})}}, breakable, before skip=4pt, after skip=4pt]
\textbf{Tactical pattern.} Treat the existing training signal as the *object of redesign* rather than augmenting it: empirically locate a structural mismatch between how labels/positives/negatives are assigned or consumed and the true generative process behind them, then surgically alter the assignment/consumption rule --- reweight (informational utility, boosting weights, negative-sample scores), reassign (shared/overlapping targets, structured negatives), geometrically decompose (project out confounder directions in representation space), or reformulate (continuous regression $\to$ discrete classification with softmax-pooled inference) --- while keeping the base architecture and loss family intact. The verb set is audit $\to$ diagnose $\to$ re-route signal, not "add an auxiliary loss" and not "build a new module." Concretely: papers replace uniform negative penalties with structured reweighting (ICLR\_2025\_1591), swap exclusive label binding for overlapping assignment (ICLR\_2024\_0793), train on the discrete labels that generated the continuous target then pool at inference (ICML\_2024\_1110), or use the data's own structural properties to directly supervise an attention mechanism that was previously trained only via downstream gradients (ICLR\_2021\_0033).

\textbf{Differentiation.} \textit{Sibling 11 ADDS a new contrastive loss term alongside the original objective; sibling 29 IMPORTS supervision from an external pre-trained source (GAN, simulator, foundation model) to substitute for missing labels; sibling 32 closes a feedback loop where the model critiques its own outputs. Cluster 20 differs from all three: it does not add a new loss term (vs 11), does not import external supervision into a previously-supervised channel (vs 29), and does not generate self-feedback (vs 32) --- it modifies the *routing* of signal that already exists in the pipeline. Q-Align (ICML\_2024\_1110) is the cleanest contrast vs sibling 11: where Supervised Contrastive Regression in cluster 11 keeps the regression target and adds a contrastive structuring term, Q-Align *changes the prediction target itself* from continuous to discrete and recovers the original target via inference-time pooling. Similarly, ICLR\_2021\_0033 (graph attention self-supervision) supervises the *existing* weighting mechanism with structural labels rather than adding a parallel contrastive head as cluster 11 would.}

\textbf{When to pick.} Pick this cluster when the gap you've identified lives in *signal routing* (which examples count as positive/negative, what target format aligns with the annotation process, which directions in the representation are confounded) rather than in *signal absence* --- i.e., when the supervision already nominally exists but is being mis-assigned, mis-aggregated, or mis-formatted. A strong signal: you can articulate the mismatch as a tractable assignment/reformulation rule that requires no new annotations and no new architecture branch.

\textbf{Tactical failure.} The recurring failure across rejects is "narrow audit, narrow lift": papers correctly diagnose a real assignment mismatch but the corrective rewiring delivers gains that read as incremental tuning rather than a foundational reframing, leaving reviewers unconvinced the redesign justifies the conceptual claim. ICLR\_2024\_0937 (RepoFusion) injects a structurally-available context into training but the gain over inference-time retrieval looks like routine engineering. ICLR\_2024\_0793 (CO-MOT) relaxes exclusive label binding only under a specific distributional regime, so the fix appears regime-bound rather than principled. ICLR\_2025\_1455 (Informative Data Selection) layers an information-theoretic reweighting on top of synthetic augmentation but the baseline already works and the bi-level objective adds machinery without commensurate accuracy lift. NeurIPS\_2023\_0632 (negative-sample scoring) and ICLR\_2025\_1591 (STARS structured negatives) both reweight a previously-uniform negative pool but rely on heuristic informativeness metrics whose generalization beyond the studied benchmark is not established. NeurIPS\_2023\_0762 (zero-shot robustification) decomposes representations geometrically but depends on hand-articulated confounder concepts whose coverage is unverifiable. The shared trap: the audit is plausible, the reroute is local, and the experiments cannot separate "the diagnosed mismatch is the real bottleneck" from "any reasonable regularizer on this dimension would have helped a similar amount."

\textbf{Examples.}
\begin{itemize}[leftmargin=12pt,topsep=2pt,itemsep=1pt]
  \item \textbf{[Oral]} \texttt{NeurIPS\_2025\_1925} \textemdash\ Recasts multimodal imbalance from a gradient/signal-rebalancing problem into a capability-gap problem and rewires training via boosting weights on the weaker classifier --- a pure redesign of how the existing supervision is *consumed across components* with no architectural change.
  \\ \textcolor{gray}{\footnotesize \textit{Lesson:} Recast multimodal imbalance from gradient/signal-rebalancing into a capability-gap problem --- rewire training via boosting weights on the weaker classifier, no architectural change.}
  \item \textbf{[Oral]} \texttt{ICML\_2024\_1110} \textemdash\ Q-Align is the canonical reformulation move: it traces continuous quality scores back to their discrete annotation source and replaces regression with discrete-label training plus softmax-pooled inference --- changing target format rather than adding any auxiliary signal.
  \\ \textcolor{gray}{\footnotesize \textit{Lesson:} Q-Align: trace continuous quality scores to their discrete annotation source and replace regression with discrete-label training plus softmax-pooled inference --- change target format, not signal.}
  \item \textbf{[Oral]} \texttt{ICLR\_2021\_0033} \textemdash\ Directly supervises the attention-weighting mechanism using the input graph's own structural properties as labels, exemplifying the cluster's audit-and-reroute pattern: the signal was already in the data, it just was not being used to train the right submodule.
  \\ \textcolor{gray}{\footnotesize \textit{Lesson:} Supervise the attention-weighting mechanism using the input graph's own structural properties as labels --- the signal was already in the data, just not training the right submodule.}
  \item \textbf{[Reject]} \texttt{ICLR\_2024\_0793} \textemdash\ CO-MOT diagnoses an exclusive-binding mismatch in MOT label assignment and relaxes it to overlapping targets, but the fix is bound to a specific distributional regime, illustrating the cluster's narrow-audit failure mode.
  \\ \textcolor{gray}{\footnotesize \textit{Lesson:} Diagnosing exclusive-binding mismatch in MOT and relaxing to overlapping targets is bound to a specific distributional regime --- the cluster's narrow-audit failure.}
  \item \textbf{[Reject]} \texttt{NeurIPS\_2023\_0632} \textemdash\ Argues negatives are informationally heterogeneous and scores them via a latent-variable metric --- a textbook signal-reweighting move that fails to convince reviewers because the informativeness metric is heuristic and benchmark-bound.
  \\ \textcolor{gray}{\footnotesize \textit{Lesson:} Latent-variable informativeness metric for negatives is heuristic and benchmark-bound --- signal-reweighting needs principled informativeness, not a learned scorer.}
  \item \textbf{[Reject]} \texttt{ICLR\_2025\_1591} \textemdash\ STARS reweights uniformly-treated negatives in skeleton self-supervision based on semantic proximity, but the redesign collides with other pipeline components whose objectives are geometrically incompatible --- the audit identified a real gap but the reroute introduced new ones.
  \\ \textcolor{gray}{\footnotesize \textit{Lesson:} Reweighting negatives by semantic proximity in skeleton SSL collides with other pipeline components --- the audit identified a real gap but the reroute introduced new ones.}
\end{itemize}
\end{tcolorbox}

\begin{tcolorbox}[colback=bgskill!50, colframe=cgreen!70, title={\textbf{C21}: Pseudo-Label Guided Representation-Objective Alignment \textcolor{gray}{\small (O6/H0/R7, conf: \textbf{medium})}}, breakable, before skip=4pt, after skip=4pt]
\textbf{Tactical pattern.} The shared tactical move is to run a clustering/structural-grouping procedure (hierarchical clustering, spectral pairwise-distance clustering, taxonomy lookup, or foundation-model semantic grouping) and harvest its discrete output as pseudo-labels that play a *dual* role: targets for predictive representation learning AND the alignment objective the embedding space must satisfy. Concretely, papers either (a) treat hierarchical/agglomerative cluster assignments at each level as multi-scale prediction targets (ICML\_2024\_1037), (b) co-optimize a clustering objective and a metric/geometric structure so that cluster assignments and embeddings converge jointly (ICML\_2024\_1061's structural-entropy-on-Lorentz, ICML\_2023\_0569's joint feature-graph diffusion), or (c) inject an external discrete signal (foundation-model categories in ICML\_2024\_1046, taxonomic ancestry in NeurIPS\_2024\_1335) as the clustering target. The defining verb is "use cluster assignments as supervision," and the defining object is "the same embedding the cluster assignments came from"---closing the representation$\leftrightarrow$objective loop with discrete labels rather than continuous similarity terms.

\textbf{Differentiation.} \textit{Versus sibling 11 (Auxiliary Contrastive Signal), which adds a *continuous* pairwise loss alongside an unchanged primary objective, this cluster commits to *discrete* pseudo-labels and folds clustering itself into the training objective so that the cluster identity IS the supervision target---not a regularizer on geometry. Versus sibling 29 (Surrogate Supervision Engineering), which repurposes a pretrained model's implicit constraints as a generic supervision oracle (e.g., 2D GANs supervising 3D reconstruction), this cluster's auxiliary signal is specifically a *partition* output that must be predicted, and the partition is the end-goal not just a teacher. Versus sibling 32 (Self-Supervised Feedback Loop), which uses the model's own critiques/loss-deltas as feedback, this cluster's signal is externally produced (clustering algorithm, taxonomy, foundation model) before being internalized. The marker is: if you remove the explicit "predict-which-cluster-this-belongs-to" head, the method collapses---whereas siblings would still train.}

\textbf{When to pick.} Pick this instantiation when the downstream task IS a partition (clustering, group discovery, mixture identification) and you have a *constructive* procedure that yields a candidate partition (hierarchical agglomeration, spectral pairwise distances, an external taxonomy, or a foundation model's category labels) that can be iterated against the representation. If the partition is only an evaluation artifact rather than something you need the embedding aligned to, prefer sibling 11's contrastive auxiliary instead.

\textbf{Tactical failure.} Rejects in this cluster fail when the pseudo-label loop is decorative rather than load-bearing: ICLR\_2023\_0199 stacks mutual-information maximization and pseudo-label regularization without showing that the clustering head materially shapes the MI geometry beyond what MI alone delivers, leaving the dual role unproven. ICLR\_2023\_0317 (ProsodyBERT) runs clustering offline once and freezes the labels, so coverage gaps and invalid feature values silently corrupt the supervision---the representation and clustering never co-evolve. NeurIPS\_2024\_1335 borrows the same hierarchical-pseudo-label idea that worked in oral ICML\_2024\_1037 but stops at "propagate contrastive signals through a taxonomy," without deriving theoretical guarantees or showing that ancestor-prototype consistency actually recovers traits unlearnable from flat labels. ICLR\_2025\_1394 imports Thomson-energy on cluster centroids as a uniformity regularizer but never closes the loop back to representation prediction; NeurIPS\_2024\_1322 and NeurIPS\_2023\_0649 substitute a distance metric or add Bernoulli randomness inside a clustering pipeline without engineering any new pseudo-label supervision at all---they sit in the cluster by topical proximity but miss the tactic's core dual-role move. ICLR\_2023\_0305 cleverly relocates clustering to gradient space but treats the recovered groups as a downstream input to robust optimization rather than as supervision that re-shapes the representation. The common thread: the discrete signal exists but does not actually drive a representation update that re-improves the partition.

\textbf{Examples.}
\begin{itemize}[leftmargin=12pt,topsep=2pt,itemsep=1pt]
  \item \textbf{[Oral]} \texttt{ICML\_2024\_1037} \textemdash\ Canonical instantiation: hierarchical clustering's multi-level cluster assignments are directly used as multi-scale prediction targets, making the partition output the supervision and the representation jointly trained to predict it at every granularity.
  \\ \textcolor{gray}{\footnotesize \textit{Lesson:} Hierarchical clustering's multi-level cluster assignments become multi-scale prediction targets --- partition output IS the supervision, jointly trained at every granularity.}
  \item \textbf{[Oral]} \texttt{ICML\_2024\_1046} \textemdash\ Externalizes the discrete pseudo-label source to a foundation model's semantic categories, exemplifying the cluster's variant where the structured supervision comes from outside but still aligns the embedding to the clustering objective.
  \\ \textcolor{gray}{\footnotesize \textit{Lesson:} Externalize the discrete pseudo-label source to a foundation model's semantic categories --- structured supervision can come from outside while still aligning embedding to clustering.}
  \item \textbf{[Oral]} \texttt{ICML\_2024\_1061} \textemdash\ Tightest representation-objective coupling: derives a structural-entropy clustering objective optimizable inside Lorentz geometry so that the cluster assignments and the hyperbolic embedding co-optimize via the same loss rather than alternating.
  \\ \textcolor{gray}{\footnotesize \textit{Lesson:} Structural-entropy clustering objective optimizable inside Lorentz geometry --- cluster assignments and hyperbolic embedding co-optimize via the same loss, no alternation.}
  \item \textbf{[Oral]} \texttt{ICML\_2023\_0482} \textemdash\ Shows the tactic on stochastic-process trajectories: pairwise-distance estimators feed spectral clustering whose assignments seed iterative refinement, breaking the chicken-and-egg of representation-vs-cluster identification.
  \\ \textcolor{gray}{\footnotesize \textit{Lesson:} Pairwise-distance estimators feed spectral clustering whose assignments seed iterative refinement --- break the chicken-and-egg of representation-vs-cluster identification.}
  \item \textbf{[Reject]} \texttt{ICLR\_2023\_0199} \textemdash\ Combines MI maximization with pseudo-label regularization but never establishes that the clustering head changes what MI would already produce---the dual role is asserted, not demonstrated.
  \\ \textcolor{gray}{\footnotesize \textit{Lesson:} MI maximization + pseudo-label regularization without establishing the clustering head changes what MI would already produce --- dual role asserted, not demonstrated.}
  \item \textbf{[Reject]} \texttt{ICLR\_2023\_0317} \textemdash\ Generates pseudo-labels via offline clustering once and freezes them, so coverage gaps in feature validity poison the supervision and the representation-cluster loop never closes during training.
  \\ \textcolor{gray}{\footnotesize \textit{Lesson:} Offline-clustering pseudo-labels frozen once mean coverage gaps in feature validity poison supervision --- the representation-cluster loop never closes during training.}
  \item \textbf{[Reject]} \texttt{NeurIPS\_2024\_1335} \textemdash\ Mirrors the hierarchical-pseudo-label tactic of ICML\_2024\_1037 with a taxonomy-aligned prototype tree but stops at contrastive consistency without theory or evidence that ancestor-prototype supervision recovers structure flat labels miss.
  \\ \textcolor{gray}{\footnotesize \textit{Lesson:} Hierarchical taxonomy-aligned prototype trees with contrastive consistency need theory or evidence that ancestor-prototype supervision recovers structure flat labels miss.}
\end{itemize}
\end{tcolorbox}

\begin{tcolorbox}[colback=bgskill!50, colframe=cgreen!70, title={\textbf{C29}: Surrogate Supervision Signal Engineering \textcolor{gray}{\small (O8/H0/R2, conf: \textbf{low})}}, breakable, before skip=4pt, after skip=4pt]
\textbf{Tactical pattern.} These papers MINE supervision from a pre-existing external artifact that implicitly encodes the target structure, then pipe that mined signal as a loss term or training target into the model that needs it. The artifacts are concrete and named: a pretrained 2D GAN's manifold (ICLR\_2021\_0017, ICLR\_2021\_0035), a differentiable renderer's gradient (ICLR\_2022\_0132), a simulator's environment state (NeurIPS\_2023\_0622), a rule-parser over text (NeurIPS\_2024\_1298), Bayes-aggregated posteriors over a known label graph (NeurIPS\_2024\_1166), a multi-view diffusion model's synthesized captures (NeurIPS\_2024\_1171), or <8\% partial correspondence anchors (ICLR\_2025\_1431). The verb is "extract/repurpose/anchor" rather than "construct" or "self-critique" --- the signal already exists somewhere structured, and the move is to wire it in as supervision.

\textbf{Differentiation.} \textit{The supervision signal here is sourced EXTERNALLY from an artifact that already encodes the target structure, whereas siblings construct signals from the model itself or from the existing label distribution. Versus Cluster 32 (Self-Supervised Feedback Loop Internalization), which closes the loop using the model's own outputs/losses/critiques (e.g., Toolformer's loss-reduction signal, Self-RAG's reflection tokens), this cluster's signal is foreign to the model --- a renderer's gradient or a GAN's manifold provides knowledge the model itself cannot produce. Versus Cluster 21 (Pseudo-Label Guided Representation-Objective Alignment), which generates discrete pseudo-labels from in-flight clustering of the same data, this cluster uses physical simulators, rule extractors, or probabilistic priors as the signal source. Versus Cluster 11 (Auxiliary Contrastive Signal), which adds a representation-space geometric term, this cluster injects structured content (positions, depths, label posteriors), not similarity geometry.}

\textbf{When to pick.} Pick this tactic when you can name a concrete external artifact --- a pretrained generator, differentiable physics module, simulator, rule parser, or sparse partial-label oracle --- that implicitly contains the answer your primary loss cannot reach. If the only signal source is the model's own outputs or in-batch clustering, pick a different sibling.

\textbf{Tactical failure.} With only 2 reject papers (ICLR\_2023\_0285, ICLR\_2023\_0328), this characterization is speculative. The shared failure pattern is that the "surrogate signal" is not clearly tied to a capability the primary loss demonstrably cannot supervise, so reviewers read it as plumbing rather than disambiguation. ICLR\_2023\_0285 replaces a derived motion signal with raw frames plus cross-sample consistency regularization --- the consistency term reads as a generic regularizer, and there is no external artifact whose structured knowledge is being mined; the move collapses into "remove a flawed signal," not "inject a sourced signal." ICLR\_2023\_0328 supervises NAS candidates via cross-layer distillation from a teacher, but the teacher is a standard distillation source, not an artifact whose latent structure is being newly exploited; reviewers saw two off-the-shelf techniques fused without arguing what unique capability the cross-layer correspondence uniquely unlocks.

\textbf{Examples.}
\begin{itemize}[leftmargin=12pt,topsep=2pt,itemsep=1pt]
  \item \textbf{[Oral]} \texttt{ICLR\_2021\_0017} \textemdash\ Canonical instance of mining supervision from a pretrained external artifact: a 2D GAN's manifold and discriminator are repurposed as a photorealistic consistency loss for 3D shape, depth, and lighting --- the 3D signal is extracted, not generated by the trained model itself.
  \\ \textcolor{gray}{\footnotesize \textit{Lesson:} A 2D GAN's manifold and discriminator repurposed as a photorealistic consistency loss for 3D shape --- mine supervision from a pretrained external artifact's structure.}
  \item \textbf{[Oral]} \texttt{NeurIPS\_2024\_1166} \textemdash\ Replaces a hard one-to-one label assignment with a Bayes-aggregated probabilistic mapping computed from a known label-set relationship --- the supervision signal is mined from a structured probabilistic prior rather than constructed by the model.
  \\ \textcolor{gray}{\footnotesize \textit{Lesson:} Replace one-to-one labels with Bayes-aggregated probabilistic mapping from a known label-set relationship --- supervision from a structured probabilistic prior.}
  \item \textbf{[Oral]} \texttt{ICLR\_2025\_1431} \textemdash\ Anchors a Schrödinger bridge with <8\% partial correspondence pairs --- supervision is sourced from a sparse external oracle and threaded through a distance-preserving constraint, rather than self-generated.
  \\ \textcolor{gray}{\footnotesize \textit{Lesson:} Anchor a Schrödinger bridge with <8\% partial correspondence pairs --- sparse external oracle threaded through a distance-preserving constraint.}
  \item \textbf{[Oral]} \texttt{ICLR\_2022\_0132} \textemdash\ Uses a differentiable renderer's gradient as the supervision source: the renderer's optimization signal is mined to define a rendering-orthogonality constraint, turning the nuisance factor's own model into the corrective signal.
  \\ \textcolor{gray}{\footnotesize \textit{Lesson:} Use the differentiable renderer's gradient as supervision --- the nuisance factor's own model becomes the corrective signal.}
  \item \textbf{[Reject]} \texttt{ICLR\_2023\_0285} \textemdash\ Discards the derived motion signal and substitutes a cross-sample consistency regularizer over raw frames --- but no external artifact's structured knowledge is being mined, so the move reads as removing supervision rather than injecting a sourced surrogate.
  \\ \textcolor{gray}{\footnotesize \textit{Lesson:} Cross-sample consistency over raw frames removes supervision rather than mining it --- no external artifact's structured knowledge is being injected.}
  \item \textbf{[Reject]} \texttt{ICLR\_2023\_0328} \textemdash\ Uses cross-layer teacher-student distillation as the auxiliary signal, but the teacher is a generic reference rather than an artifact whose latent structure encodes a specific capability the primary loss cannot reach, so reviewers read it as plumbing rather than disambiguation.
  \\ \textcolor{gray}{\footnotesize \textit{Lesson:} Cross-layer teacher-student distillation reads as plumbing --- the teacher must be an artifact whose latent structure encodes a specific capability the primary loss cannot reach.}
\end{itemize}
\end{tcolorbox}

\begin{tcolorbox}[colback=bgskill!50, colframe=cgreen!70, title={\textbf{C32}: Self-Supervised Feedback Loop Internalization \textcolor{gray}{\small (O6/H4/R2, conf: \textbf{low})}}, breakable, before skip=4pt, after skip=4pt]
\textbf{Tactical pattern.} The shared tactic is to take a component that normally lives OUTSIDE the model --- tool invocation, retrieval decision, evaluation oracle, intermediate supervision, or documentation --- and internalize it through one of two concrete moves: (a) reify the external decision as special tokens in the model's own output vocabulary so it can be predicted via the standard LM head (Self-RAG's reflection tokens, Toolformer's API-call tokens, ToolkenGPT's toolkens), or (b) substitute the model's own loss, output, or critique for the missing supervisor --- using perplexity drop as a filter for self-generated tool calls (Toolformer), final response quality as a distant signal for intermediate modules (Hexa), verbal-trace explanations as the debugging oracle (Self-Debug), or recursive self-critique as the reward (RCI). The unifying verb is "fold externally-supplied feedback into the model's native generation/scoring machinery"; the unifying object is a previously-external pipeline component (tool, critic, doc, retriever) that becomes a token, a filter on self-generated data, or a re-prompted role of the same model.

\textbf{Differentiation.} \textit{Unlike sibling 11, which adds an auxiliary CONTRASTIVE loss over representation pairs (different points in the same batch), cluster 32 generates its supervision from the model's own forward passes --- the auxiliary signal is a token, a perplexity number, or a self-written critique, never a pairwise distance. Unlike sibling 20, which audits and REWEIGHTS existing labeled signals to fix a mismatch, cluster 32 has no labels at all to audit; it manufactures signal where none existed. Unlike sibling 21, whose pseudo-labels come from clustering algorithms or hierarchical priors imposed on the data, cluster 32's pseudo-labels come from re-prompting the same model in a different role (generator vs. critic). Unlike sibling 29, which exploits an EXTERNAL pre-trained model's implicit knowledge (a 2D GAN supervising 3D, an auxiliary foundation model auditing another), cluster 32 closes the loop on the model itself --- producer and supervisor share weights. The signature test: if you delete every external module/teacher and the loss still works because the model is filtering or critiquing its own outputs, you are in cluster 32.}

\textbf{When to pick.} Pick this tactic when the capability gap is a discrete pipeline decision (when to call a tool, whether output is supported, which adapter to load) AND the base model is already strong enough that its own confidence/loss/critique is a usable signal --- i.e., you have no oracle but you do have a competent generator. Avoid it when the model lacks the underlying competence the self-loop assumes, since the feedback then reinforces its own errors.

\textbf{Tactical failure.} The two rejects in this cluster show the central risk: a self-generated signal that is asserted to be reliable but never adversarially stress-tested against degeneration or correlation collapse. ICLR\_2024\_0850 (Hexa) bootstraps intermediate retrieval/summarization modules from final-response quality, but the paper concedes that the loop can reinforce poor intermediate outputs --- reviewers want a demonstration that the bootstrapping fixed point is non-degenerate, not just empirically better than zero supervision. ICLR\_2025\_1471 (LLM Critics for code) substitutes a gold-patch-guided LLM critic for execution, but the failure mode is exactly the one reviewers anticipate: semantic-similarity critique can rank syntactically plausible but execution-wrong patches highly, and the paper's correlation with execution metrics is treated as the headline result rather than as the threshold that would justify replacing execution. Both rejects illustrate that "the model's own X is a good enough signal" is the load-bearing claim of the entire tactic, and skipping rigorous validation of that claim --- adversarial cases, fixed-point analysis, comparison against the oracle the signal is meant to replace --- is what kills these papers. With only 2 rejects this pattern is suggestive rather than conclusive (confidence=low).

\textbf{Examples.}
\begin{itemize}[leftmargin=12pt,topsep=2pt,itemsep=1pt]
  \item \textbf{[Oral]} \texttt{ICLR\_2024\_0945} \textemdash\ Cleanest instance of move (a): retrieve/critique decisions are reified as reflection tokens in the LM vocabulary, turning a multi-component RAG pipeline into a single end-to-end-trained generator that emits its own pipeline control signals.
  \\ \textcolor{gray}{\footnotesize \textit{Lesson:} Reify retrieve/critique decisions as reflection tokens in the LM vocabulary --- multi-component RAG becomes single end-to-end-trained generator emitting its own pipeline control signals.}
  \item \textbf{[Oral]} \texttt{NeurIPS\_2023\_0744} \textemdash\ Cleanest instance of move (b): the model's own perplexity-reduction is the filter that decides which self-generated API calls become training data, with no human or task-reward signal in the loop.
  \\ \textcolor{gray}{\footnotesize \textit{Lesson:} Model's own perplexity-reduction is the filter for self-generated API-call training data --- no human or task-reward signal in the loop.}
  \item \textbf{[Oral]} \texttt{NeurIPS\_2023\_0669} \textemdash\ Pure self-critique exemplar --- the same frozen model is recursively re-prompted as producer then critic, substituting its internal evaluation capacity for an external reward function or labeled correction set.
  \\ \textcolor{gray}{\footnotesize \textit{Lesson:} Recursively re-prompt the same frozen model as producer then critic --- internal evaluation capacity substitutes for an external reward function.}
  \item \textbf{[Oral]} \texttt{ICLR\_2025\_1435} \textemdash\ Pushes the loop one level outward: the model uses its own tool-use failures as a signal to rewrite the documentation it then re-reads, making the external scaffold itself a learnable artifact updated by self-feedback.
  \\ \textcolor{gray}{\footnotesize \textit{Lesson:} Use tool-use failures to rewrite the documentation the model re-reads --- make the external scaffold a learnable artifact updated by self-feedback.}
  \item \textbf{[Reject]} \texttt{ICLR\_2024\_0850} \textemdash\ Bootstrap loop where final-response quality scores self-generated intermediate retrieval/summarization candidates --- the load-bearing claim that the loop won't reinforce bad intermediates is asserted but not rigorously characterized, the canonical failure of this tactic.
  \\ \textcolor{gray}{\footnotesize \textit{Lesson:} Bootstrap loops where final-response quality scores intermediate retrieval/summarization candidates need rigorous characterization that the loop won't reinforce bad intermediates.}
  \item \textbf{[Reject]} \texttt{ICLR\_2025\_1471} \textemdash\ Replaces execution-based code evaluation with an LLM critic guided by gold patches --- exactly the 'self-/learned-signal substitutes for the oracle it claims to replace' move, and reviewers found the correlation evidence insufficient to justify the substitution.
  \\ \textcolor{gray}{\footnotesize \textit{Lesson:} Replacing execution-based code evaluation with an LLM critic guided by gold patches needs correlation evidence --- self-/learned-signal substituting for an oracle is exactly the failure shape.}
\end{itemize}
\end{tcolorbox}

\subsection{Align Heterogeneous Sources in Shared Space}
\textit{Parent methodology id: \texttt{shared\_latent\_alignment}; 1 tactical cluster.}

\begin{tcolorbox}[colback=bgskill!50, colframe=cgreen!70, title={\textbf{C24}: Unified Cross-Modal Latent Space Alignment \textcolor{gray}{\small (O45/H22/R52, conf: \textbf{high})}}, breakable, before skip=4pt, after skip=4pt]
\textbf{Tactical pattern.} Project heterogeneous input modalities (or task formats) into one shared latent space so a single model can operate uniformly across them, instantiated through three concrete moves: (a) discrete tokenization that converts each modality --- speech waveforms, video frames, time-series patches, image pixels, MIDI events, brain signals --- into a unified token vocabulary consumed by one autoregressive or masked transformer (VideoPoet, InfinityStar, BEATs, Mu$^2$SLAM, BEiT, Population Transformer, Video-LaVIT); (b) contrastive or masked-reconstruction objectives that pull modality-specific encoders toward shared anchors --- pretrained text embeddings, fixed vocabulary vectors, or each other (ProtST, TEST, CLIP decomposition, VTM, SimMTM, Cross-Modal Fine-Tuning); (c) lightweight projection layers (Q-Former, MosIT adapters, mapping nets) between frozen pretrained components that align feature distributions without retraining the heavy encoders (BLIP-2, NExT-GPT, GILL, 3D-LLM). The recurring verb is "tokenize-or-project all sources into one space, then run one operator over that space."

\textbf{Differentiation.} \textit{Singleton under parent methodology --- no sibling clusters exist.}

\textbf{When to pick.} Pick this instantiation when the heterogeneity is across input modalities or data formats (audio/video/text/time-series/3D/protein/EEG) and the goal is a single unified model that operates over all of them at inference --- not just transferring from one to another or aligning two domains of the same modality. The signal is "we have N qualitatively different input streams and want one backbone to consume any of them."

\textbf{Tactical failure.} The dominant reject failure is "unification-as-aesthetic": authors stack a frozen pretrained encoder, a contrastive/masked alignment loss, and a lightweight adapter on a new dataset, claim a unified cross-modal model, but cannot show the unification unlocks behavior the parts did not already deliver in isolation --- reviewers see additive component combination and ask which loss is doing the work (MULAN ICLR\_2025\_1515 explicitly diagnosed this; LaVie ICLR\_2024\_0878, DiffTell ICLR\_2025\_1399, Aya ICLR\_2025\_1357, face forgery joint embedding NeurIPS\_2023\_0688, MoE token merging ICLR\_2025\_1588 all triggered the same critique). A second mode is "shared-vocabulary-as-claim": papers tokenize two modalities into a joint vocabulary (unified music ICLR\_2025\_1638, ECG instruction tuning ICLR\_2025\_1414, ViT-for-MTSC ICLR\_2023\_0391, pinyin diffusion ICLR\_2024\_0783) without demonstrating that the joint vocabulary captures cross-modal dependencies the modality-specific models miss --- the unification is presentation, not mechanism. A third mode is "asymmetric cyclic alignment": back-translation or symbolic autoencoding across modalities with very different dimensionality/length (sign-language ICLR\_2024\_0986, symbolic AE ICLR\_2025\_1599, person re-ID ICLR\_2023\_0158) where the structural mismatch makes the cycle objective insufficient as the only alignment signal.

\textbf{Examples.}
\begin{itemize}[leftmargin=12pt,topsep=2pt,itemsep=1pt]
  \item \textbf{[Oral]} \texttt{ICML\_2024\_1143} \textemdash\ Built a unified discrete tokenizer over video, audio, and text and trained one LLM on the joint token stream --- the canonical instantiation of tokenize-everything-into-one-vocabulary.
  \\ \textcolor{gray}{\footnotesize \textit{Lesson:} Build a unified discrete tokenizer over video, audio, and text and train one LLM on the joint stream --- canonical tokenize-everything-into-one-vocabulary instantiation.}
  \item \textbf{[Oral]} \texttt{ICML\_2023\_0433} \textemdash\ Diagnosed that cross-modal transfer fails not from architecture but from distribution mismatch in the shared latent, then aligned distributions via a lightweight network --- exemplifies the projection-into-shared-space move.
  \\ \textcolor{gray}{\footnotesize \textit{Lesson:} Diagnose that cross-modal transfer fails from distribution mismatch in the shared latent (not architecture), then align distributions via lightweight network --- projection-into-shared-space.}
  \item \textbf{[Oral]} \texttt{ICLR\_2023\_0385} \textemdash\ Embedded heterogeneous task labels (depth maps, normals, segmentation) and inputs into one token space and replaced task-specific heads with non-parametric matching in that space.
  \\ \textcolor{gray}{\footnotesize \textit{Lesson:} Embed heterogeneous task labels and inputs into one token space and replace task-specific heads with non-parametric matching --- task-format unification.}
  \item \textbf{[Oral]} \texttt{ICML\_2023\_0511} \textemdash\ Aligned protein sequence encoders with biomedical text encoders via contrastive objectives over paired data, building a cross-type shared embedding that is the textbook contrastive instantiation.
  \\ \textcolor{gray}{\footnotesize \textit{Lesson:} Align protein sequence and biomedical text encoders via contrastive over paired data --- textbook contrastive instantiation across types.}
  \item \textbf{[Reject]} \texttt{ICLR\_2025\_1515} \textemdash\ Combined two protein encoders with a contrastive alignment loss and a lightweight adapter --- components individually validated, but the joint design exposed no new mechanism beyond what each part already provides.
  \\ \textcolor{gray}{\footnotesize \textit{Lesson:} Two protein encoders + contrastive alignment + lightweight adapter individually validated, but the joint design exposes no new mechanism beyond what each part provides.}
  \item \textbf{[Reject]} \texttt{ICLR\_2025\_1638} \textemdash\ Tokenized symbolic MIDI and acoustic waveform into a shared vocabulary for one autoregressive model, but did not show the joint vocabulary unlocks cross-modal capability beyond what the per-modality models already deliver --- surface unification, no mechanism.
  \\ \textcolor{gray}{\footnotesize \textit{Lesson:} Tokenizing MIDI and waveform into a shared vocabulary needs to unlock cross-modal capability beyond per-modality models --- surface unification, no mechanism.}
  \item \textbf{[Reject]} \texttt{ICLR\_2024\_0986} \textemdash\ Used cyclic back-translation to align sign-video and text without paired data, but the asymmetric structural gap (sequence length, dimensionality) made the cycle objective insufficient on its own --- exemplifies the asymmetric-cyclic failure mode.
  \\ \textcolor{gray}{\footnotesize \textit{Lesson:} Cyclic back-translation between sign-video and text fails when asymmetric structural gap (sequence length, dimensionality) makes the cycle objective insufficient.}
\end{itemize}
\end{tcolorbox}

\subsection{Aggregate Over Diverse Candidates}
\textit{Parent methodology id: \texttt{candidate\_aggregation\_over\_diversity}; 1 tactical cluster.}

\begin{tcolorbox}[colback=bgskill!50, colframe=cgreen!70, title={\textbf{C25}: Ensemble Aggregation over Diverse Candidates \textcolor{gray}{\small (O6/H3/R3, conf: \textbf{low})}}, breakable, before skip=4pt, after skip=4pt]
\textbf{Tactical pattern.} The shared tactical move is: pick a step in the pipeline that currently emits ONE output (a single decoded path, a single retrieved item, a single hyperparameter pick, a single agent's answer, a single draft head, a single trained subnetwork) and replace it with K independently generated candidates whose errors are deliberately decorrelated --- then collapse them via majority vote (ICLR\_2023\_0339), iterated peer-conditioning (ICLR\_2025\_1509), debate-and-revise rounds (ICML\_2024\_1047), weighted least-squares combination with provable target bounds (ICLR\_2023\_0165), Bayesian-network factor assembly (ICLR\_2025\_1360), pseudo-label averaging across foundation models (ICML\_2025\_1710), oracle-selected best-of-N from a population (NeurIPS\_2025\_1845), MoE-decoupled draft heads (ICML\_2025\_1734), or weight-averaging across pruned copies (ICML\_2023\_0475). The verbs are "sample/instantiate/parameterize K", then "vote/weight/synthesize/select"; the engineering effort goes into making the K candidates structurally diverse (different sampling temperatures, different experts, different sub-networks, different model families) so aggregation actually denoises rather than averaging within a single failure mode.

\textbf{Differentiation.} \textit{Singleton under parent methodology --- no sibling clusters exist.}

\textbf{When to pick.} Pick this tactic when (a) the task has a verifier or compressible answer space (discrete labels, executable code, a learnable scorer, a quality oracle) so a vote/selection rule is well-defined, and (b) you can produce candidates whose errors are *not* perfectly correlated --- via temperature sampling, heterogeneous models, decoupled heads, or differently-pruned sub-networks --- because aggregation only pays off when the diverse candidates fail in different places.

\textbf{Tactical failure.} Reject papers in this cluster fail by being unable to attribute the headline gain to the aggregation move itself. ICLR\_2023\_0401 stacks two retrievers plus a re-ranker plus a secondary training procedure, and reviewers cannot disentangle whether the hybrid candidate set or the extra training step is doing the work --- the "two complementary components" framing reads as additive composition of validated parts rather than a tested claim about diversity. ICML\_2025\_1734 (Jakiro) targets speculative-decoding draft diversity via MoE heads, but the contribution narrows to an engineering swap inside an existing pipeline (EAGLE), so reviewers question whether the diversity gain meaningfully exceeds what tuned single-head sampling already provides. NeurIPS\_2025\_1832 (AdaCoT) is the weakest fit: it reframes CoT as a triggering decision rather than aggregating multiple CoT samples, and reviewers see the "Pareto-optimal" framing as wrapper RL on top of a known reasoning enhancement rather than a new aggregation primitive. The shared failure is: aggregation looks promising on paper but the gain leaks into a tuned auxiliary stage (extra training, RL controller, MoE plumbing), leaving the diversity claim unfalsifiable.

\textbf{Examples.}
\begin{itemize}[leftmargin=12pt,topsep=2pt,itemsep=1pt]
  \item \textbf{[Oral]} \texttt{ICLR\_2025\_1509} \textemdash\ MoA is the canonical instance of this tactic: take K heterogeneous LLMs, feed each one all peers' raw outputs as context, layer the synthesis --- aggregation IS the architecture, not a post-hoc trick.
  \\ \textcolor{gray}{\footnotesize \textit{Lesson:} MoA: K heterogeneous LLMs each fed all peers' outputs as context, layered synthesis --- aggregation IS the architecture, not a post-hoc trick.}
  \item \textbf{[Oral]} \texttt{ICLR\_2023\_0165} \textemdash\ Provides the theoretical backbone of the cluster: instead of selecting a single best DA model, derive a weighted-least-squares aggregation whose weights are estimable from source data alone with provable target-error bounds --- turns aggregation into a principled selector replacement.
  \\ \textcolor{gray}{\footnotesize \textit{Lesson:} Weighted-least-squares aggregation with weights estimable from source data alone, provable target-error bounds --- turn aggregation into a principled selector replacement.}
  \item \textbf{[Oral]} \texttt{ICML\_2023\_0475} \textemdash\ Demonstrates the most aggressive form of the tactic: extract an ensemble of diverse pruned subnetworks from a SINGLE training pass via model-soup weight averaging, collapsing what was an O(k)-runs ensemble construction into O(1).
  \\ \textcolor{gray}{\footnotesize \textit{Lesson:} Extract an ensemble of diverse pruned subnetworks from a SINGLE training pass via model-soup averaging --- collapse O(k)-runs to O(1).}
  \item \textbf{[Oral]} \texttt{NeurIPS\_2025\_1845} \textemdash\ Inference-time variant of the tactic for RL: a converged policy's mode is suboptimal, so sample a diverse population of trajectories and let a quality oracle select the best --- aggregation breaks the performance ceiling without retraining.
  \\ \textcolor{gray}{\footnotesize \textit{Lesson:} A converged policy's mode is suboptimal; sample diverse trajectories and let a quality oracle select --- break the performance ceiling without retraining.}
  \item \textbf{[Reject]} \texttt{ICLR\_2023\_0401} \textemdash\ Combines two parallel retrievers and a joint re-ranker into a hybrid candidate set, but the aggregation gain is entangled with a layered training procedure, so the independent contribution of diversity-via-combination cannot be isolated.
  \\ \textcolor{gray}{\footnotesize \textit{Lesson:} Two retrievers + joint reranker entangles aggregation gain with a layered training procedure --- diversity-via-combination contribution can't be isolated.}
  \item \textbf{[Reject]} \texttt{ICML\_2025\_1734} \textemdash\ Replaces a single MLP draft head with MoE-decoupled heads to generate decorrelated speculative candidates --- exactly the diversity-injection move --- but the diversity gain is hard to separate from generic MoE capacity gains in the verification pipeline.
  \\ \textcolor{gray}{\footnotesize \textit{Lesson:} MoE-decoupled draft heads is the diversity-injection move, but diversity gain entangles with generic MoE capacity gains in the verification pipeline.}
  \item \textbf{[Reject]} \texttt{NeurIPS\_2025\_1832} \textemdash\ Loosely fits the parent: instead of aggregating diverse CoT samples, it learns a Pareto-optimal binary trigger for whether to invoke CoT at all, so the aggregation primitive collapses into a single gating decision and the cluster's diversity-suppression mechanism is absent.
  \\ \textcolor{gray}{\footnotesize \textit{Lesson:} Learning a Pareto-optimal binary trigger for whether to invoke CoT collapses aggregation into a single gating decision --- diversity-suppression is absent.}
\end{itemize}
\end{tcolorbox}

\subsection{Localize via Causal Probing}
\textit{Parent methodology id: \texttt{mechanistic\_probing\_with\_intervention}; 2 tactical clusters.}

\begin{tcolorbox}[colback=bgskill!50, colframe=cgreen!70, title={\textbf{C37}: Localized Component Attribution for Targeted Intervention \textcolor{gray}{\small (O20/H1/R16, conf: \textbf{high})}}, breakable, before skip=4pt, after skip=4pt]
\textbf{Tactical pattern.} The shared move is **locate-then-act**: papers identify a minimal, addressable subset of internal entities (specific neurons, attention heads, parameters, latent SAE features, or individual training examples) that disproportionately mediate a target behavior or output property, then use that attribution as a *handle* for downstream action --- editing weights/activations to fix or transfer behavior (Cones, Value Activation Editing, attention-head safety surgery, SAE entity-recognition transfer to refusal), filtering/scoring data (TRAK, Snapshot Influence, Data Shapley in One Training Run, trajectory-conditioned influence), or producing perceptually-meaningful explanations rooted in localized components (Listenable Maps, CLUE, Self-Interpretable TS counterfactuals). The typical pipeline is: (1) define an attribution score over components, (2) verify localization by ablation/transfer, (3) demonstrate that the same handle yields a *new operational capability* --- not just a saliency map. Verbs and objects are concrete: "score per neuron via gradient norm and zero-out the top-k", "project Jacobian to estimate per-example influence and rank training data", "patch SAE feature direction to control downstream refusal".

\textbf{Differentiation.} \textit{Versus sibling cluster 38 (Causal Probing for Latent Algorithm Discovery), the question being answered is fundamentally different. Cluster 38 starts with a hypothesized algorithm (e.g., gradient descent in-context, word2vec-style vector arithmetic, n-step lookahead) and uses interventions to *verify or falsify* that the model internally implements it --- the deliverable is a mechanistic claim. Cluster 37 makes no such algorithmic hypothesis: it treats the behavior as opaque and asks only "*which components are responsible?*", then immediately converts that attribution into an actionable lever (edit a value neuron to fix arithmetic, knock out 0.006\% of heads to break safety, prune training samples by influence score). Where 38 papers terminate at "yes, the model implements A," 37 papers terminate at "and now you can use this localization to edit/transfer/select." A 38 paper would never just rank training examples by influence; a 37 paper would never derive a constructive existence proof of an algorithm.}

\textbf{When to pick.} Pick this tactic when the practical goal is to *manipulate* a behavior (suppress hallucination, transplant a capability, remove poisoned data, edit a fact) rather than to characterize what the model is computing. The signal is: you can articulate a downstream action you'd take *if* you knew which components/data were responsible.

\textbf{Tactical failure.} Reject papers in this cluster cluster around three pitfalls. (1) **Attribution-without-payoff** --- proposing a new attribution score with elegant axiomatic properties but no demonstration that the localization enables a novel intervention: NeurIPS\_2023\_0692 (formal feature attribution metric), ICLR\_2023\_0187 (Fourier-consistent attribution), ICLR\_2023\_0363 (effective-coalition baselines for IG), and ICLR\_2023\_0159 (Sobol region attribution) all derive better scores but stop at fidelity benchmarks. (2) **Critique without construction** --- formally refuting an existing attribution method without offering a localization handle that unlocks new capability: NeurIPS\_2023\_0577 (Shapley refutation) is the prototypical case. (3) **Method-transfer without leverage** --- porting an existing attribution technique to a new architecture/domain without showing the localization buys anything specific to that domain: ICLR\_2023\_0230 (LRP to graph RL), ICLR\_2024\_0928 (NMF on penultimate activations), ICLR\_2023\_0161 (counterfactual attention for medical imaging), NeurIPS\_2023\_0644 (GOAt training-free GNN attribution), and ICML\_2025\_1725 (Shapley for amino-acid selection) all stay at descriptive saliency. The throughline: the cluster's Oral papers always close the loop with an *operational* downstream use (edit, transfer, filter, refuse); the rejects stop at "we computed a faithfulness-improving score."

\textbf{Examples.}
\begin{itemize}[leftmargin=12pt,topsep=2pt,itemsep=1pt]
  \item \textbf{[Oral]} \texttt{ICML\_2023\_0430} \textemdash\ Canonical locate-then-edit instance: gradient-norm scoring isolates per-concept neurons in a diffusion model, then those exact neurons are flipped to drive customized generation --- attribution becomes the editing handle.
  \\ \textcolor{gray}{\footnotesize \textit{Lesson:} Gradient-norm scoring isolates per-concept neurons in a diffusion model; flip those exact neurons to drive customized generation --- attribution becomes the editing handle.}
  \item \textbf{[Oral]} \texttt{ICLR\_2025\_1526} \textemdash\ Demonstrates the cluster's signature payoff --- discovering safety lives in <0.006\% of attention heads and using that localization to perform targeted ablation --- rather than merely claiming safety is decodable.
  \\ \textcolor{gray}{\footnotesize \textit{Lesson:} Safety lives in <0.006\% of attention heads; targeted ablation uses that localization to perform precise intervention --- locate-then-act with a measurable payoff.}
  \item \textbf{[Oral]} \texttt{ICML\_2024\_1051} \textemdash\ Identifies specific FFN value neurons that store arithmetic associations and directly edits their activations at inference to correct outputs, exemplifying the locate$\to$intervene loop without retraining.
  \\ \textcolor{gray}{\footnotesize \textit{Lesson:} Specific FFN value neurons store arithmetic associations; edit their activations at inference to correct outputs --- locate$\to$intervene loop without retraining.}
  \item \textbf{[Oral]} \texttt{ICML\_2023\_0560} \textemdash\ TRAK shows the data-side instantiation of the tactic: random-projected Jacobian gives per-example attribution scores at scale, which then drive concrete operations like behavior debugging and training-set surgery.
  \\ \textcolor{gray}{\footnotesize \textit{Lesson:} TRAK: random-projected Jacobian gives per-example attribution at scale --- drives debugging and training-set surgery, the data-side instantiation of the tactic.}
  \item \textbf{[Reject]} \texttt{NeurIPS\_2023\_0577} \textemdash\ Pure critique of Shapley-based explanation without offering a localization handle that enables any downstream intervention --- the cluster's pattern requires the attribution to *do* something operational.
  \\ \textcolor{gray}{\footnotesize \textit{Lesson:} Pure critique of Shapley-based explanation without offering a localization handle that enables downstream intervention --- attribution must DO something operational.}
  \item \textbf{[Reject]} \texttt{ICLR\_2023\_0187} \textemdash\ Proposes Fourier-consistent feature attribution as a more axiomatic score but never converts the more-faithful localization into editing, transfer, or selection --- it stops at the metric.
  \\ \textcolor{gray}{\footnotesize \textit{Lesson:} More-faithful Fourier-consistent attribution that never converts to editing, transfer, or selection stops at the metric.}
  \item \textbf{[Reject]} \texttt{ICLR\_2023\_0230} \textemdash\ Adapts LRP to GNN-actor RL agents and reports relevance maps, but the localization is never used to edit behavior, prune samples, or transfer capability --- method-transfer without leverage.
  \\ \textcolor{gray}{\footnotesize \textit{Lesson:} LRP-on-GNN-actor relevance maps without using localization to edit, prune, or transfer is method-transfer without leverage.}
\end{itemize}
\end{tcolorbox}

\begin{tcolorbox}[colback=bgskill!50, colframe=cgreen!70, title={\textbf{C38}: Causal Probing for Latent Algorithm Discovery \textcolor{gray}{\small (O9/H5/R4, conf: \textbf{medium})}}, breakable, before skip=4pt, after skip=4pt]
\textbf{Tactical pattern.} The shared move is hypothesis-driven algorithm identification: pick a specific *named* known procedure (gradient descent, ridge regression, beam-search look-ahead, Word2Vec-style vector arithmetic, board-state world model, Fourier multiplication) as a candidate explanation for the system's behavior, then triangulate three evidence types against it --- (a) constructive proof or symbolic derivation that the architecture *can* implement the candidate, (b) behavioral discrimination experiments designed so competing candidate algorithms produce *distinguishable* outputs (e.g., GD vs. ridge vs. least squares on under-determined cases), and (c) probe-then-intervene verification where a linear/nonlinear probe decodes the candidate's intermediate quantities (board state, future box position, ideology score, relation vector) and a causal edit on those representations flips the downstream behavior in the predicted direction. Verbs are: name a known algorithm $\to$ prove implementability $\to$ decode its intermediate state $\to$ causally intervene on that state $\to$ check the prediction.

\textbf{Differentiation.} \textit{Cluster 37 localizes components in order to *use* the localization downstream --- editing concept neurons, attributing training-data influence, distilling local explanations --- so the question being answered is "which internal piece is responsible for this behavior so I can act on it?" Cluster 38 inverts the goal: localization and intervention are not instrumental but evidentiary, deployed to settle "is this system silently executing known algorithm X?" The candidate algorithm is named *before* the probe is built (ridge regression in ICLR\_2023\_0397, Othello board state in ICLR\_2023\_0227, Fourier multiplication in ICML\_2025\_1703, beam-search look-ahead in ICML\_2025\_1819), and the intervention's job is to falsify that named hypothesis rather than to enable a downstream edit. Cluster 37 papers rarely commit to a specific symbolic algorithm; cluster 38 papers rise or fall on whether the named algorithm survives a discriminating test.}

\textbf{When to pick.} Pick this tactic when you can name a *specific* prior algorithm (or world-model schema) that the system might be implicitly running and you can construct test inputs on which that algorithm gives a distinguishably different output than its rivals. If you only know "some structure is in there" but cannot name the candidate procedure or its intermediate state in closed form, you should be in cluster 37, not here.

\textbf{Tactical failure.} Rejects in this cluster fail by skipping one of the three evidentiary legs and presenting the remaining ones as if they sufficed. NeurIPS\_2023\_0670 names the candidate algorithm (Word2Vec-style vector arithmetic) and shows correlational alignment, but the intervention does not discriminate it from generic linear-probe-friendly geometry, so reviewers read the "simple algorithm" claim as redescription rather than falsification. ICML\_2025\_1819 has the look-ahead candidate and quantitative ground truth from chess, but stays descriptive --- it characterizes representational depth without a causal account of *how* the look-ahead is computed, so the ground-truth precision masks the absence of a discriminating intervention. NeurIPS\_2024\_1149 reduces the test to a linear classifier on frozen activations and treats classifier accuracy as a proxy for "structured representation," omitting both a named algorithm and any causal edit, which is exactly the weak-probe-as-evidence shortcut reviewers reject. NeurIPS\_2024\_1150 goes furthest in the wrong direction: a purely behavioral self-identification test with no internal probing at all, confusing operationalization with mechanism. The recurring pattern is offering evidence that is *consistent with* the candidate algorithm rather than *discriminating against* its alternatives.

\textbf{Examples.}
\begin{itemize}[leftmargin=12pt,topsep=2pt,itemsep=1pt]
  \item \textbf{[Oral]} \texttt{ICLR\_2023\_0227} \textemdash\ Canonical instance of the tactic --- names the candidate latent structure (the Othello board state), trains a probe to decode it, then causally edits the activations and shows downstream moves shift in the predicted direction, completing the hypothesize-decode-intervene loop in a fully controlled setting.
  \\ \textcolor{gray}{\footnotesize \textit{Lesson:} Othello: name the candidate (board state), train a probe to decode it, causally edit activations and show downstream moves shift --- complete the hypothesize-decode-intervene loop.}
  \item \textbf{[Oral]} \texttt{ICLR\_2023\_0397} \textemdash\ Names ridge regression / gradient descent as the candidate algorithms, gives a constructive proof the transformer can implement them, and discriminates between candidates by matching ICL outputs to GD/ridge/LS predictions across depth and noise --- the textbook three-leg algorithm-identification triangulation.
  \\ \textcolor{gray}{\footnotesize \textit{Lesson:} Name ridge regression / GD as candidates, give a constructive proof the transformer can implement them, discriminate via depth and noise --- three-leg algorithm-identification triangulation.}
  \item \textbf{[Oral]} \texttt{ICLR\_2025\_1457} \textemdash\ Picks model-based planning as the named candidate inside a model-free agent, decodes future box/goal positions with linear probes, and causally intervenes on those probe directions to flip the agent's plan --- using intervention specifically to argue the planning algorithm is implicitly executed, not merely correlated.
  \\ \textcolor{gray}{\footnotesize \textit{Lesson:} Name model-based planning as the candidate, decode future positions with linear probes, intervene to flip the plan --- argue the planning algorithm is implicitly executed via causal evidence.}
  \item \textbf{[Oral]} \texttt{ICML\_2025\_1703} \textemdash\ Names the Fourier Multiplication Algorithm as the candidate mechanism for grokking and falsifies the architecture-specific story by reproducing the same phenomenon in a non-neural system that shares only the theorized algorithmic substrate.
  \\ \textcolor{gray}{\footnotesize \textit{Lesson:} Name the Fourier Multiplication Algorithm and falsify the architecture-specific story by reproducing in a non-neural system sharing only the algorithmic substrate.}
  \item \textbf{[Reject]} \texttt{NeurIPS\_2023\_0670} \textemdash\ Names Word2Vec-style vector arithmetic as the candidate algorithm but provides evidence that is consistent with the claim rather than discriminating it from alternative linear encodings --- the intervention does not falsify rival hypotheses about the same activations.
  \\ \textcolor{gray}{\footnotesize \textit{Lesson:} Evidence consistent with Word2Vec-style arithmetic but not discriminating it from alternative linear encodings --- intervention does not falsify rival hypotheses.}
  \item \textbf{[Reject]} \texttt{ICML\_2025\_1819} \textemdash\ Has the named candidate (multi-step look-ahead) and precise chess ground truth but stays at the characterization stage, omitting the causal-intervention leg that distinguishes mechanism from descriptive correlation.
  \\ \textcolor{gray}{\footnotesize \textit{Lesson:} Named candidate (multi-step look-ahead) + precise ground truth but stays at characterization --- omits the causal-intervention leg distinguishing mechanism from correlation.}
  \item \textbf{[Reject]} \texttt{NeurIPS\_2024\_1149} \textemdash\ Substitutes a linear classifier's accuracy on frozen activations for both a named algorithm and a causal intervention, leaving the generalization-vs-memorization question undecidable by the evidence offered.
  \\ \textcolor{gray}{\footnotesize \textit{Lesson:} Linear classifier accuracy on frozen activations substitutes for both a named algorithm and a causal intervention --- leaves generalization-vs-memorization undecidable.}
\end{itemize}
\end{tcolorbox}

\subsection{Audit Evaluation via Controlled Perturbation}
\textit{Parent methodology id: \texttt{evaluation\_validity\_via\_controlled\_perturbation}; 1 tactical cluster.}

\begin{tcolorbox}[colback=bgskill!50, colframe=cgreen!70, title={\textbf{C35}: Benchmark Validity Auditing via Controlled Perturbation \textcolor{gray}{\small (O46/H29/R22, conf: \textbf{high})}}, breakable, before skip=4pt, after skip=4pt]
\textbf{Tactical pattern.} The tactic is item-level (not framework-level) controlled variation: take an existing benchmark B that the field already trusts, name a specific confound C (training-set similarity, modality shortcut, format/style preference, capability gap, surface memorization, frame-of-reference ambiguity, single-frame solvability, temporal contamination), then construct minimally-different variants of B's individual items where C is held fixed while the claimed target capability T is varied, or vice versa. Concrete verb-object instantiations recur: parameterize templates and inject semantically-irrelevant elements (GSM-Symbolic), construct matched pairs that equalize capability and vary only the subtle dimension (RM-Bench), filter or down-weight items solvable without one input modality (NeurIPS\_2024\_1161, EMMA), freeze one component and sweep the other (Cambrian-1), inject deliberately trivial baselines to measure how much "claimed-hard" signal they already capture (NeurIPS\_2024\_1334, ICLR\_2025\_1653), apply algebraic/temporal transformations that preserve task structure but defeat memorization (LLM-SRBench, LiveCodeBench), or stratify the test set along a previously-implicit axis (ICLR\_2025\_1561 similarity-aware splits). The shared move is reattribution: the headline score on B is decomposed into (capability-driven) and (C-driven) components, and the field's prior conclusion is restated only on the residual.

\textbf{Differentiation.} \textit{Singleton under parent methodology --- no sibling clusters exist.}

\textbf{When to pick.} Pick this instantiation when you can name one specific confound and construct minimally-different variants of the existing benchmark's items --- if you cannot point to such a confound or your "variants" require an entirely new task design, the move drifts toward generic benchmark-building rather than auditing. The strongest signal is that you can predict which prior conclusions will flip under the controlled variants before running them.

\textbf{Tactical failure.} The dominant failure mode in this cluster's rejects is that the "controlled variant" is itself confounded, so the audit cannot logically isolate what it claims to isolate. ICLR\_2024\_0958 (SynBench) builds a synthetic proxy whose distributional assumptions are never ablated, so any correlation with downstream quality could come from the proxy or from the confound. ICML\_2025\_1783 (quality-bias mitigation) and NeurIPS\_2024\_1162 (synthetic-artifact injection for IE) similarly inject controlled stimuli without demonstrating the injection itself is uncorrelated with the property being measured. A second pattern is "audit without falsification": ICLR\_2023\_0387 (conceptual consistency) and ICLR\_2024\_0956 (style over substance) propose decomposed metrics but stop at descriptive measurement, never showing that prior published conclusions actually reverse under the new evaluation --- reviewers read this as a methodology proposal rather than an audit. A third pattern is bundling the audit with a fix (ICLR\_2025\_1402 direct-judgement DPO, ICLR\_2025\_1474 LBG blocking, NeurIPS\_2024\_1291 RLHF prompt-template debias, ICLR\_2024\_0869 compositional decision-making), which dilutes both contributions; the successful Oral papers in this cluster (RM-Bench, GSM-Symbolic, ICLR\_2024\_0985 unprocessing) almost always publish the audit alone and let the rank-reversal carry the paper. A fourth pattern is over-specific benchmarks (ICLR\_2025\_1422 ERiC-UP\textasciicircum{}3, ICLR\_2023\_0295 dynamic vision, ICLR\_2023\_0393 visual transformation telling, NeurIPS\_2025\_1841 motion-centric eval) that audit a niche existing eval but produce no transferable methodological lesson --- reviewers want the audit to expose a confound that generalizes beyond the specific benchmark touched.

\textbf{Examples.}
\begin{itemize}[leftmargin=12pt,topsep=2pt,itemsep=1pt]
  \item \textbf{[Oral]} \texttt{ICLR\_2025\_1443} \textemdash\ Quintessential parameterized-template tactic: re-instantiates GSM8K items under value substitution, complexity scaling, and semantically-valid-but-logically-irrelevant element injection to isolate surface-pattern reliance from genuine generalization.
  \\ \textcolor{gray}{\footnotesize \textit{Lesson:} Re-instantiate GSM8K under value substitution + complexity scaling + irrelevant-element injection --- isolate surface-pattern reliance from genuine generalization via parameterized templates.}
  \item \textbf{[Oral]} \texttt{ICLR\_2025\_1565} \textemdash\ Canonical matched-pair construction: pairs are equalized on capability (both from strong models) and varied only on the subtle quality/style dimension, directly controlling out the capability-gap confound that made prior reward benchmarks too easy.
  \\ \textcolor{gray}{\footnotesize \textit{Lesson:} Pairs equalized on capability (both from strong models) and varied only on subtle quality/style --- directly control out the capability-gap confound.}
  \item \textbf{[Oral]} \texttt{NeurIPS\_2024\_1161} \textemdash\ Modality-shortcut filter applied to the existing eval: runs LLM-only inference to identify items answerable from text alone, then filters or down-weights them so that residual scores reflect genuine cross-modal capability.
  \\ \textcolor{gray}{\footnotesize \textit{Lesson:} LLM-only inference identifies items answerable from text alone; filter or down-weight them --- modality-shortcut filter applied to the existing eval.}
  \item \textbf{[Oral]} \texttt{ICLR\_2025\_1630} \textemdash\ Direct ablation of the named confound: fine-tunes older base models on task-specific format data, shows headline performance gaps collapse, and reattributes 'emergence' to format-conditioning rather than capability scaling.
  \\ \textcolor{gray}{\footnotesize \textit{Lesson:} Fine-tune older base models on task-specific format data; show headline gaps collapse --- direct ablation of the named confound (test-task fine-tuning).}
  \item \textbf{[Reject]} \texttt{ICLR\_2024\_0958} \textemdash\ Builds a synthetic-data proxy without ablating whether the proxy's distributional assumptions actually capture the binding constraint --- the controlled variant is itself confounded, so any correlation with downstream quality cannot be cleanly attributed.
  \\ \textcolor{gray}{\footnotesize \textit{Lesson:} Synthetic-data proxy without ablating that distributional assumptions capture the binding constraint --- the controlled variant is itself confounded.}
  \item \textbf{[Reject]} \texttt{ICLR\_2024\_0956} \textemdash\ Proposes the right decomposition (style vs. substance) but stops at descriptive measurement instead of demonstrating that prior LLM-as-judge rankings actually flip under the decomposed evaluation, reading as a methodology proposal rather than a binding audit.
  \\ \textcolor{gray}{\footnotesize \textit{Lesson:} Right decomposition (style vs substance) but stops at descriptive measurement --- needs to demonstrate prior LLM-as-judge rankings flip under the decomposed evaluation.}
  \item \textbf{[Reject]} \texttt{ICLR\_2025\_1402} \textemdash\ Bundles the judge-bias audit with a DPO-based remediation, diluting both contributions; the cluster's successful papers almost always publish the audit alone and let the rank-reversal carry the paper.
  \\ \textcolor{gray}{\footnotesize \textit{Lesson:} Bundling the audit with a DPO-based remediation dilutes both contributions --- successful papers publish the audit alone and let the rank-reversal carry the paper.}
\end{itemize}
\end{tcolorbox}

\subsection{Sharpen a Single Analytic Step}
\textit{Parent methodology id: \texttt{analytic\_refinement\_of\_loose\_bounds}; 1 tactical cluster.}

\begin{tcolorbox}[colback=bgskill!50, colframe=cgreen!70, title={\textbf{C15}: Tight Bound Recovery via Targeted Proof Refinement \textcolor{gray}{\small (O60/H1/R33, conf: \textbf{high})}}, breakable, before skip=4pt, after skip=4pt]
\textbf{Tactical pattern.} Locate a single load-bearing inequality, complexity parameter, or concentration tool inside a mature proof --- one that downstream rates pass through unchanged --- and replace it with a strictly sharper drop-in substitute, leaving the algorithm and the rest of the derivation intact. Concrete moves observed: swap a worst-case maximum information gain bound for an algorithm-trajectory-specific MIG (ICML\_2025\_1728, NeurIPS\_2025\_1890); replace the structural diameter with a tighter functional parameter (span) as the complexity index of an MDP (NeurIPS\_2024\_1306); replace L$\infty$ score accuracy with L2 in diffusion convergence proofs (ICLR\_2023\_0334); upgrade Freedman's scalar martingale concentration to its Hilbert-space analogue to match the algorithm's algebraic ambient space (NeurIPS\_2023\_0731 [sic: NeurIPS\_2024\_1307]); construct a new lower-bound hard instance tailored to weighted-average iterates rather than worst-case ones (ICML\_2023\_0555). The verbs are "replace," "substitute," "propagate" --- not "redesign." The papers explicitly claim the algorithm is unchanged and the cascade through existing proof scaffolding is the contribution.

\textbf{Differentiation.} \textit{Singleton under parent methodology --- no sibling clusters exist.}

\textbf{When to pick.} Pick this tactic when you can name a specific lemma, inequality, or complexity parameter inside a published analysis whose looseness you can quantify (e.g., worst-case-vs-trajectory, scalar-vs-functional-space, diameter-vs-span) AND when downstream theorems consume that object as a black box --- so a sharper replacement cascades automatically without rewriting the algorithm or its proof structure.

\textbf{Tactical failure.} Rejects in this cluster fail in three recurring ways. First, the "sharpening" is perceived as a routine port of known machinery rather than a genuinely tighter object: NeurIPS\_2024\_1182 (Adafactor) and ICLR\_2025\_1625 (parameter-free Adam/AdaGrad) adapt existing Adam/DoG proof templates with bookkeeping for low-rank or preconditioned terms, but reviewers read this as mechanical extension because the substitute object is not strictly sharper than what it replaces --- it merely covers a new optimizer. Second, the targeted "loose step" turns out not to be load-bearing: ICLR\_2023\_0245 derives gradient-like properties for hard thresholding, but reviewers note these properties were already established in the iterative-hard-thresholding literature, so the cascade improves nothing downstream. Third, the substitution introduces auxiliary assumptions that nullify the perceived tightness: ICML\_2025\_1674 (idealized Polyak SPS*) replaces global Lipschitz with a local expected-gradient bound that critics view as covertly assuming what it claims to prove, and NeurIPS\_2023\_0586 (adaptive quasi-Newton with cubic regularization) introduces cubic regularization to enable the rate, but reviewers see it as changing the algorithm rather than sharpening the analysis. Across these, the diagnostic is: was the swap audibly *tighter on the same object*, or was it a re-derivation under different assumptions / for a different algorithm?

\textbf{Examples.}
\begin{itemize}[leftmargin=12pt,topsep=2pt,itemsep=1pt]
  \item \textbf{[Oral]} \texttt{ICML\_2025\_1728} \textemdash\ Prototype of the tactic: prove a tighter maximum-posterior-variance bound, then propagate that single sharpened lemma through existing GP-bandit regret analyses unchanged, simultaneously improving multiple downstream rates.
  \\ \textcolor{gray}{\footnotesize \textit{Lesson:} Prove a tighter maximum-posterior-variance bound, then propagate that single sharpened lemma through existing GP-bandit regret unchanged --- multiple downstream rates improve simultaneously.}
  \item \textbf{[Oral]} \texttt{NeurIPS\_2025\_1890} \textemdash\ Replaces the worst-case MIG appearing in GP-UCB regret with an algorithm-trajectory-specific MIG bound and substitutes it into the existing analysis --- a textbook instance of swapping one bottleneck object for a sharper one.
  \\ \textcolor{gray}{\footnotesize \textit{Lesson:} Replace worst-case MIG in GP-UCB with algorithm-trajectory-specific MIG --- substitute one bottleneck object for a sharper one without touching the algorithm.}
  \item \textbf{[Oral]} \texttt{NeurIPS\_2024\_1306} \textemdash\ Identifies that diameter is a loose proxy and span is the structurally correct functional complexity parameter for average-reward MDPs, then re-derives matching upper and lower bounds in span --- a single-parameter substitution that cascades to minimax optimality.
  \\ \textcolor{gray}{\footnotesize \textit{Lesson:} Diameter is a loose proxy; span is the structurally correct functional complexity parameter for average-reward MDPs --- single-parameter substitution cascades to minimax optimality.}
  \item \textbf{[Oral]} \texttt{ICML\_2023\_0555} \textemdash\ Constructs a refined lower-bound hard instance tailored to weighted-average iterates (rather than the worst-case variance assumed in prior lower bounds), tightening the lower-bound side of an existing analysis without touching the algorithm.
  \\ \textcolor{gray}{\footnotesize \textit{Lesson:} Construct a refined lower-bound hard instance tailored to weighted-average iterates --- tighten the lower-bound side without touching the algorithm.}
  \item \textbf{[Reject]} \texttt{ICLR\_2023\_0245} \textemdash\ Targets gradient-like properties of the hard thresholding operator as the loose step, but reviewers identify the properties as already established in the literature --- the substitution exists but does not actually sharpen anything downstream.
  \\ \textcolor{gray}{\footnotesize \textit{Lesson:} Targeting properties already established in the literature is substitution without sharpening --- the targeted property must be genuinely loose downstream.}
  \item \textbf{[Reject]} \texttt{NeurIPS\_2024\_1182} \textemdash\ Adapts Adam-style convergence proofs to Adafactor by bounding the rank-1 approximation error, but reads as mechanical proof-template porting rather than identifying and replacing a genuinely loose primitive.
  \\ \textcolor{gray}{\footnotesize \textit{Lesson:} Adapting Adam-style proofs to Adafactor by bounding rank-1 approximation error reads as mechanical proof-template porting, not loose-primitive replacement.}
  \item \textbf{[Reject]} \texttt{ICML\_2025\_1674} \textemdash\ Relaxes global Lipschitz/smoothness to a local expected-gradient bound to unify smooth and Lipschitz Polyak analyses, but reviewers read the unifying assumption as assuming what is to be proved rather than as a strictly sharper substitute.
  \\ \textcolor{gray}{\footnotesize \textit{Lesson:} Relaxing global Lipschitz/smoothness to a local expected-gradient bound that 'unifies' analyses can be reading as assuming what is to be proved.}
\end{itemize}
\end{tcolorbox}
